% concatenated from 25_03-general-clique-factors-arxiv.tex
%• 05D10: "Ramsey theory" (MSC2020), • 05C35: "Extremal problems in graph theory" (MSC2020)
\documentclass[reqno, 12pt]{amsart}
%\usepackage[utf8]{inputenc}

\title{The complete picture for clique factors in randomly perturbed graphs}
\author{ Sylwia Antoniuk} 
\address{Department of Discrete Mathematics, Faculty of Mathematics and CS, Adam Mickiewicz University,
	Poznań, Poland}
\email{sylwia.antoniuk@amu.edu.pl}
\author {Nina Kam\v{c}ev}
\address{Department of Mathematics, Faculty of Science, University of Zagreb, Zagreb, Croatia}
\email{nina.kamcev@math.hr}
%
\author{
	Christian Reiher}
\address{Fachbereich Mathematik, Universität Hamburg, Hamburg, Germany}
\email{christian.reiher@uni-hamburg.de}
\author{Tadej Petar Tukara
}\address{Department of Mathematics, Faculty of Science, University of Zagreb, Zagreb, Croatia}
\email{tadej.tukara@math.hr}

\thanks{S.~Antoniuk was supported by Narodowe Centrum Nauki, grant 2024/53/B/ST1/00164.}
\thanks{N.~Kam\v cev and T.P.~Tukara are supported by the Croatian Science Foundation, project number HRZZ-IP-2022-10-5116 (FANAP)}

\makeatletter
\let\origsection=\section \def\section{\@ifstar{\origsection*}{\mysection}} 
\def\mysection{\@startsection{section}{1}\z@{.7\linespacing\@plus\linespacing}{.5\linespacing}{\normalfont\scshape\centering\S\hspace{1pt}}}
%\mhttps://www.overleaf.com/project/67c60e817c6bc036ad291da2#akeatother        

\usepackage{amsmath,amssymb,amsthm}
\usepackage{mathrsfs}
\usepackage{mathabx}\changenotsign
\usepackage{dsfont}
\usepackage{url}
\usepackage{colortbl}

\usepackage{xcolor}
\usepackage[backref]{hyperref}
\hypersetup{
	colorlinks,
	linkcolor={red!60!black},
	citecolor={green!60!black},
	urlcolor={blue!60!black}
}

%\usepackage[inline]{showlabels}
\usepackage[open,openlevel=2,atend]{bookmark}

%\usepackage[abbrev,msc-links,backrefs]{amsrefs}

\usepackage{doi}
\renewcommand{\doitext}{DOI\,}

%\renewcommand{\PrintDOI}[1]{\doi{#1}}

%\renewcommand{\eprint}[1]{\href{http://arxiv.org/abs/#1}{arXiv:#1}}

%%%%%****************bibliography
\usepackage{natbib}
%\NatBibNumeric
\def\bibfont{\small}%
\def\bibsep{\smallskipamount}%
\def\bibhang{24pt}%
\def\BIBand{and}%
\def\newblock{\ }%
\bibpunct[, ]{[}{]}{,}{n}{}{,}%

\usepackage{hyperref}
\usepackage[T1]{fontenc}
\usepackage{lmodern}
\usepackage[babel]{microtype}
\usepackage[english]{babel}

\usepackage{tikz}
\usetikzlibrary{calc,decorations.pathmorphing,decorations.pathreplacing}
\pgfdeclarelayer{background}
\pgfdeclarelayer{foreground}
\pgfdeclarelayer{front}
\pgfsetlayers{background,main,foreground,front}

\linespread{1.2}
\usepackage{geometry}
\geometry{left=27.5mm,right=27.5mm, top=25mm, bottom=25mm}

\numberwithin{equation}{section}
\numberwithin{figure}{section}
\renewcommand{\thefootnote}{\fnsymbol{footnote}} 

\usepackage{enumitem}
\def\rmlabel{\upshape({\itshape \roman*\,})}
\def\RMlabel{\upshape(\Roman*)}
\def\alabel{\upshape({\itshape \alph*\,})}
\def\Alabel{\upshape({\itshape \Alph*\,})}
\def\nlabel{\upshape({\itshape \arabic*\,})}

\def\Hlabel{\upshape({H\arabic*})}



\theoremstyle{plain}
\newtheorem{thm}{Theorem}[section]
\newtheorem{fact}[thm]{Fact}
\newtheorem{prop}[thm]{Proposition}
\newtheorem{clm}[thm]{Claim}
\newtheorem{claim}[thm]{Claim}
\newtheorem{prob}[thm]{Problem}
\newtheorem{cor}[thm]{Corollary}
\newtheorem{lemma}[thm]{Lemma}
\newtheorem{conclusion}[thm]{Conclusion}
\newtheorem{assertion}[thm]{Assertion}
\newtheorem{remark}[thm]{Remark}

\theoremstyle{definition}
\newtheorem{rem}[thm]{Remark}
\newtheorem{dfn}[thm]{Definition}
\newtheorem{definition}[thm]{Definition}
\newtheorem{exmp}[thm]{Example}
\newtheorem{quest}{Question}[section]
\newtheorem{question}{Question}[section]
\newtheorem{conj}[thm]{Conjecture}




\let\eps=\varepsilon
\let\theta=\vartheta
\let\rho=\varrho
\let\phi=\varphi
\let\vsigma=\varsigma
\let\Ups=\varUpsilon


\def\NN{\mathds N}
\def\ZZ{\mathds Z} 
\def\PP{\mathds P}
\def\QQ{\mathds Q}
\def\RR{\mathds R}

\def\ind{\mathds 1}

\def\cA{{\mathcal A}}
\def\cB{{\mathcal B}}
\def\cC{{\mathcal C}}
\def\cD{{\mathcal D}}
\def\cE{{\mathcal E}}
\def\cF{{\mathcal F}}
\def\cH{{\mathcal H}}
\def\cK{{\mathcal K}}
\def\cL{{\mathcal L}}
\def\cM{{\mathcal M}}
\def\cP{{\mathcal P}}
\def\cQ{{\mathcal Q}}
\def\cR{{\mathcal R}}
\def\cS{{\mathcal S}}
\def\cT{{\mathcal T}}
\def\cU{{\mathcal U}}
\def\cV{{\mathcal V}}
\def\cX{{\mathcal X}}
\def\cY{{\mathcal Y}}

\def\fA{\mathfrak{A}}
%\def\fA1{\mathfrak{A}_1}
%\def\fA2{\mathfrak{A}_2}
\def \fD{\mathfrak{D}}
\def\fB{\mathfrak{B}}
\def\fd{\mathfrak{d}}
\def\ff{\mathfrak{f}}
\def\fH{\mathfrak{H}}
\def\fR{\mathfrak{R}}
\def\fT{\mathfrak{T}}
\def\fS{\mathfrak{S}}
\def\fU{\mathfrak{U}}
\def\fX{\mathfrak{X}}
\def\fx{\mathfrak{x}}

\def\ccA{{\mathscr{A}}}
\def\ccB{{\mathscr{B}}}
\def\ccC{{\mathscr{C}}}
\def\ccD{{\mathscr{D}}}
\def\ccH{{\mathscr{H}}}
\def\ccI{{\mathscr{I}}}
\def\ccL{{\mathscr{L}}}
\def\ccM{{\mathscr{M}}}
\def\ccP{{\mathscr{P}}}
\def\ccQ{\mathscr{Q}}
\def\ccR{\mathscr{R}}
\def\ccS{\mathscr{S}}
\def\ccT{{\mathscr{T}}}
\def\ccW{\mathscr{W}}
\def\ccZ{\mathscr{Z}}


\def\Is{I_\star}
\def\Js{J_\star}
\def\Ks{K_\star}
\def\Ls{L_\star}
\def\ps{p_\star}
\def\mus{\mu_\star}

\def\hp{\hat p}

\def\cAh{\cA_h}
\def\cAs{\cA_\star}
\def\cPb{\cP_\bullet}
\def\Pb{P_\bullet}

\usepackage{enumitem}
\def\rmlabel{\upshape({\itshape \roman*\,})}
\def\RMlabel{\upshape(\Roman*)}
\def\alabel{\upshape({\itshape \alph*\,})}
\def\Alabel{\upshape({\itshape \Alph*\,})}
\def\nlabel{\upshape({\itshape \arabic*\,})}
\renewcommand\labelenumi{(\roman{enumi})}
\renewcommand\theenumi\labelenumi
\def\sb{\mathbf{s}}
\def\Sb{\mathbf{S}}
\def\tb{\mathbf{t}}
\def\Tb{\mathbf{T}}
\def\Ub{\mathbf{U}}

%nina probability
\newcommand{\plog}{\log^{O(1)}}
%\newcommand{\degree}{\operatorname{deg}}
\newcommand\given[1][]{\:#1  \vert  \:}
\newcommand{\Bin}{\textrm{Bin}}
\newcommand{\Unif}{\textrm{Unif}}
\newcommand{\Var}[1]{\textrm{Var} \left[ #1 \right]}
\newcommand{\Cov}[1]{\textrm{Cov} \left( #1 \right)}
\newcommand{\pr}[1]{\mathbb{P} \left[ #1 \right]}
\newcommand{\er}[1]{\mathbb{E} \left[ #1 \right]}
\newcommand{\E}{\mathbb{E} }
\def \Gnp{G_{n,p}}
\newcommand{\tG}{\Gamma}

%nina macros
\def \Qone{Q} %this is denoted Q_1 in Christian's notes 

\def \nst{n_{*}}


% comments
\newcommand{\AM}[1]{\textcolor{purple}{ $\blacktriangleright$\ {\sf AM: #1}\
		$\blacktriangleleft$ }}

\newcommand{\nknov}[1]{\textcolor{green}{ $\blacktriangleright$ NK:} \ {\sf  #1}\
		\textcolor{green}{$\blacktriangleleft$ }} %used for the tiling lemma in the singular case, 



\DeclareMathOperator{\ex}{ex}
\DeclareMathOperator{\codeg}{codeg}
\DeclareMathOperator{\Forb}{Forb}


\let\polishlcross=\l
\def\l{\ifmmode\ell\else\polishlcross\fi}


\makeatletter
\def\moverlay{\mathpalette\mov@rlay}
\def\mov@rlay#1#2{\leavevmode\vtop{   \baselineskip\z@skip \lineskiplimit-\maxdimen
		\ialign{\hfil$\m@th#1##$\hfil\cr#2\crcr}}}
\newcommand{\charfusion}[3][\mathord]{
	#1{\ifx#1\mathop\vphantom{#2}\fi
		\mathpalette\mov@rlay{#2\cr#3}
	}
	\ifx#1\mathop\expandafter\displaylimits\fi}
\makeatother

\newcommand{\dcup}{\charfusion[\mathbin]{\cup}{\cdot}}
\newcommand{\bigdcup}{\charfusion[\mathop]{\bigcup}{\cdot}}



\def\tand{\ \text{and}\ }
\def\qand{\quad\text{and}\quad}
\def\qqand{\qquad\text{and}\qquad}

\newcommand{\vrhup}[1]{\scaleobj{0.6}{\scalerel*{\rightharpoonup}{#1}}}
\newcommand{\nrhup}{\mathord{\scaleobj{0.6}{\scalerel*{\rightharpoonup}{x}}}}
\newcommand{\wrhup}{\scaleobj{0.6}{\scalerel*{\rightharpoonup}{W}}}
\def\vseq#1{\ThisStyle{  \mathord{\vbox{\offinterlineskip\ialign{    \hfil##\hfil\cr
					$\SavedStyle{}_{\smash{\vrhup#1}}$\cr
					\noalign{\kern-0.7\scriptspace}
					$\SavedStyle#1$\cr}}}}}
\def\seq#1{\ThisStyle{  \mathord{\vbox{\offinterlineskip\ialign{    \hfil##\hfil\cr
					$\SavedStyle{}_{\smash{\nrhup}}$\cr
					\noalign{\kern-0.5\scriptspace}
					$\SavedStyle#1$\cr}}}}}
\def\wseq#1{\ThisStyle{  \mathord{\vbox{\offinterlineskip\ialign{    \hfil##\hfil\cr
					$\SavedStyle{}_{\smash{\wrhup#1}}$\cr
					\noalign{\kern-0.7\scriptspace}
					$\SavedStyle#1$\cr}}}}}

\let\vec=\seq

\let\st=\colon
\def\bl{\bigl(}
\def\br{\bigr)}

\def \absorber{\cA}

\def\tinyskip{\vspace{2pt plus 0.7pt minus 0.7pt}}

\let\setminus=\smallsetminus
\let\emptyset=\varnothing
\let\vn=\varnothing

\let\lra=\longrightarrow
\let\to=\lra


\makeatletter
\newcommand{\pushright}[1]{\ifmeasuring@#1\else\omit\hfill$\displaystyle#1$\fi\ignorespaces}
\newcommand{\pushleft}[1]{\ifmeasuring@#1\else\omit$\displaystyle#1$\hfill\fi\ignorespaces}
\makeatother

\let\N=\NN
\let\R=\RR
\let\df=\fd
\let\Hc=\cH
\let\Kc=\cK
\let\Pc=\cP
\let\Jn=J
\let\pst=\ps


\newcommand{\npl}{n^{\ell}p^{\binom \ell 2}}
\newcommand{\Rbinn}{\R^{\binom {[n]}{2}}}

\newcommand{\norm}[1]{\left \|#1 \right \|}
\newcommand{\cutnorm}[1]{\left \|#1 \right \|_\square}

\newcommand{\emb}{\textup{Emb}}

\newcommand{\nk}[1]{\textcolor{orange}{ $\blacktriangleright$\ {\sf NK: #1}\
		$\blacktriangleleft$ }}
\newcommand{\nka}[1]{\textcolor{blue}{ $\blacktriangleright$\ {\sf NK: #1}\
		$\blacktriangleleft$ }} %10 apr 25 -new proof
\newcommand{\nkmay}[1]{\textcolor{blue}{  {\sf NK: #1}\
		}} %10 apr 25 -new proof
\newcommand{\tpt}[1]{\textcolor{purple}{ \ {\sf  #1}\ }}
\def\nkb#1{\textcolor{violet}{#1}} 
\def\nksing#1{\textcolor{blue}{#1}}
\def\sa#1{\textcolor{cyan}{#1}}
\def\nksimpler#1{\textcolor{red}{#1}} 
\newcommand{\todo}[1]{\textcolor{orange}{ $\blacktriangleright$\ {\sf To do: #1}\
		$\blacktriangleleft$ }}

\begin{document}
	
	
	
	
	\keywords{Random graphs, randomly perturbed graphs, thresholds, clique factors}
	\subjclass[2020]{05C80 (primary), 05C35, 05D40 (secondary)}
	
	%Old abstract----------------------------------------------------
	%\begin{abstract}
		%   A \emph{randomly perturbed graph} $G^p = G_\alpha \cup \Gnp$ is obtained by taking a~deterministic $n$-vertex graph $G_\alpha = (V, E)$ with minimum degree $\delta(G)\geq \alpha n$  
		%and adding the edges of the binomial random graph $\Gnp$ defined on the same vertex set $V$. For which value $p$ (depending on $\alpha$) does the graph $G^p$ contain a $K_r$-factor (a spanning collection of vertex-disjoint $K_r$-copies) with high probability?
		
		%   The  order of magnitude of the \textit{minimal} value of $p$ has been determined  whenever $\alpha \neq 1- \frac{s}{r}$ for an integer $s$ (see Han, Morris, and Treglown [RSA, 2021] and Balogh, Treglown, and Wagner [CPC, 2019]).
		
		%   We establish the minimal probability $p_s$ (up to a constant factor) for all values of $\alpha = 1-\frac{s}{r} \leq \frac 12$, and show that the \textit{threshold} exhibits a polynomial jump at $\alpha = 1-\frac{s}{r}$ compared to the surrounding intervals. An extremal example $G_{\alpha}$ which shows that $p_s$ is optimal up to a constant factor differs from the previous (usually multipartite) examples in containing a pseudorandom induced subgraph.
	%\end{abstract} 
    
	%New abstract - for RS&A----------------------------------

    \begin{abstract}
        A \emph{randomly perturbed graph} $G^p = G_\alpha \cup \Gnp$ is obtained by taking a~deterministic $n$-vertex graph $G_\alpha = (V, E)$ with minimum degree $\delta(G)\geq \alpha n$ and adding the edges of the binomial random graph $\Gnp$ defined on the same vertex set $V$. For which value $p$ (depending on $\alpha$) does the graph $G^p$ contain a $K_r$-factor -- a spanning collection of vertex-disjoint copies of $K_r$ -- with high probability?
        The  order of magnitude of the \textit{minimum} such $p$ is known whenever $\alpha \neq 1- \frac{s}{r}$ for an integer $s$ (see Han, Morris, and Treglown [RSA, 2021] and Balogh, Treglown, and Wagner [CPC, 2019]).
        In earlier work, the first three authors determined this threshold probability $p_s$ up to a constant factor for all values of $\alpha = 1-\frac{s}{r}\leq \frac 12$. Here, we complete the picture by establishing $p_s$ in the remaining case $\alpha > \frac12$. 

        A key ingredient in our approach is an extremal result of independent interest: we prove a fractional stability version of a \emph{tiling} theorem due to Shokoufandeh and Zhao.
        
        %We establish the minimal probability $p_s$ (up to a constant factor) for all values of $\alpha = 1-\frac{s}{r} \leq \frac 12$, and show that the \textit{threshold} exhibits a polynomial jump at $\alpha = 1-\frac{s}{r}$ compared to the surrounding intervals. An extremal example $G_{\alpha}$ which shows that $p_s$ is optimal up to a constant factor differs from the previous (usually multipartite) examples in containing a pseudorandom induced subgraph.
    \end{abstract}

    
	\maketitle



%    \tableofcontents
	
	%************************* Introduction ****************************
	
	\section{Introduction}
	
	
	
%	\nk{How I think we should advertise the two papers: the first paper discovers the phenomenon, proves the lower bound and develops some of the machinery. The present paper relies on an important advance in extremal combinatorics to prove the full thing.}


    Randomly perturbed graphs have been introduced by Bohman, Frieze and Martin~\cite{bfm03,bfm04}. Since their introduction, they have been studied extensively because they interpolate between purely random graphs and highly structured deterministic graphs, providing a~useful framework for understanding fundamental graph properties such as Hamiltonicity and the existence of other spanning structures (see, for example~\cite{adrrs21,balog16,bhkm19,bmpp20,jk20,kks17,kst06}). In this paper, we further advance the study of clique factors in randomly perturbed graphs.
    
    Formally, a \emph{randomly perturbed graph} $G^p = G_\alpha \cup \Gnp$ is obtained by taking a deterministic $n$-vertex graph $G_\alpha = (V, E)$ satisfying minimum degree condition $\delta(G)\geq \alpha n$ and adding to it the edges of the binomial random graph $\Gnp$ defined on the same vertex set $V$. Given an increasing graph property $\mathcal{P}$, for any $\alpha\in(0,1)$ one would like to establish the \emph{threshold} density $p=p(n)$ for this property, that is the lowest density\footnote{As usual, this threshold is not uniquely defined and we are only interested in determining the threshold up to a constant factor for our specific property.} $p$ which, \emph{with high probability} (w.h.p.), ensures that $G_\alpha \cup \Gnp$ satisfies $\mathcal{P}$. Here, as usual, w.h.p.~means with probability tending to 1 when $n\rightarrow\infty$.

    A \emph{$K_r$-factor} in a graph $G$ is a collection of vertex disjoint copies of the complete graph $K_r$ covering all the vertices of $G$. Throughout the paper, we will assume that the number of vertices of the \emph{host graph} $G_\alpha$, denoted by $n$, is divisible by $r$. Depending on $\alpha >0$, we are interested in  determining the threshold density for the existence of a $K_r$-factor in $G_\alpha \cup \Gnp$. Both `extreme cases' ($p=0$ and $\alpha=0$) constitute classical and difficult results. Firstly, resolving a~conjecture of Erd\H os, Hajnal and Szemer\'edi \cite{Hajnal1970} have shown that for sufficiently large $n$, any $n$-vertex graph $G$ with $\delta(G) \geq \left(1-\frac{1}{r} \right)n$ contains a $K_r$-factor. Hence, for $\alpha \geq 1- \frac 1r$, no random edges are needed. On the other hand, a celebrated result of Johansson, Kahn and Vu~\cite{jkv08} states that $p = Cn^{-2/r}(\log n)^{2/(r(r-1))}$ is the threshold for the containment of a $K_r$-factor in $\Gnp$. Hence, this is the desired threshold for the case $\alpha=0$.

    Considering the $K_r$-factor problem in the model of randomly perturbed graphs reveals surprising phenomena, and requires an interesting synergy between combinatorial and probabilistic techniques. The threshold for the appearance of a $K_r$-factor in the randomly perturbed graph $G_\alpha \cup \Gnp$ has been determined for almost all values of $\alpha$: first for $\alpha \in \left(0,\frac{1}{r}\right)$ by Balogh, Treglown, and Wagner~\cite{btw18}, and subsequently for the remaining intervals of the form $\alpha \in \left(1-\frac{s}{r}, 1-\frac{s-1}{r}\right)$ by Han, Morris, and Treglown~\cite{hmt21}.
	
	\begin{thm}
		[Clique-factors in randomly perturbed graphs~\cite{btw18}, \cite{hmt21}]
		\label{thm:hmt}
		For all integers $1 < s \leq r$ and any $\alpha \in \left(1-\frac{s}{r}, 1-\frac{s-1}{r} \right)$, there exists $C$ such that the following holds. For any $n$-vertex graph $G_\alpha$ with minimum degree at least $\alpha n$, w.h.p.\ $G_\alpha \cup \Gnp$ with $p = Cn^{-2/s}$ contains a $K_r$-factor.
	\end{thm}

    The density $p = Cn^{-2/s}$ is necessary, and understanding the extremal example $G_\alpha$ sheds light on the problem. 
    In particular, the extremal example demonstrates that for $\alpha \in \left(1-\frac{s}{r}, 1-\frac{s-1}{r} \right)$, the bottleneck for the threshold density $p$ is the existence of $\Omega(n)$ copies of $K_{s}$ in the random graph $\Gnp$. Namely, let $G_{\alpha}$ consist of an independent set $A$ of size $(1-\alpha) n > \frac {(s-1)n}{r}$, and all the remaining edges (see Figure~\ref{fig:extremal_1}). If $G_{\alpha} \cup \Gnp$ is to contain a $K_r$-factor $\cF$, then $\cF$ cannot contain only cliques $K$ with $|K \cap A| \leq s-1$; in fact, at least $\Omega(n)$ cliques from $\cF$ have to intersect $A$ in at least $s$ vertices. Thus, $\Gnp[A]$ has to contain $\Omega(n)$ copies of $K_{s}$, which in turn requires edge density $p = \Omega (n^{-2/s})$.

    \begin{figure}
		\begin{tikzpicture}[scale=0.8]
			\centering
			\filldraw[color=lightgray] (0,1.3) -- (4,1.6) -- (4,-1.6) -- (0,-1.3) -- (0,1.3);
			\draw[fill=white] (0,0) ellipse (1.1cm and 1.9cm);
			\draw[fill=lightgray] (4,0) ellipse (1.3cm and 2.1cm);
			\node at (0,-2.2) [below] {$A$};
			\node at (4,-2.2) [below] {$V\setminus A$};
		\end{tikzpicture}
		\caption{\label{fig:extremal_1} The extremal construction $G_{\alpha}$ which, crucially, contains an independent set of size $(1-\alpha)n$.}
	\end{figure}

    What about the \textit{transition points}, that is $\alpha$ of the form $\alpha = 1-\frac sr$\,?
    B{\"o}ttcher, Parczyk, Sgueglia, and Skokan~\cite{bpss23} considered the only interesting transition point for $r=3$ (the triangle factor), $\alpha = \frac 13$, showing that the threshold for a $K_3$-factor in $G_{\frac 13} \cup \Gnp$ is $p = \frac{C \log n}{n}$. Instructed by this result, the example from Figure~\ref{fig:extremal_1} and the Johansson--Kahn--Vu Theorem, one might guess the following: as $\alpha$ decreases from $\alpha > 1-\frac{s}{r}$ to precisely $\alpha = 1- \frac sr$, the threshold density `gains' a~polylogarithmic factor.

 %    \begin{thm}
	% 	[Triangle-factors in randomly perturbed graphs~\cite{bpss23}]
	% 	There exists $C$ such that for any $n$-vertex graph $G$ with minimum degree at least $\delta(G)\geq n/3$, w.h.p.\ $G \cup \Gnp$ with $p = C\log n/n$ contains a $K_3$-factor.
	% \end{thm}

    %This was the first result on clique-factors in randomly perturbed graphs when the minimum degree condition is of the form $\alpha = 1 - \frac{s}{r}$ for some integers $s$ and $r$. 
    %We call this case the \emph{transition point}. It is worth noting that in case of all clique-factors, there is a~polynomial jump in threshold probability when we switch from one interval determining $\alpha$, to the consecutive one. The result in~\cite{bpss23} states that for triangle-factors ($r=3$), at the only interesting transition point $\alpha=1/3$, we need an extra $\log n$ factor. This would suggest a similar behaviour in other transition points.
    
    However, Böttcher et al.~\cite{bpss23} have already pointed out a rather surprising polynomial jump in the threshold probability when $r=4$ and $s=3$ ($\alpha=1/4$), and asked about the correct threshold. In~\cite{akr24}, the first three authors have pinpointed this threshold, and generalised it to an infinite family of transition points. The table below summarises the thresholds for $r=4$.

    \smallskip
    
    \begin{center}
           \begin{tabular}{| c | c |c | >{\columncolor{gray!15}}c | c|c| c | c |}
    \hline
    & & & & & & & \\[-3ex]
    $\alpha$ & 0 & $\left(0, \frac{1}{4}\right)$ & $\frac{1}{4}$ & $\left(\frac{1}{4}, \frac{1}{2}\right)$ & $\frac{1}{2}$ & $\left(\frac{1}{2}, \frac{3}{4}\right)$ & $\left[\frac{3}{4}, 1\right]$ \\ [0.5ex]
    \hline
    & & & & & & & \\[-3ex]
    $K_4\text{-factor}$ threshold & $n^{-\frac{2}{4}} \log^{\frac{1}{6}} n$ & $n^{-\frac{2}{4}}$ & $n^{-\frac{3}{5}}$ & $n^{-\frac{2}{3}}$ & $n^{-1} \log n$ & $n^{-1}$ & $0$ \\ [0.4ex]
    \hline
  \end{tabular}
      \end{center}

    \smallskip
    
    %In \cite{akr24}, we have determined this threshold to be of order $n^{-3/5}$, so strictly between $n^{-2/3}$ and $n^{-1/2}$ which is the order of magnitude of the threshold probability for $\alpha < 1/4$ and $\alpha > 1/4$, given by Theorem~\ref{thm:hmt}. What is more, we have addressed the problem of determining the threshold for $K_r$-factors in the randomly perturbed graphs for general $r$ and at the transition points.

    Specifically, in~\cite{akr24}, the transition points $\alpha = 1-\frac sr$ were considered for all $r \geq 3$ and $1<s<r$. The threshold was determined whenever $s \geq r/2$ (equivalently, $\alpha \leq 1/2$), and a~lower bound was found for all values of $s$. 

    The extremal construction establishing the lower bound from~\cite{akr24} is a partly pseudorandom host graph $G_\alpha$. This is a completely new approach, as in most results on the randomly perturbed graphs the 0-statements are proved using some highly structured graphs such as complete multipartite graphs. Rather than describing the construction here, we refer the reader to~\cite{akr24}.
    
    Before stating the results, we introduce the crucial functions. Let 
	\(\varphi(s) = \frac{2s}{(s-1)(s+2)}\)
	and 
	\[p_s = p_s(n) = \begin{cases} 
		\frac{\log n}{n} & \text{ if } s=2 \\
		n^{-\varphi(s)} & \text{ if } s\geq 3.
	\end{cases} \]

    The following theorem was proved in~\cite{akr24}. For $r=3$ and $s=2$, both statements were earlier established by B{\"o}ttcher, Parczyk, Sgueglia, and Skokan~\cite{bpss23}.
    \begin{thm}[Clique-factors at the transition points \cite{akr24}]
    \label{thm:main_in_AKR}
		Let $r$ and $s$ be integers with $1 < s <r$, and let $n$ be divisible by $r$.
        \begin{enumerate}
            \item For any $\eps>0$, there is a~constant $c>0$ and a graph $G$ with minimum degree at least $\left(1 - \frac{s}{r} \right)n$ such that taking $p = cp_s(n)$ the probability that the graph $G^p = G \cup \Gnp$ has a~$K_r$-factor is less than $\eps$.
            \item If $s \geq r/2$, there exists $C>0$ such that for $p = Cp_s(n)$ the following holds. If $G$ is an $n$-vertex graph with minimum degree at least $\left(1 - \frac{s}{r} \right)n$, then  w.h.p.\,the graph $G^p = G \cup \Gnp$ has a~$K_r$-factor.
        \end{enumerate}
	\end{thm}
	 
	In~\cite{akr24} it was conjectured that statement (ii) extends to all values of $s$. In the present paper, we confirm this conjecture, also answering the questions raised in~\cite{bpss23,hmt21}. 

    \begin{thm}\label{t:main}
		For any integers $r$ and $s$ with $1 < s < r$, there exists $C>0$ such that the following holds. Let $G$ be an $n$-vertex graph with minimum degree at least $\left(1 - \frac{s}{r} \right)n$. For $n$ divisible by $r$ and $p = Cp_s(n)$, w.h.p.\,the graph $G^p = G \cup \Gnp$ has a~$K_r$-factor.
	\end{thm}

    In the remainder of the introduction, we highlight an extremal result that is of independent interest, and other proof ideas that require treating the case $s \geq r/2$ in a separate article.

    \subsection*{A Dirac--type problem}
       
    How does the Hajnal--Szemere\'edi Theorem generalise to $H$-factors, for a given graph $H$\,? Shokoufandeh and Zhao~\cite{sz03} proved that if $\delta(G) \geq \left(1-\frac{1}{\chi_{cr}(H)} \right)n$, then $G$ contains a collection of vertex-disjoint copies of $H$ covering all but a constant number of vertices; here, $\chi_{cr}(H)$ is the \textit{critical chromatic number} of $H$. This result was previously conjectured by Koml\'os~\cite{komlos00}. The minimum degree condition is optimal for every graph $H$, and the presence of a constant number of uncovered vertices is unavoidable (see K\"uhn, Osthus~\cite{ko09}).
        
        
    We show the \textit{fractional stability version} of~\cite{sz03}. To be precise, a \emph{weighted graph} is a pair $(F, w_F)$, where $F$ is a~graph and $w_F: V(F) \to \R_{\geq 0}$ is a~weight function.
    For a~family of weighted graphs $\cF$, an \emph{$\cF$-packing} in $G$ is a~collection of graphs $(F_i)_{i \in I}$, where $F_i \in \cF$, with corresponding embeddings $(f_i)_{i \in I}$ into $G$, such that for any $x \in V(G)$, 
    $$\sum_{i \in I: Im (f_i) \ni x} w_{F_i}(f_i^{-1}(x)) \leq 1.$$
    In other words, for each $x \in V(G)$ the total weight in $(F_i)_{i\in I}$ of all vertices mapped to $x$ is at most 1. For now, we only consider the case $\cF=\{H\}$, and thus refer to $H$-packings\footnote{The terms fractional packing, fractional tiling and partial fractional factor are also often used in the literature.}.
    For a weighted graph $H$ and $\phi > 0 $, denote by $\phi\ltimes H$ the weighted graph obtained by multiplying all the vertex weights of $H$ by $\phi$.

    The following theorem states that if $G$ is `sufficiently far' from the extremal example in the Shokoufandeh--Zhao theorem, then an $H$-packing with `constant residue' can be found, even under a relaxed minimum-degree assumption.

        \begin{thm}\label{t:sz-stability}
        Let $H$ be a graph with critical chromatic number $\chi_{cr}(H)$.  There are integers $b,C>0$ such that for all $\gamma>0$, there exists a constant $\delta >0$ such that the following holds for sufficiently large $n$. If $G$ is an $n$-vertex graph with $\delta(G)\geq \left(1-\frac{1}{\chi_{cr(H)}}-\delta\right)n$ in which every vertex set of size $\frac{n}{\chi_{cr}(H)}$ spans at least $\gamma n^2$ edges, then $G$ contains a $(b^{-1} \ltimes H)$-packing $(f_i)_{i \in I}$ with 
         $$n- |I||V(H)|b^{-1}\leq C.$$\end{thm}
        An asymptotic version of this statement was shown by Hladk\'y, Hu and Piguet~\cite[Theorem 1.3]{hhp19}. However, in the result of~\cite{hhp19}, the $H$-copies can cover all but $O(\delta n)$ vertices (where, crucially, $\delta$ also appears in the minimum-degree assumption on $G$).

        Theorem~\ref{t:sz-stability} is the key to our proof, and it might have further applications in extremal combinatorics. It is proved in Section~\ref{sec:sz-stability}.
        
          %Let $r = ms+t$. (Here $s$ may not be an integer, that does not matter for the proof). Then the weighted graph $T$ (see Section 4) contains a $\{q \ltimes H\}$ packing for some rational $q$. (This is not hard to see, but I also think it's shown in \cite{sz03}, except that they call $T$ a "bottle graph", and use different notation for $m,r,s,t$.)

        
         
         %. Let $H(m, s, t)$ be the complete multipartite graph with $m$ classes of size $s$ and one class of size $t$, where $t \leq s$. By definition, $\chi_{cr}(H(m,s,t)) = \frac{r}{s}$. (The weighted graph $T$ from Section 4 is exactly $r^{-1} \ltimes H(m,s,t)$.

     %    \textbf{Fact 1.} Let $H$ be a graph with $\chi_{cr}(H) = \frac rs$, and chromatic number $m$ Then $H(m,s,t)$ has a $\{q \ltimes H\}$-packing with residue 0.


   %     So I'd like to check if it implies the full version. For that, we need to answer two questions.

   %     Let $T_1$ be the graph with $m$ classes of size $s_1, \ldots, s_m \geq t$, and one class of size $t$, with a total of $ms+t$ vertices, so also $\chi_{cr}(H_1) = s/r$.

    %    1) Is it true that every graph $H$ with $\chi_{cr}(H) = s/r$ embeds into $H_1$? (I think this is true, but pls check, and the answer might be in~\cite{sz03}).

    %    **new question**
  %      2) Importantly, does $Q^{(h)}_m$ have a fractional $T_1$ factor, for every $s_1, ..., s_m \geq t$?


    
    \subsection*{Proof ideas} 
    Let us mention a few proof ingredients, which also justifies treating the cases $s \geq r/2$ and $s< r/2$ separately. In short, the present proof requires substantial new results and ideas, including Theorem~\ref{t:sz-stability}. On the other hand, the present proof does not imply the result of~\cite{akr24}. Finally, part of the article~\cite{akr24} was devoted to developing tools for constructing `spread embeddings' in regular partitions and applying the recently proved (fractional) Kahn--Kalai conjecture~\cite{fknp21,pp24}. We now use these tools (Lemma~\ref{l:Z-conf.-Kt-factor}) as a~black box.

    Recall that $r = ms+t$ with $t<s$, and that we have to prove the statement for every deterministic graph $G$ with $\delta(G) \geq 1 - \frac sr$. We distinguish a few cases depending on the structure of $G$. Note that the extremal example is \textit{close to} the graph in Figure~\ref{fig:extremal_1}, where $|A| = \frac{sn}{r}-o(n)$, and $G[A]$ induces a sparse random-like graph~\cite{akr24}. %\sa{This is not true. In case $s\geq 3$, the extremal example is the complete bipartite graph with extra edges (pseudorandom part) added in one of the parts to get the min degree condition.}

    Most of the work and innovation lies in the \textit{non-extremal case}, where $G$ does not contain an almost independent set of size roughly $\frac{sn}{r}$. In this case, we apply Szemer\'edi's Regularity Lemma~\cite{szemeredi78}, and wish to find a useful substructure in the \textit{reduced} graph. The useful structure is essentially given by Theorem~\ref{t:sz-stability}.

    Even upon finding this structure, we are not done using standard techniques. Namely, a~lot of work is devoted to covering the leftover vertices and ensuring exactly the desired part sizes in the regular partition. This is also a major challenge in~\cite{bpss23,akr24}, and whenever exact minimum degree conditions are needed. The absorption strategy for resolving this issue is different from~\cite{akr24}, and it relies heavily on the fact that, since $\delta(G) > \frac 12$, the graph is connected, and any pair of vertices has many common neighbours (see, e.g.~Definition~\ref{def:absorber}).

    In the extremal case, the host graph $G$ contains an almost independent set $A$ of size roughly $\frac{sn}{r}$, and, due to the minimum degree condition, most edges between $A$ and $V\setminus A$ are present. Therefore, in constructing a $K_r$-factor, most of the copies of $K_r$ use exactly $s$ vertices from $A$, and the respective sets should span cliques in the random graph $\Gnp$. However, a number of copies of $K_{s+1}$ in $G[A]$ may also be required, which is the bottleneck for our construction and which in turn determines the threshold density $p_s$. 

     An additional difficulty in the present paper are \textit{intermediate cases}. For instance, suppose that $G$ consists of an independent set $A$ of size roughly $\frac{sn}{r}$, but the remaining graph $G[V \setminus A]$ does not have a particular structure. The main idea for resolving this case is an inductive approach: recalling that $G^p = G \cup G_{n,p}$, we can find a $K_{r-s}$-factor in $G^p[V \setminus A]$ (which roughly satisfies the \textit{correct} minimum-degree condition and is \textit{non-extremal}), and a $K_s$-factor in $G^p[A]$. The aim and the difficulty is then to construct those factors in a way that they can be combined into a $K_r$-factor. An additional hurdle is that $|A|$ can differ from $\frac{sn}{r}$. Most of these issues are addressed in Section~\ref{sec:pf-of-main}. In particular, a~stronger conclusion is needed in the \textit{non-extremal} case (Theorem~\ref{t:non-extremal-conforming}).
    
    Finally, a careful reader may notice that the number of possible cases is roughly $\lfloor r/s \rfloor$, and the case distinction (Lemma~\ref{l:4.3chr}) is rather refined.
    
      
    \subsection*{Paper structure} 
    The article is organised as follows. Preliminary definitions and results are stated in Section~\ref{sec:prelim}. Section~\ref{sec:pf-of-main} contains the main lemma statements which roughly correspond to the outlined case distinction, as well as a deduction of the main result, Theorem~\ref{t:main}. The afore-mentioned lemmas are proved in Sections~\ref{sec:non-extremal} and~\ref{sec:partition-lemma}. Section~\ref{sec:sz-stability} contains the proof of Theorem~\ref{t:sz-stability}, and further extremal consequences of this result.	
	
%************************* Preliminaries ****************************
	
    \section{Preliminaries}
	   \label{sec:prelim}
    Let $m, s, t$ be integers such that $r = ms + t$ with $1\leq t \leq s$. Throughout the paper, we refer to the case $s = t = r/(m+1)$ as the \textit{singular} case, and in this case some steps are treated separately. We will often use an auxiliary parameter $g$ defined as
    \[ g = \begin{cases}
        0 & \text{if} \quad t<s; \\
        1 & \text{if} \quad t=s.
    \end{cases}\]

    We use the same notation as in \cite{akr24}. In particular, for any set $X$ and an integer $\ell$, by ${X \choose \ell}$ we denote the family of all $\ell$-element subsets of $X$.
    For a collection of graphs $\cC$, by $V(\cC)$ we denote the set of vertices of the graphs in $\cC$. Let $e_G(X,Y)$ be the number of edges in $G$ whose one endpoint lies in $X$ and the other in $Y$. If $X\cap Y\neq\emptyset$, then any edge $uv$ with $u,v\in X\cap Y$ is counted twice in $e_G(X,Y)$.

    \begin{definition}
        For a~vertex set $Z$, we call a collection $\cK$ (usually a $K_r$-factor) \textit{$Z$-conforming} if for all $K \in \cK$, $K$ contains at most one element of $Z$.
    \end{definition}
    
    
    %For a graph $G=(V,E)$ we denote by $V(G)$ and $E(G)$ the set of vertices and the set of edges of $G$, and by $v(G)$ and $e(G)$ the sizes of the respective sets. For a subset $S\subset V$, by $e_G(S)$ we denote the number of edges spanned in $G$ by vertices of $S$, and by $N_G(S)$ we denote the set of common neighbours in graph $G$ of vertices in $S$. Next, $G(V,p)$ denotes the usual binomial random graph defined specifically on the vertex set $V$, and $G_{n,m,p}$ is the random bipartite graph with partition classes of size $n$ and $m$, respectively. Moreover, $G[V_1,V_2]$ is the bipartite subgraph of $G$ induced by partition classes $V_1$ and $V_2$. 
	
    %For a family of graphs $\mathcal{P}$, a \emph{$\mathcal{P}$-factor} in a graph $G$ is a collection $\mathcal{F}$ of pairwise vertex-disjoint subgraphs of $G$, each isomorphic to an element of $\mathcal{P}$, which together span $V(G)$. Sometimes we refer to a collection of such pairwise vertex-disjoint subgraphs of $G$ as a \textit{partial factor}.
	
%************************* Probabilistic bounds ****************************
	
    \subsection{Probabilistic bounds}

    The following lemma is proved in \cite{akr24} for $s\geq 3$, and for $s=2$ follows from classic results on the independence number of $\Gnp$ (found, for instance, in~\cite[Theorem 7.4]{jlr00}).

    \begin{lemma}\label{l:ks-in-lin-sized}
        For $s \geq 2$ and $\eps >0$, let  $p = n^{-2/s+\eps}$. Then, w.h.p.\,any subset $S \subset V(\Gnp)$ with $|S| \geq \frac{n}{\log^2 n}$ contains a copy of $K_s$.
    \end{lemma}
    %why did i not state this with n/ \sqrt{log n}? because the assumption is p= n^{-2/s+eps} sto je previse za slucaj s=2.

    The next lemma is a well-known consequence of Janson's Inequality.
    
    \begin{lemma}
	\label{l:akr:ks-s-everywhere}
        Let  $s$ be an integer, and $\nu >0$. Let $V$ be an $n$-vertex set, and let $\cC$ be a~family of $s$-element subsets of $V$ with $|\cC|\geq \nu n^{s}$. For $p = Cn^{-2/s}$ and $C$ sufficiently large, the probability that no member of $\cC$ induces a copy of $K_s$ in $G(V,p)$ is at most $e^{-c n^2 p}$, where $c$ is a constant depending on $s$.
    \end{lemma}

    The following lemma is a generalisation of the previous one, in which a slightly more complex subgraph is sought in $\Gnp$.
    
    \begin{lemma}
        \label{l:many_ks_and_one_ks-ks}
        Let $s,q\ge 2$ be integers, and let $\nu>0$.  
Define $F_q$ to be the graph on vertex set $[qs-1]$ in which the sets
$\{1,\dots,s\}$ and $\{s,\dots,2s-1\}$ induce copies of $K_s$, and for each
$i=3,\dots,q$, the set $\{(i-1)s,\dots,is-1\}$ induces a copy of $K_s$,
with no other edges. (Thus $F_q$ consists of $q$ copies of $K_s$, exactly two
sharing exactly one vertex, and all others pairwise vertex-disjoint.)

Let $V$ be an $n$-vertex set and let $\cC\subset V^{qs-1}$ satisfy
$|\cC|\ge \nu n^{qs-1}$. For $p=Cn^{-2/s}$ and $C$ sufficiently large, the probability that no member of $\cC$ induces a copy of $F_q$ in $G(V,p)$ is at most $e^{-c n^2 p}$, where $c=c(s,q)>0$ is a constant.
    \end{lemma}
    \begin{proof}
        [Proof sketch.] We apply induction on $q$, noting that $F_{q+1}$ is isomorphic to the disjoint union of $F_q$ and $K_s$.
        
        Suppose first that $q=2$. For this case, the lemma was essentially proved in~\cite[Lemma 9.2]{akr24}. It is an application of Janson's Inequality, and even though the setup in~\cite{akr24} is slightly more specific, the calculation extends straightforwardly to the present statement.

        Now assume that the Lemma holds for $q$, and let $\cC \subset V^{(q+1)s-1}$ be as in the statement. Using double exposure, we sample $\Gnp$ as $G_1\cup G_2$, where $G_1, G_2 \sim G_{n,p_1}$ are independent and $1-p = (1-p_1)^2$. Applying a standard pruning argument and Lemma~\ref{l:akr:ks-s-everywhere}, w.h.p.~$G_1$ contains an $s$-tuple $\textbf{v}$ %$\textbf{v}=v_{(q-1)s}, \ldots, v_{qs-1}$ 
        which induces a $K_s$-copy and extends to at least $\nu'n^{qs-1}$ members of $\cC$. Then we can apply the induction hypothesis to this subcollection of $\cC$ extending $\textbf{v}$ and the graph $G_2$, finding a copy of $F_{q+1}$ in $G_1 \cup G_2$, as desired.
    \end{proof}

    The next lemma, proved in \cite{akr24}, is the key ingredient in the main proof and turns out to be the bottleneck for the threshold probability $p_s$ when $s\geq 3$.
	
     \begin{lemma}[\cite{akr24},        Lemma~4.3]\label{l:5.5chr}
        For all $s\geq 2$, there are constants $\gamma, C>0$ such that the following holds. Let $G$ be an $n$-vertex graph with minimum degree $h\leq \gamma n$. W.h.p.\,the graph $G^p = G \cup \Gnp$, with $p = Cp_s(n)$, contains a collection of $h$ vertex-disjoint copies of $K_{s+1}$. 
    \end{lemma}

    In the current paper we will need the following $Z$-conforming strengthening of Lemma~\ref{l:5.5chr}.
    
    \begin{lemma}\label{l:5.5chr-stronger}
        For all $s\geq 2$ and $\delta\ll\beta\ll 1/s$, there are constants $\gamma, C>0$ such that the following holds. Let $G$ be an $n$-vertex graph with minimum degree $h\leq \gamma n$. Let $Z\subset V(G)$ be such that\footnote{In the application, these will be vertices with degree less than $|A_i|-cn$ into some part $A_i$, where $c$ is some small constant like $c <\beta r^{-3}$. $\beta$ is the constant from the partitioning lemma.} $|Z| \leq \delta n$, and suppose that for $v \in Z$, $\deg_G(v) \geq \beta n$. W.h.p.\,the graph $G^p = G \cup \Gnp$, with $p = Cp_s(n)$, contains a collection of $h$ vertex-disjoint copies of $K_{s+1}$ such that each of those copies contains at most one vertex of $Z$.
    \end{lemma}
    
    \begin{proof}
        Assume first that $|Z|\geq h$. Then we may greedily find a collection of vertex-disjoint copies of $K_{s+1}$, each using exactly one vertex from $Z$. Indeed, we find copies of $K_{s+1}$ one-by-one. At the point of processing a vertex $v\in Z$, its number of neighbours in $G$ distinct from vertices in all previous $K_{s+1}$-copies is at least 
        \begin{equation}
                \beta n-|Z|(s+1)\geq n\left(\beta-(s+1)\delta\right)\geq \frac{\beta n}{2}
            \label{eq:Z-min-degree}
        \end{equation}
        for sufficiently small $\delta$. Thus, by Lemma~\ref{l:akr:ks-s-everywhere}, the probability that there is no $K_s$-copy in $N_G(v)\setminus Z$, vertex-disjoint from all previous copies, is at most $e^{-cn}$ for a sufficiently small constant $c>0$. Taking the union bound, the probability that we will fail at any stage of this process is $o(1)$.
        
        Suppose now that $|Z|<h$. Note that $G-Z$ has minimum degree at least $h-|Z|\leq\gamma n$, so we may apply Lemma~\ref{l:5.5chr} to $G-Z$, which yields a collection $\cC$ of $h-|Z|$ vertex-disjoint $K_{s+1}$-copies that are disjoint from $Z$. Now, as in the above argument (noting that $|V(\cC)|\leq \gamma (s+1)n\leq \beta n/4$, so a bound slightly weaker than \eqref{eq:Z-min-degree} holds), this collection can be greedily extended to a collection of size $h$.
    \end{proof}

    \begin{lemma}\label{l:Z-conf.-Kt-factor}
        For all $s\geq t\geq 2$ and $\delta \ll \beta \ll 1/s$, there is a constant $C>0$ such that the following holds. Let $G$ be an $n$-vertex graph, with $n$ divisible by $t$, and $Z\subset V(G)$ be a~subset with $|Z| \leq \delta n$ such that for any vertex $v\in Z$, $\deg_G(v) \geq \beta n$. W.h.p.~the graph $G^p=G\cup \Gnp$, with $p=Cp_s(n)$, has a $K_t$-factor in which each clique contains at most one vertex from $Z$.
    \end{lemma}

    \begin{proof}
        Using the same approach as in the proof of Lemma~\ref{l:5.5chr-stronger}, we first greedily find a~collection $\cC$ of vertex-disjoint $K_t$-copies, each using exactly one vertex from $Z$. We then cover the remaining vertices in $V(G)\setminus V(\cC)$ by a $K_t$-factor using only the edges in $\Gnp$. We can do this since $p$ is strictly above the threshold for a $K_t$-factor (see \cite{jkv08}, Theorem~2.1).
    \end{proof}	
	
%************************* Regularity ****************************

    \subsection{Regularity}\label{sec:prelim-reg}
	
    %Some of our proofs rely heavily on the regularity method. One phenomenon that our paper, as well as~\cite{psss22}, demonstrates, is that the concept of regular pairs works very well in conjunction with the machinery of spread measures introduced in Section~\ref{sec:prelim-kk}.

    In this section, we recall some standard terminology and results regarding regularity.  
    Let $G$ be a graph, and let $V_1, V_2 \subset V(G)$ be disjoint subsets of the vertices of $G$. For non-empty sets $X_1\subseteq V_1$, $X_2 \subseteq V_2$, we define the \emph{density of $G[X_1,X_2]$} as $d_G(X_1, X_2)\coloneq \tfrac{e_G(X_1,X_2)}{|X_1||X_2|}$.
    Given $\eps, d > 0$, we say that a pair $(V_1,V_2)$ is \emph{$(\eps,d)$-regular} in $G$ if for $i = 1,2$ and for all sets $X_i \subseteq V_i$  with $|X_i|\geq \eps |V_i|$, we have $| d_G(X_1,X_2)-d| < \eps$. We say that a pair $(V_1,V_2)$ is $(\eps,d^+)$-regular if it is $(\eps,d')$-regular for some $d'\geq d$. We often omit the subscript $G$ if it is clear from the context.

    We will use the following standard lemmas, which can be found, for example, in~\cite{abcd22,ks95}. We begin by recalling the Slicing Lemma.
	
    \begin{lemma}[\cite{abcd22}, Lemma~2.7]\label{l:slicing}
        Let $0 < \eps < \beta$, $d \leq 1$ and let $(V_1,V_2)$ be an $(\eps,d)$-regular pair. Then for any pair $(U_1, U_2)$ such that for $i=1,2$, $U_i\subset V_i$ with $|U_i| \geq \beta |V_i|$, the pair $(U_1,U_2)$ is $(\eps', d')$-regular with $\eps' = \max\{\eps/\beta, 2\eps \}$ and some $d'>0$ satisfying $|d' -d|\leq \eps$.
    \end{lemma}

    The following lemma is similar. Although the proof is standard, we have not found a~self-contained statement in the literature, so we prove it here.

    \begin{lemma}\label{l:slicing-adding}
        Let $0<\xi<\eps \ll 1$, $d\leq 1$ and let $(U_1, U_2)$ be an $(\eps,d)$-regular pair. Then for any pair $(V_1,V_2)$ such that for $i=1,2$, $U_i\subset V_i$ with $|V_i\setminus U_i|\leq \xi |U_i|$, the pair $(V_1,V_2)$ is $(\eps', d')$-regular with $\eps'= \max\{\xi/\eps, 6\eps\}$ and some $d'>0$ satisfying $|d'-d|<3\eps $.
    \end{lemma}
        
            
    \begin{proof}
        Let $A_i\subseteq V_i$ with $|A_i|\geq \eps' |V_i|$. It is easy to see that $|A_i\cap U_i|\geq (\eps'-\xi)|U_i|\geq \eps |U_i|$ and $|A_i\cap U_i| \geq |A_i|-|V_i\setminus U_i|\geq |A_i|-\xi|U_i|\geq|A_i|-\xi|V_i|\geq |A_i|(1-\xi/\eps')\geq |A_i|(1-\eps).$ 
        We have 
        \begin{align*}
            e_G(A_1,A_2)&\geq e_G(A_1\cap U_1, A_2 \cap U_2) \geq d(1-\eps)|A_1\cap U_1||A_2\cap U_2|\\
            &\geq d (1-\eps)^3|A_1||A_2| \geq (d-3\eps)|A_1||A_2|
        \end{align*}
        and, since $|A_i\setminus U_i|\leq \xi|V_i|\leq \frac{\xi}{\eps'}|A_i|$,
        \begin{align*}
            e_G(A_1,A_2) &\leq e_G(A_1\cap U_1, A_2 \cap U_2)+\frac{2\xi}{\eps'}|A_1||A_2| \\
            &\leq \left(d+\eps+\frac{2\xi}{\eps'}\right)|A_1||A_2| \leq (d+3\eps)|A_1||A_2|.
        \end{align*}
        
        As for the density bound, we have 
        \[d'=\frac{e(V_1,V_2)}{|V_1||V_2|} \geq \frac{e(U_1,U_2)}{(1+\xi)^2|U_1||U_2|} \geq (1-3\xi)d \geq d - 3\xi\]
        and 
        \[d'\leq \frac{e(V_1,V_2)}{|V_1||V_2|}\leq  \frac{e(U_1,U_2)+2\xi|U_1||U_2|+\xi^2|U_1||U_2|}{|U_1||U_2|}\leq d+3\xi.\]
    \end{proof}
	
    To state the version of the Regularity Lemma which we will apply, we recall the notion of regular partitions and the reduced graph.
    We say that the partition $V_0\dcup V_1 \dcup \ldots \dcup V_\ell$ of $V(G)$ is an \emph{$\eps$-regular partition} if $|V_0| \leq \eps |V(G)|$, $|V_1| = \ldots = |V_{\ell}|$, and for all but at most $\eps \ell^2$ pairs $(i,j) \in [\ell] \times [\ell]$, $i\neq j$, the pair $(V_i, V_j)$ is $\eps$-regular. The \emph{$(\eps, d)$-reduced graph} $R$ with respect to the above partition is a graph on vertex set $[\ell]$ such that the edges $ij$ correspond precisely to the $(\eps, d')$-regular pairs $(V_i,V_j)$ with density $d' \geq d$.  
	
    \begin{lemma}[\cite{abcd22}, Lemma~2.6]\label{l:regularity}
        For all $0 < \eps \leq 1$ and $\ell_0 \in \N$, there exists $L, n_0 \in \N$ such that for every $0<d<\alpha<1$, every graph $G$ on $n > n_0$ vertices with minimum degree $\delta(G)\geq \alpha n$ has an $\eps$-regular partition $V_0 \dcup V_1 \dcup \ldots \dcup V_\ell$ with $(\eps,d)$-reduced graph $R$ on $\ell$ vertices such that $\ell_0 \leq \ell \leq L$ and $\delta(R)\geq (\alpha - d - 2\eps)\ell$.
    \end{lemma}

    Next, to find spanning subgraphs, it is convenient to strengthen the notion of regular pairs to super-regular pairs.   
    A pair of sets $(V_1, V_2)$ in a graph $G$ is \emph{$(\eps, d, \vartheta)$-super-regular} if it is $(\eps, d)$-regular, and for $i \neq j$, every vertex $v \in V_i$ has at least $\vartheta|V_j|$ neighbours in $V_j$. As usual, we refer to an $(\eps, d, d-\eps)$-super-regular pair as \emph{$(\eps, d)$-super-regular}. We will say that a $k$-tuple $(V_1,\dots, V_k)$ is $(\eps,d)$-regular (resp.~$(\eps,d,\vartheta)$-super-regular) if each pair in this tuple is $(\eps,d)$-regular (resp.~$(\eps,d,\vartheta)$-super-regular).

    The following lemma allows us to pass from regular $k$-tuples to super-regular $k$-tuples with appropriate parameters by discarding some vertices.

    \begin{lemma}[\cite{abcd22}, Lemma~2.9]\label{l:super-reg-k-tuple}
        Let $k\in\mathbb{N}$ and $0<\eps<d<1$ with $\eps\leq\frac{1}{2k}$. If $(V_1,\dots, V_k)$ is an $(\eps,d^+)$-regular $k$-tuple of vertex-disjoint sets of size $n$, then there are subsets $\tilde V_i \subseteq V_i$, $i\in[k]$ with $|\tilde V_i|=\lceil (1-k\eps)n\rceil$ for all $i\in[k]$ so that the $k$-tuple $(\tilde V_1,\dots,\tilde V_k)$ is $(2\eps, (d-\eps)^+, d-k\eps)$-super-regular.
    \end{lemma}
	
    The next lemma allows us to discard some edges in a super-regular pair and maintain super-regularity with appropriate parameters.
	
    \begin{lemma}[\cite{abcd22}, Lemma~2.12]\label{l:super-reg-min}
        Let $0<\eps<1$ and let $G$ be a bipartite graph with parts $V_1, V_2$ of size $n$, for sufficiently large $n$. Let $\vartheta \in [4 \eps, 1]$. If $(V_1, V_2)$ is $(\eps^2,d', \vartheta-\eps^2)$-super-regular, then there is a spanning subgraph $G' \subset G$ such that $(V_1,V_2)$ is~$(4\eps,\vartheta, \vartheta - 4\eps)$-super-regular in~$G'$.
	\end{lemma}
	
    \subsection{Other auxiliary lemmas}
	
    \begin{lemma}[\cite{akr24}, Lemma~7.3]
	%this is lemma 7.2 from cfrag 4, but with no remainder and V_i replaced by V_{i,1}U... U V_{i,s}
	\label{l:stronger-regular-pairs-factors}
        For integers $m, s, r$  and constants $ \vartheta \gg \eps$, there exists a constant $C$ such that the following holds. Let $G$ be an $nr$-vertex graph and $(V_{i,j})$ be a collection of $n$-vertex pairwise disjoint subsets of $G$, with $i \in [m]$ and $j \leq s^*(i) \leq s$. Suppose that for $i \neq i'$, each $(V_{i,j}, V_{i',j'})$ is an $(\eps, d,  \vartheta)$-super-regular pair in $G$ with some $d \geq \vartheta$. For $r = s^*(1)+\dots+s^*(m)$ and $p = C{n^{-2/s}}\left(\log n \right)^{2/(s(s-1))}$, w.h.p.~the graph $G \cup G_{nr, p}$ contains a $K_r$-factor in which every copy of $K_r$ uses exactly one vertex of each $V_{i,j}$.
    \end{lemma}
	
    We will need the following $Z$-conforming strengthening of the above lemma.
	
    \begin{lemma}\label{l:st-reg-pairs-factors}
        Let $r=ms+t$ with $t \in [s]$. For all $\vartheta \gg \eps $, there exists a constant $C$ such that the following holds. Let $G$ be an $(m+1)$-partite $n$-vertex graph with vertex classes $V_1 \dcup \dots \dcup V_{m+1} = V$ satisfying 
		$$|V_i| = \frac{s|V|}{r} \ \text{ for } i \in [m] \text{ \quad and \quad } |V_{m+1}| = \frac{t|V|}{r}.$$ Suppose that each pair $(V_i, V_j)$, $i \neq j$, is $( \eps, d, \vartheta)$-super-regular with $d \geq \vartheta$. Let $Z\subset V$ be given, with $|Z\cap V_i| \leq \eps |V_i|$ for $i \in [m+1]$. For $p = C{n^{-2/s}}\left(\log n \right)^{2/(s(s-1))}$, if $n=|V|$ is divisible by $r$, then w.h.p.\,the graph $G \cup \Gnp$ contains a $Z$-conforming $K_r$-factor.
    \end{lemma}
    
    \begin{proof}
        We partition the sets $V_i$ as follows. For $i\in[m]$, partition the set $V_i$ into sets $V_{i,1},\dots,V_{i,s}$ by first putting all vertices in $Z\cap V_i$ in $V_{i,1}$, and then splitting the remaining vertices of $V_i\setminus Z$ between sets $V_{i,1},\dots,V_{i,s}$, so that each of them has the same size $\frac{|V|}{r}$, uniformly at random. We partition the set $V_{m+1}$ analogously, but into $t$ sets. By Lemma~\ref{l:slicing}, all pairs $(V_{i,j}, V_{i',j'})$, $i\neq i
        '$, are $(s\eps,(d-\eps)^+)$-regular. Moreover, since we split the vertices randomly, w.h.p.~the degree of any vertex in $V_{i,j}$ into $V_{i',j'}$, $i\neq i'$, is at least $\frac{\vartheta}{2}|V_{i',j'}|$. Indeed, consider a vertex $x\in V_{i,j}$. Since the pair $(V_i, V_{i'})$ is $(\eps,d,\vartheta)$-super-regular, $x$ has at least $\vartheta|V_{i'}|$ neighbours in $V_{i'}$. Let $X\subset V_{i'}$ consist of exactly $\vartheta|V_{i'}|$ arbitrarily chosen neighbours of $x$. Then $|X\cap V_{i',j'}|$ is a hypergeometric random variable with expectation 
        \[\mu \geq \frac{(1-\eps)|V_{i',j'}|\vartheta|V_{i'}|}{|V_{i'}|} = (1-\eps)\vartheta|V_{i'j'}| = \frac{(1-\eps)\vartheta}{r}n. \]
        Thus, using Chernoff-type bounds for the hypergeometric distribution (see, e.g.~\cite{skala2013hypergeometric}), the probability that $|X\cap V_{i',j'}|<\frac{\vartheta}2|V_{i'j'}|$ is $e^{-\Omega(n)}$, so the assertion follows by a union bound over all vertices. 
        In particular, the pairs $(V_{i,j}, V_{i',j'})$, $i\neq i
        '$, are $\left(s\eps,(d-\eps)^+,\frac{\vartheta}{2}\right)$-super-regular.

        We can now use Lemma~\ref{l:stronger-regular-pairs-factors} to get the desired $Z$-conforming $K_r$-factor.
    \end{proof}

%************************* Proof of the main theorem ****************************

    \section{Proof of the main theorem}
        \label{sec:pf-of-main}

    In this section, we present a proof of Theorem~\ref{t:main}. Before we proceed, we state the key auxiliary results, the first two of which is the following treatment of the fully non-extremal case. 
    
    \begin{thm}
        \label{t:non-extremal-conforming}
        For any integers $r$ and $s$ with $2 \leq s < r$, let $r^{-1} \gg \gamma \gg \delta\gg \zeta \gg C^{-1} $. Let $G$ be an $n$-vertex graph with minimum degree at least $\left(1 - \frac{s}{r} - \delta\right)n$. Suppose that every subset of $V(G)$ of order $\frac{sn}{r}$ contains at least $\gamma n^2 $ edges. Let $Z\subset V(G)$ with $|Z|\leq \zeta n$ be given.  For $n$ divisible by $r$ and $p = Cp_s(n)$, w.h.p.\,the graph $G^p = G \cup \Gnp$ has a~$Z$-conforming~$K_r$-factor.
    \end{thm}

    \begin{lemma}
        \label{l:annoying-non-extremal}
        For any integers $r$ and $s$ with $2 \leq s < r$, let $r^{-1} \gg \gamma\gg \delta \gg \eta \gg C^{-1} $. Let $G$ be an $n$-vertex graph with minimum degree at least $\left(1 - \frac{s}{r} - \delta\right)n$. Suppose that every subset of $V(G)$ of order $\frac{sn}{r}$ contains at least $\gamma n^2 $ edges. For $p = Cp_s(n)$, w.h.p.\,the graph $G^p = G \cup \Gnp$ contains a collection $\cC$ of $\eta n$ vertex-disjoint copies of $K_{r+1}$.
    \end{lemma}
    
    The proofs of Theorem~\ref{t:non-extremal-conforming} and Lemma~\ref{l:annoying-non-extremal} can be found in section~\ref{sec:non-extremal}. We will also need the following lemma which extracts the structure of our initial deterministic graph $G$. To keep this section focused, we defer its proof to section \ref{sec:partition-lemma}.

        \begin{lemma}
            \label{l:4.3chr} Let $r = ms+t$ be as before. There is a constant $\beta_0>0$ such that for all $\beta \leq \beta_0$ there is a $\gamma>0$ for which the following holds. If $\gamma_1 \ll \gamma_2 \ll \ldots \ll \gamma_m \ll \gamma$ and $G$ is an $n$-vertex graph with minimum degree at least $\left( 1-\frac sr \right)n$, then for some $h \in \{0\}\cup[m]$ there exists a partition $V(G) = A_1 \dcup A_2 \dcup \dots \dcup A_{h+1}$ satisfying the following conditions.
            \begin{enumerate}
                \item\label{partition(i)} For $i \in [h]$, $|A_i| = \left(\frac{s}{r} \pm \gamma_h\right)n$.
                \item\label{partition(ii)} Every vertex $x \in V(G) \setminus A_{h+1}$ has at most $4 \beta n$ non-neighbours in $A_{h+1}$.
                \item\label{partition(iii)} Every vertex $x \in A_{h+1}$ has at least $\beta n$ neighbours in $A_i$, for any $i \in [h]$.
                \item\label{partition(iv)} If $h<m$ and $X \subset A_{h+1}$ with $|X| =sn/r$, then $e(X) > \gamma_{h+1}^2 n^2$. 
                \item\label{partition(v)} For every distinct $i, j \in [h]$, every vertex $x$ from $A_i$ has at least $\frac 13 |A_j|$ neighbours in $A_j$.
                \item\label{partition(vi)} For every distinct $i, j \in [h+1]$, there are at most $\gamma_h|A_i| |A_j|$ missing edges between $A_i$ and $A_j$.
            \end{enumerate}
        \end{lemma}


    

    Now we can prove the main result, that is Theorem~\ref{t:main}.

   \begin{proof}[Proof of Theorem~\ref{t:main}]
    
    Let $\beta$ be the constant given by Lemma~\ref{l:4.3chr}.
    Apply Lemma~\ref{l:4.3chr} with constants 
    \[0\ll \gamma_1 \ll  \gamma_2  \ll \ldots \ll \gamma_{m+1} \ll \beta \ll 1,\] to get a partition 
    \[V(G) = A_1 \dcup \dots \dcup A_{h+1},\]
    for some $h \in \{0\} \cup [m]$. To be more specific, $\gamma_{m}$ should be small enough for Lemma~\ref{l:st-reg-pairs-factors} to hold for $\left((2r\gamma_m)^{1/3}, \frac12, \frac{\beta}3\right)$-super-regular pairs. 
    Moreover, for any $h\in[m]$, the constants $\gamma_h$ and $\gamma_{h+1}$ should satisfy the following conditions
      \begin{equation}
        %m\gamma_m \ll \beta^2, \quad
        400r^4\gamma_h\ll \beta^2 \quad \text{ and } \quad
          20r^3 \gamma_h \ll \beta\gamma_{h+1}^2.
          %\quad \frac{2m \gamma_h}{\beta} < \zeta^{1/10} ,
          %\quad \gamma_h \leq \eta \quad \text{and} %\quad 
          %\frac{\gamma_h m}{\beta}< \delta.
          \label{eq:gamma-h}
      \end{equation}
    Finally, $\gamma_1$ should also be sufficiently small for Theorem~\ref{t:non-extremal-conforming} to hold and the constant $C$ should be sufficiently large for all the invoked lemmas to hold. 

    In case $h=0$ (the fully non-extremal case), we are done using Theorem~\ref{t:non-extremal-conforming} with $\gamma_{\ref{t:non-extremal-conforming}}=\gamma_1$, $\zeta_{\ref{t:non-extremal-conforming}}=0$, and $Z_{\ref{t:non-extremal-conforming}}=\emptyset$.
    
    Hence, assume $h>0$. Instead of sampling $\Gnp$ with $p = Cp_s(n)$ at once, we will expose four independent graphs $G_i \sim G(n, p/5)$ for $i \in [4]$. The union $G_1 \cup G_2 \cup G_3\cup G_4$ can be viewed as a subgraph of $\Gnp$.
    Furthermore, we assume that
    \begin{itemize}
        \item[($\ast$)] any subset of $G_1$ of size at least $\frac{n}{\sqrt {\log n}}$ contains a copy of $K_s$,
    \end{itemize}
     which occurs w.h.p.~by Lemma~\ref{l:ks-in-lin-sized}.

    In order to control the vertices that have too many non-neighbours in other parts, for $i\in[h+1]$, we define the following sets 
    \[Z_i=\left\{v\in A_i : |A_j\setminus N_G(v)| > \beta |A_j| \text{ for some } j \in [h]\setminus\{i\}\right\}.\] 
    Note that by property \ref{partition(vi)} in Lemma~\ref{l:4.3chr}, we have 
     \begin{equation}
        \label{eq:zi-bound}
         |Z_i|\leq m \cdot \frac {\gamma_h}{\beta}|A_i|. 
     \end{equation}
    Let $Z=\bigcup_{i\in [h+1]}Z_i$. The following standard claim will be used throughout the proof.
        
    \begin{claim}
        \label{c:extend-by-Ks}
        Suppose that $K\subset V(G)$ is a set of vertices containing at most one vertex from $Z$, and such that $|K|\leq r$. Suppose that $K\cap A_j=\emptyset$ for some $j\in [h+1]$. Let $A_j'\subseteq A_j$ be such that $|A_j'|\geq \frac 56 |A_{j}|$. Then there is a $K_s$-copy in $G_1[(N_G(K)\cap A_j')\setminus Z_j]$. 
    \end{claim} 
         
    \begin{proof}
        First, suppose that $j\in [h]$. Since $K$ is disjoint from $A_j$, each vertex in $K\setminus Z$ has at most $\beta|A_j|$ non-neighbours in $|A_j|$. If there is a vertex in $K\cap Z$, then it has at most $\frac 23 |A_j|$ non-neighbours in $|A_j|$. Therefore $|(N_G(K)\cap A_j')\setminus Z_j|\geq |A_j|\left(\frac 56-\frac 23-r\beta-\frac{m\cdot \gamma_m}{\beta}\right)\geq \frac 17 |A_j|$. 
            
        If $j=h+1$, using property \ref{partition(ii)} from Lemma~\ref{l:4.3chr}, we similarly get that $|(N_G(K)\cap A_{h+1}')\setminus Z_{h+1}|\geq \frac 12 |A_{h+1}|$. The claim now follows from ($\ast$).
    \end{proof}

    The rest of the proof is divided into three steps. We first remove a small number of copies of $K_r$ in order to balance the sizes of the partition sets $A_i$, $i\in[h+1]$. Let $A_i' \subset A_i$ be the remaining parts. In the second step, we factorize the part induced by $A_{h+1}'$ with appropriate cliques of size $r-sh$. Finally, in the third step we extend this $K_{r-sh}$-factor to a~$K_r$-factor covering the remaining parts and such that each clique in this factor uses exactly $s$ vertices from any set $A_i'$, $i\in[h]$.

    \smallskip

    {\bf Step 1.~Balancing the sizes of the partition sets}

    \smallskip
        
    Let $g_1, g_2, \ldots, g_{h+1}$ be integers such that 
    \[|A_i|=\frac{sn}{r}+g_i, \text{ for } i\in[h] \quad \text{ and } \quad |A_{h+1}|=\frac{(m-h)s+t}{r}\cdot n+g_{h+1}.\] 
    Notice that $\sum_{i \in [h+1]}g_i = 0$ and, by property \ref{partition(i)} in Lemma~\ref{l:4.3chr}, 
    \begin{equation}\label{eq:g_i-sizes}
        |g_i| \leq \gamma_h n, \text{ for } i\in[h] \quad \text{ and } \quad |g_{h+1}|\leq h\gamma_hn.
    \end{equation}
    Let $\sigma_i=g_i/|g_i|$ be the sign of $g_i$, and set \[g=\sum_{i=1}^{h+1}|g_i|/2 \quad \text{ and } \quad I_1 = \{i\in[m+1]: g_i>0\}.\]
    Let $I_2 = [m+1]\setminus I_1$.

    \smallskip
    Consider first the case $h=m$, which is slightly easier to handle. 
    
    \begin{claim}
        \label{c:size-adj}
        If $h=m$, then w.h.p.~there is a collection $\cC$ of $g$ vertex-disjoint $Z$-conforming copies of $K_r$ in $G\cup G_1$ such that 
        \begin{itemize}
            \item for each $i\in [m]$, there are exactly $|g_i|$ cliques $K\in \cC$ with $|K\cap A_i|=s+\sigma_i$, while all other cliques $\tilde K\in \cC$ satisfy $|\tilde K\cap A_i|= s$,
            \item there are exactly $|g_{m+1}|$ cliques $K\in \cC$ with $|K\cap A_{m+1}|=t+\sigma_{m+1}$, while all other cliques $\tilde K\in\cC$ satisfy $|\tilde K\cap A_{m+1}|=t$.
        \end{itemize}
    \end{claim}
    
    \begin{proof}[Proof of Claim~\ref{c:size-adj}] 
        For any $i\in I_1\setminus\{m+1\}$ we have $\delta (G[A_i])\geq g_i$. Moreover, for any $v\in Z_i$, since there is an index $j\in [m]\setminus\{i\}$ such that $|A_j\cap N_G(v)|< (1-\beta) |A_j|$, then by the minimum-degree assumption on $G$, we have $|A_i\cap N_G(v)|\geq \beta |A_i|$. Hence, by Lemma~\ref{l:5.5chr-stronger} with $\delta_{\ref{l:5.5chr-stronger}}=m\gamma_m/\beta$, $\beta_{\ref{l:5.5chr-stronger}}=\beta$, and $\gamma_{\ref{l:5.5chr-stronger}}=\gamma_m$, w.h.p.~there is a collection $\cC_i$ of $g_i$ vertex-disjoint $Z_i$-conforming $K_{s+1}$-copies in $(G \cup G_1)[A_i]$. %and recall that \textit{$Z_i$-conforming} means that each of those copies contains at most one vertex of $Z_i$. 
        
        If $m+1\in I_1$, we have two cases.
        
        \textbf{Case 1.} If $t<s$, let $\cC_{m+1}$ be a collection of $g_{m+1}$ vertex-disjoint $Z_i$-conforming $K_{t+1}$-copies in $(G\cup G_1)[A_{m+1}]$, which can be found greedily using ($\ast$), since $t+1 \leq s$. Here we can actually make the copies in $\cC_{m+1}$ to be vertex-disjoint from $Z_i$.

        \textbf{Case 2.} If $t=s$, by \eqref{eq:g_i-sizes} we have $|A_{m+1}| = \frac{sn}{r}+g_{m+1} \leq \frac {sn}r + m\gamma_mn$ and $\delta(G[A_{m+1}])\geq g_{m+1}$, so we can proceed to find a collection $\cC_{m+1}$ of $g_{m+1}$ vertex-disjoint $Z_{m+1}$-conforming $K_{s+1}$-copies using Lemma~\ref{l:5.5chr-stronger}, as in the case $i\in [m]$, but with $\gamma_{\ref{l:5.5chr-stronger}}=m\gamma_m$.

        Using Claim~\ref{c:extend-by-Ks}, we can now extend the cliques in $\cC = \bigcup_{i\in I_1}\cC_i$ to vertex-disjoint $K_{r+1}$-copies. To this end, we perform the following process. For each $K\in \cC_i$, go through the indices $j\in [m]\setminus \{i\}$ in order, and for each $j$, find a $K_s$-copy in $(A_j \cap N_G(K))\setminus Z_j$ (also avoiding the previously used vertices) and add it to $K$. Similarly, if $i\neq m+1$, find a $K_t$-copy in $(A_{m+1}\cap N_G(K))\setminus Z_{m+1}$ and use it to extend $K$.  This will turn $\cC$ into a collection of vertex-disjoint $Z$-conforming copies of $K_{r+1}$. Since $r|\cC|< \frac 16|A_j|$ for each $j\in [m+1]$,  Claim~\ref{c:extend-by-Ks} is applicable in each step. %We point out that we have also just proved the following claim, which will be used in case $h<m$ in a slightly different form.

        %\begin{claim}\label{c:r+1-cliques-extremal}
        %    There is a collection $\cC=\bigcup_{i\in I_1}\cC_i$ of vertex-disjoint $Z$-conforming $K_{r+1}$-copies in $G\cup G_1$ such that $|\cC_i|=g_i$ for every $i\in I_1$. Moreover, for $i \in I_1 \setminus \{m+1\}$ and $K \in \cC_i$,
         %   $$|V(K)\cap A_j|=
          %  \begin{cases}
          %      s+1, & j=i,\\
          %      t, & j=m+1,\\
          %      s, & j \in [m]\setminus\{i\}.
          %  \end{cases}.
          %  $$
          %  Moreover, if $g_{m+1}>0$, then every $K \in \cC_{m+1}$  intersects $A_{m+1}$ in $t+1$ vertices, and all other sets $A_j$ in $s$ vertices.
        %\end{claim}
        
         Finally, since $\sum_{j\in I_2}|g_j|=g=|\cC|$, to every $K\in \cC$ we can assign an index $j_K\in I_2$ so that each index $j\in I_2$ is assigned to exactly $g_j$ cliques in $\cC$. Note that if $K\in\cC_i$, then $j_K\neq i$. Now, for each $K\in \cC$, remove from $K$ a single vertex belonging to $A_{j_K}$, to reduce a~$K_{r+1}$-copy to a $K_r$-copy. Then the resulting collection $\cC$ satisfies the conditions stated in Claim~\ref{c:size-adj}.
    \end{proof}

    Suppose now that we are in the intermediate case, i.e.~$0<h<m$. Let 
    \begin{equation}\label{eq:r'}
        r'=r-hs = (m-h)s+t \quad \text{ and } \quad \tilde{n} = \frac{r'}{r}n,
    \end{equation}
    so that $|A_{h+1}| = \tilde{n}+g_{h+1}$.
    The following auxiliary claim ensures that the graph $G[A_{h+1}]$ is fully non-extremal for appropriate parameters even after removing $O(\gamma_{h+1}^2n)$ vertices. %Hence we can apply to it Theorem~\ref{t:non-extremal-conforming} and Lemma~\ref{l:annoying-non-extremal}.

    \begin{claim} \label{c:min-degree-in-A_h+1} Let $A \subset A_{h+1}$ with $||A|-\tilde{n}| \leq \frac 19 \gamma_{h+1}^2 n$. Then 
        \begin{equation*} %\label{eq:min-degree-in-A_h=1}
            \delta(G[A])\geq\left(1-\frac{s}{r'}\right)|A| - ||A|-\tilde{n}|.
        \end{equation*}
        Secondly, if $S \subset A$ with $|S| \geq \frac{s}{r'}|A|$, then $e_G(S) \geq \frac 12 \gamma_{h+1}^2 n^2$.
    \end{claim}
    
    \begin{proof}[Proof of Claim~\ref{c:min-degree-in-A_h+1}]
        For any $x\in A$, we have
        \begin{align*}
            \deg_{G[A]}(x) & \geq \left(1-\frac{s}{r}\right)n-(n-|A|) = |A| - \frac{s}{r}n = |A| - \frac{s}{r'} \tilde{n} \\
            & \geq \left(1 - \frac{s}{r'}\right)|A| - ||A|-\tilde{n}|,
        \end{align*}
        so the minimum degree condition for $G[A]$ holds as required.

        The second part of the claim follows directly from property \ref{partition(iv)} of Lemma~\ref{l:4.3chr} if $|A| \geq \tilde{n}$, since in this case $|S|\geq\frac{s}{r'}|A| \geq \frac{s}{r}n$. 
        So assume that $|A|<\tilde{n}$.
        Let $S \subset A$ with $|S| \geq \frac{s}{r'}|A|$. Let $X$ be any subset of $A$ containing $S$ and such that $|X|= \frac{s n}{r} = \frac{s\tilde{n}}{r'}$. Note that
        $$|X \setminus S| \leq \frac{s}{r'}\left(\tilde{n}-|A|\right)\leq \frac 19 \gamma_{h+1}^2 n.$$
        Again, by property \ref{partition(iv)} of Lemma~\ref{l:4.3chr}, $X$ induces at least $\gamma_{h+1}^2 n^2$ edges in $G$. At most $n|X \setminus S| $ of those edges have an endpoint outside of $S$, so we conclude that
        $$e_G(S) \geq \frac 12 \gamma_{h+1}^2 n^2.$$
    \end{proof}
    
    Our aim now is to find a collection of $K_r$-copies analogous to that in Claim~\ref{c:size-adj}. To do this, we will expose the graphs $G_1$ and $G_2$.

    \begin{claim}
        \label{c:size-adj_h}
        If $0<h<m$, then w.h.p.~there is a collection $\cC$ of $g$ vertex-disjoint $Z$-conforming copies of $K_r$ in $G\cup G_1\cup G_2$ such that 
        \begin{itemize}
            \item for each $i\in [h]$, there are exactly $|g_i|$ cliques $K\in \cC$ with $|K\cap A_i|=s+\sigma_i$, while all other cliques $\tilde K\in \cC$ satisfy $|\tilde K\cap A_i|= s$,
            \item there are exactly $|g_{h+1}|$ cliques $K\in \cC$ with $|K\cap A_{h+1}|=r'+\sigma_{m+1}$, while all other cliques $\tilde K\in\cC$ satisfy $|\tilde K\cap A_{h+1}|=r'$.
        \end{itemize}
    \end{claim}

    \begin{proof}[Proof of Claim~\ref{c:size-adj_h}]
    Let $A = A_{h+1}\setminus Z_{h+1}$ and consider first the graph $G[A]$. Recall that by \eqref{eq:zi-bound} we have
    \[|Z_{h+1}|\leq m\cdot\frac{\gamma_h}{\beta}|A_{h+1}|.\] 
    Therefore, by Claim~\ref{c:min-degree-in-A_h+1}, we have
    $$\delta(G[A]) \geq \left(1-\frac{s}{r'} -\frac{2m \gamma_h}{\beta}\right)|A|,$$
    and any set $S \subset A$ with $|S|\geq \frac{s|A|}{r'}$ induces in $G[A]$ at least $\frac 12 \gamma_{h+1}^2 n^2$ edges. In other words, one can think of $G[A]$ as a fully non-extremal graph.
    
    We now expose the random graph $G_2$. By the properties above and using \eqref{eq:gamma-h}, we may apply to $G[A]$ Theorem~\ref{t:non-extremal-conforming} with $r_{\ref{t:non-extremal-conforming}}=r'$, $\delta_{\ref{t:non-extremal-conforming}}=2m\gamma_h/\beta$, $\gamma_{\ref{t:non-extremal-conforming}}=\frac12\gamma_{h+1}^2$, and $Z_{\ref{t:non-extremal-conforming}}=\emptyset$, after perhaps removing at most $r'-1$ vertices for divisibility issues. W.h.p., we then get in $(G\cup G_2)[A]$:
    \begin{itemize}
        \item a collection $\mathcal{F}$ of vertex-disjoint copies of $K_{r'}$ covering all but at most $r'-1$ vertices of $A$.
    \end{itemize}
    Similarly, in case $g_{h+1}>0$, applying to $G[A]$ Lemma~\ref{l:annoying-non-extremal} with $r_{\ref{l:annoying-non-extremal}}=r'$, $\gamma_{\ref{l:annoying-non-extremal}}=\frac12\gamma_{h+1}^2$, $\delta_{\ref{l:annoying-non-extremal}}=2m\gamma_h/\beta$ and $\eta_{\ref{l:annoying-non-extremal}}=h\gamma_h$, so that by \eqref{eq:g_i-sizes} we have $g_{h+1}\leq \eta_{\ref{l:annoying-non-extremal}}n$, w.h.p.~we get in $(G\cup G_2)[A]$:
    \begin{itemize}
        \item a collection $\cC_{h+1}$ of $g_{h+1}$ vertex-disjoint copies of $K_{r'+1}$.
    \end{itemize}
    
    The first collection will be used to extend some cliques $K$ spanned by vertices in $A_1\dcup \dots \dcup A_h$ to copies of $K_{r+1}$. The latter one will only be used in case $g_{h+1}>0$ to deal with surplus vertices in $A_{h+1}$. 

    Now we expose the graph $G_1$. Suppose first that $g_{h+1} > 0$. Using the same strategy as in the proof of Claim~\ref{c:size-adj}, we can extend cliques in $\cC_{h+1}$ to copies of $K_{r+1}$ in $G\cup G_1\cup G_2$, each having exactly $s$ vertices in $A_i$, for $i\in[h]$. Let $\cC_{h+1}^+$ denote this collection of $K_{r+1}$-copies.
    
    Next, using again the same strategy as in the proof of Claim~\ref{c:size-adj}, we can find a collection $\cC = \bigcup_{i\in I_1 \setminus \{h+1\}}\cC_i$ of vertex-disjoint $Z$-conforming $K_{hs + 1}$-copies in $(G\cup G_1)[A_1\dcup \dots \dcup A_h]$, such that any clique $K\in \cC_i$ contains $s+1$ vertices from $A_i$ and $s$ vertices from the other sets $A_j$, $j\in[h]\setminus\{i\}$.
    We will now use the partial $K_{r'}$-factor $\cF$ to turn those $K_{hs+1}$-copies into copies of $K_{r+1}$, one by one. 
    In particular, for each clique $K\in \cC$, we can find a clique $K'\in \cF$ such that the vertices in $V(K)\cup V(K')$ span a copy of $K_{r+1}$ in $G \cup G_1 \cup G_2$. Indeed, by removing $K'$ in each step, by \eqref{eq:g_i-sizes} we will have removed at most $|\cC|=g\leq h\gamma_h n$ cliques from $\cF$. On the other hand, by property \ref{partition(ii)} of Lemma~\ref{l:4.3chr}, the number of non-neighbours of $K$ in $A_{h+1}$ satisfies $|A_{h+1}\setminus N_G(K)|\leq 4\beta (hs+1)n$, and this is an upper bound for the number of cliques in $\cF$ that cannot be used as an extension for $K$. Therefore, after forbidding those cliques, as well as those cliques we have already used in the process, we still have at least 
    \[ |\cF| -(4\beta (hs+1)+h\gamma_h)n \geq \lfloor n'/r'\rfloor-(\gamma_{h+1}+4\beta (hs+1)+h\gamma_h)n\geq 1\] 
    other cliques in $\cF$, so we can find an appropriate extension $K'\in \cF$ for $K$. Let $\cC'$ be the new collection of those extended copies and set $\cC^+ =\cC'\cup \cC_{h+1}^+$.

    Finally, using again the same strategy as in the proof of Claim~\ref{c:size-adj}, we can turn the collection $\cC^+$ into the desired collection of copies of $K_r$ by carefully removing one vertex from each clique in $\cC^+$.
    \end{proof}

    Let $\cC$ be the collection of copies of $K_r$ given by Claim~\ref{c:size-adj} in case $h=m$ or by Claim~\ref{c:size-adj_h} in case $0<h<m$. We remove from $A_1\dcup\dots\dcup A_{h+1}$ the vertices belonging to cliques in $\cC$. Let $A_i' \subset A_i$, $i\in[h+1]$, be the vertex sets obtained after this removal. Let $n' = n-gr$, and note that due to the properties of the collection $\cC$ we have
    \begin{equation}\label{eq:A_{h+1}'}
        |A_i'| = \frac{s n'}{r} \text{ for } i \in [h] \quad \text{ and } \quad |A_{h+1}'| = \frac{(r-hs)n'}{r}=\frac{r'n'}{r}.
    \end{equation}
    Note also that the sets $A_1', \ldots, A_{h+1}'$ still satisfy properties \ref{partition(ii)}--\ref{partition(vi)} from Lemma~\ref{l:4.3chr}, with slightly modified constants, as detailed below. Set \[G'=G[A_1'\dcup\dots\dcup A_{h+1}'],\]  and proceed to the next step.

    \smallskip

    {\bf Step 2.~Covering $A_{h+1}'$ with a $K_{r-hs}$-factor $\cP$}

    \smallskip
    
    We now expose the random graph $G_3$ in order to find an appropriate $K_{r-hs}$-factor in $(G'\cup G_3)[A_{h+1}']$. Moreover, since we later want to extend this factor to a $K_r$-factor, we need to carefully treat the vertices with many non-neighbours outside $A_{h+1}'$. Thus, let 
    \[Z' = \{v \in A_{h+1}': |A_j' \setminus N_G(v)| \geq \frac{\beta}{10r}|A_j'| \text{ for some } j\in[h]\}.\] 
    Using property \ref{partition(vi)} from Lemma~\ref{l:4.3chr}, we have 
    \[ |Z'| \leq 20hr \beta^{-1} \gamma_h |A_{h+1}'| .\] 

    Consider first the case $h=m$. Note that $r-hs=t \leq s$. Moreover, for any vertex $v\in Z'$, there exists $j\in[h]$, such that 
    \[ \text{deg}_{G'[A_{h+1}']}(v) \geq \frac{\beta}{20r}|A_j'| \geq \frac{\beta}{20rs}|A_{h+1}'|.\] 
    Hence, using Lemma~\ref{l:Z-conf.-Kt-factor} applied to the graph $G'[A_{h+1}']$ with $\beta_{\ref{l:Z-conf.-Kt-factor}}=\frac{\beta}{20rs}$ and $\delta_{\ref{l:Z-conf.-Kt-factor}}=20mr\beta^{-1}\gamma_h$, by~\eqref{eq:gamma-h} w.h.p.~we can find in $(G'\cup G_3)[A_{h+1}']$ a $Z'$-conforming $K_t$-factor $\cP$.

    \smallskip
    
    Now consider the case $0<h<m$ and recall by \eqref{eq:r'} that $r-hs = r'$ and $|A_{h+1}|=\tilde{n}+g_{h+1}$ with $\tilde{n}=\frac{r'}{r}n$. Then by \eqref{eq:A_{h+1}'} and \eqref{eq:g_i-sizes}, we have
    \begin{align*} 
        ||A_{h+1}'|-\tilde{n}| = |A_{h+1}'|\left|1- \frac{n}{n'}\right| \leq \frac{2rg}{n}|A_{h+1}'| \leq 2rh\gamma_h |A_{h+1}'|.
    \end{align*}
     Using Claim~\ref{c:min-degree-in-A_h+1} applied with $A=A_{h+1}'$, we get
    \[\delta(G'[A_{h+1}'] \geq \left(1-\frac{s}{r'}-2rh\gamma_h\right)|A_{h+1}'|\] 
    and any set $X\subset A_{h+1}'$ with $|X|=\frac{s|A_{h+1}'|}{r'}$ satisfies 
    \[e_{G'}(X) \geq \frac12\gamma_{h+1}^2n^2 > \gamma_{h+1}^2|A_{h+1}'|.\] 
    Hence, using Theorem~\ref{t:non-extremal-conforming} with $r_{\ref{t:non-extremal-conforming}}=r'$, $Z_{\ref{t:non-extremal-conforming}}=Z'$, $\zeta_{\ref{t:non-extremal-conforming}}=20hr\beta^{-1}\gamma_h$ and $\gamma_{\ref{t:non-extremal-conforming}} = \gamma_{h+1}^2$, we conclude that w.h.p.~the graph $(G \cup G_3)[A_{h+1}']$ contains a $Z'$-conforming $K_{r'}$-factor $\mathcal{P}$.
   %note that the in this application we have delta (the defect in min degree) < (r+1)gamma_h  << zeta. but this is actually weaker than the non-conforming statement, where delta >> zeta.

   \smallskip

    {\bf Step 3.~Extending $\cP$ to a $K_{r}$-factor}

    \smallskip

    Now form the graph $\tilde G = G'/\mathcal{P}$ by \textit{collapsing} each clique $B \in \mathcal{P}$ into a single vertex; that is, the vertex set $A_{h+1}'$ is replaced by the set $\tilde A_{h+1}$ which is in one-to-one correspondence with the collection $\mathcal{P}$, and a vertex $v_B \in \tilde A_{h+1}$ representing the clique $B\in\cP$ is adjacent in $\tilde{G}$ to $v \in A_1'\dcup\dots\dcup A_h'$ if and only if $V(B) \subset N_{G'}(v)$. We will apply Lemma~\ref{l:st-reg-pairs-factors} to the graph $\tilde G$ after verifying the super-regularity conditions. Before we proceed, set
    \[ \tilde r = sh+1 \quad \text{ and } \quad \tilde{n}=|V(\tilde G)|,\]
    and note that
    \[ \tilde n = \frac{sh+1}{r}n' = \frac{\tilde r}{r}n'. \]
    In particular,
    \[ |A_{i}'| = \frac{sn'}{r}=\frac{s\tilde n}{\tilde r} \text{ for } i\in[h] \quad \text{ and } \quad |\tilde A_{h+1}|=\frac{n'}{r}=\frac{\tilde{n}}{\tilde r}.\]

    To verify the super-regularity conditions, notice first that by properties \ref{partition(v)} and \ref{partition(vi)} of Lemma~\ref{l:4.3chr}, each pair $(A_i,A_j)$, $i,j\in[h]$, is $(\gamma_h^{1/3} , d, 1/3)$-super-regular for some $d\geq 1-\gamma_h > 1/2$. Passing to subsets $A_i'\subset A_i$, and using Lemma~\ref{l:slicing}, we get that the pairs $(A_i', A_j')$, $i,j\in[h]$, are $(2\gamma_h^{1/3} , d', 1/4)$-super-regular for some $d'\geq 1/2$. 
    
    Consider now the pair $(A_i', \tilde A_{h+1})$, $i\in[h]$. We first check the degree condition for super-regularity. Take any vertex $v_B\in \tilde A_{h+1}$. The clique $B$ contains at most one vertex from $Z'$, and by property \ref{partition(iii)} of Lemma~\ref{l:4.3chr} such a vertex has at least $\frac12\beta n$ neighbours in $A_i'$. On the other hand, by the definition of $Z'$, the remaining vertices in $B$ have each at most $\frac{\beta}{10r}|A_i'|$ non-neighbours in $A_i'$. Thus, $v_B$ has at least \[\frac12\beta n - \frac{\beta(r-sh)}{10r}|A_i'| \geq \frac{\beta}{3}|A_i'|\]
    neighbours in $A_i'$. On the other hand, property \ref{partition(ii)} of Lemma~\ref{l:4.3chr} yields that any vertex $x\in A_i'$ has at most $4\beta n$ non-neighbours in $A_{h+1}'$ and the same is true for $\tilde A_{h+1}$.

    As for the regularity, by property \ref{partition(vi)} of Lemma~\ref{l:4.3chr} there are at most \[2\gamma_h|A_i'||A_{h+1}'|=2(r-sh)\gamma_h|A_i'||\tilde A_{h+1}|\] missing edges between $A_i'$ and $A_{h+1}'$, and this is also an upper bound for the number of missing edges between $A_i'$ and $\tilde A_{h+1}$. We conclude that the pair $(A_i', \tilde A_{h+1})$ is $((2r\gamma_h)^{1/3}, d'', \beta/3)$-super-regular for some $d'' \geq 1-2r\gamma_h \geq 1/2$.

    Finally, we expose the random graph $G_4$. By Lemma~\ref{l:st-reg-pairs-factors}, the graph $\tilde G\cup G_4[A_1'\dcup\dots\dcup A_h'\dcup \tilde A_{h+1}]$ has a $K_{hs+1}$-factor in which each clique uses exactly $s$ vertices from any set $A_i'$, $i\in[h]$, and one vertex from $\tilde A_{h+1}$.  This factor corresponds in turn to a $K_r$-factor in $G \cup G_3 \cup G_4\left[\bigcup_{i \in [h+1]} A_{i}'\right]$. Together with the $K_r$-cliques in the collection $\cC$ constructed in Step 1.~of the proof, we obtain the desired $K_r$-factor in $G\cup G_1\cup \dots \cup G_4 \subset G\cup \Gnp$.

    \end{proof}
    
    
    \section{The main extremal ingredient}

     In this section, we introduce a family of auxiliary graphs which, together with a suitable weighted graph packing, will play crucial role in splitting the parts of the regular partition in the non-extremal case into smaller parts of appropriate sizes and preserving important properties of the regular pairs. To get the desired sub-partition, we will also need to use some carefully constructed absorbers. 
	
    \subsection{The auxiliary graph $Q_h^{(m)}$}

    \begin{definition}
        For every $h\in [m]$ define the graph $Q_h^{(m)}$ that consists of disjoint 
        \begin{itemize}
            \item $h$-cliques $L_1,\dots,L_{s-t},N_1,\dots N_{(m-h)s+t}$ and
            \item $(m-h+1)$-cliques $M_1,\dots,M_{(m-h)s+t}$
        \end{itemize}
        together with all edges from $L_i$ to $M_j$ and all edges from $M_j$ to $N_j$, for any $i\in [s-t]$, $j\in[(m-h)s+t]$. 
    \end{definition}

    \begin{remark}\label{rem:Q_h}
        Note that $|V(Q_h)|=(m+1-h)r$. Moreover, in case $t=0$, $Q_m^{(m)}$ is just a~union of $s$ vertex-disjoint copies of $K_m$.
    \end{remark}

    When $m$ is clear from the context, we will use the notation $Q_h\coloneq Q_h^{(m)}$. We will also refer to the cliques $L_1,\dots, L_{s-t}$ as the $L$-sets, or to the $L$-part of $Q_h$. Similarly, we will refer to the $M$- and $N$-sets, or $M$- and $N$-part.

   \begin{figure}
    \begin{tikzpicture}[scale=0.8]
        \centering

        \coordinate (v1) at (0,2);
        \coordinate (v2) at (0,-2);
        \coordinate (v3) at (4,3);
        \coordinate (v4) at (4,1);
        \coordinate (v5) at (4,-3);
        \coordinate (v6) at (8,3);
        \coordinate (v7) at (8,1);
        \coordinate (v8) at (8,-3);
        
        \draw[line width=3] (v3) -- (v6);
        \draw[line width=3] (v4) -- (v7);
        \draw[line width=3] (v5) -- (v8);
        \draw[line width=3] (v1) -- (v3);
        \draw[line width=3] (v1) -- (v4);
        \draw[line width=3] (v1) -- (v5);
        \draw[line width=3] (v2) -- (v3);
        \draw[line width=3] (v2) -- (v4);
        \draw[line width=3] (v2) -- (v5);
        
        \draw[fill=white] (0,2) circle [radius=0.5];
        \draw[fill=white] (0,-2) circle [radius=0.5];

        \draw[fill=white] (4,3) circle [radius=0.7];
        \draw[fill=white] (4,1) circle [radius=0.7];
        \draw[fill=white] (4,-3) circle [radius=0.7];
        
        \draw[fill=white] (8,3) circle [radius=0.5];
        \draw[fill=white] (8,1) circle [radius=0.5];
        \draw[fill=white] (8,-3) circle [radius=0.5];
        
        \node at (0,0) {$\vdots$};
        \node at (4,-1) {$\vdots$};
        \node at (8,-1) {$\vdots$};

        \node at (0,4) [above] {$K_h$-copies};
        \node at (4,4) [above] {$K_{m-h+1}$-copies};
        \node at (8,4) [above] {$K_h$-copies};
           
        \node at (0,-4) [below] {$L$-sets};
        \node at (4,-4) [below] {$M$-sets};
        \node at (8,-4) [below] {$N$-sets};

        \draw [decorate,decoration={brace,amplitude=7pt,mirror,raise=4ex}]
  (-0.5,2.5) -- (-0.5,-2.5) node[midway,xshift=-4em]{$s-t$};
        \draw [decorate,decoration={brace,amplitude=7pt,mirror,raise=4ex}]
  (8.5,-3.5) -- (8.5,3.5) node[midway,xshift=6em]{$(m-h)s+t$};

    \end{tikzpicture}
\caption{The graph $Q_h^{(m)}$. Thick edges represent complete bipartite graphs.}
\end{figure}

    \begin{lemma}\label{l:6.2chr}
    Let $r = ms+t$ with $0 \leq t <s$. There exists a constant $C>0$ such that for all $\gamma>0$ the following holds. There exists a constant $\delta >0$ and $n_0\in\mathbb{N}$ such that every graph $G$ on $n$ vertices, $n\geq n_0$, with $\delta(G)\geq \left(1-\frac{s}{r}-\delta\right)n$, satisfies one of the following two conditions.
    %
    \begin{enumerate}[label=(\roman*)]
        \item There are vertex-disjoint copies of $K_{m+1},Q_1^{(m)},\dots, Q_m^{(m)}$ in $G$ covering all but at most $C$ vertices, or
        \label{l:6.2chr(i)}
        \item $G$ has an independent set of size $\left(\frac{s}{r}-\gamma\right)n$. \label{l:6.2chr(ii)}
        \end{enumerate}
    \end{lemma}

    Note that by a simple reparametrization of $m$, we immediately get the following corollary which will be useful in the singular case $t=s$.
    
    \begin{cor}\label{cor:6.2chr}
    Let $r = (m+1)s$. There exists a constant $C>0$ such that for all $\gamma>0$ the following holds. There exists a constant $\delta >0$ and $n_0\in\mathbb{N}$ such that every graph $G$ on $n$ vertices, $n\geq n_0$, with $\delta(G)\geq \left(1-\frac{1}{m+1}-\delta\right)n$, satisfies one of the following two conditions.
    
    \begin{enumerate}[label=(\roman*)]
        \item There are vertex-disjoint copies of $K_{m+2},Q_1^{(m+1)},\dots, Q_{m+1}^{(m+1)}$ in $G$ covering all but at most $C$ vertices, or
        \item $G$ has an independent set of size $\left(\frac{1}{m+1}-\gamma\right)n$.
        \end{enumerate}
    \end{cor}
    
    \begin{proof}[Proof of Lemma~\ref{l:6.2chr}]
         Let $C>r^4$ be sufficiently large for the proofs of Claim~\ref{c:p-A-edges-in-m+1} and Claim~\ref{c:p-B-edges} (more specifically, the contained applications of the K\H{o}vari--S\'os--Tur\'an Theorem) to hold, and let $\gamma >0$. For the remaining constants ($\beta, \xi, \eta, \delta$), we have explicit restrictions 
          \[ \beta<r^{-(r+1)}, \quad \xi < \min \{ r^{-3}, \gamma \},\quad 2\eta \leq \xi^2 \beta, \quad \delta \leq \min\left\{\frac{\xi-2\xi^2}{m+1}, \eta \right\}.\]
        Finally, let $n_0$ be sufficiently large.
        Let $G$ be as in the statement of the Lemma and set \[\alpha=1-\frac{s}{r}=\frac{(m-1)s+t}{ms+t},\] so that $\delta(G)\geq (\alpha-\delta)n$. Let \[\Pc=\{K_1,\dots,K_{m+1},Q_1,\dots, Q_m\}.\]
        Note that evidently $G$ has at least one $\Pc$-factor (for instance, each vertex can be covered by its own copy of $K_1$). Given a $\Pc$-factor $\cF$ of $G$, we let $k_i^\cF$ be the number of copies of $K_i$ in $\cF$ (for $i\in [m+1]$) and $q_h^\cF$ be the number of copies of $Q_h$ in $\cF$ (for $h\in [m]$). Now, set
        $$\phi_h^\cF=k_h^\cF+(s-t)q_h^\cF \quad \text
            {for all }\ h\in [m],$$
        $$\phi_{m+1}^\cF=k_{m+1}^\cF+\sum_{h=1}^m((m-h)s+t)q_h^\cF.$$
        We call the vector
    \[(\phi_{m+1}^\cF,\phi_m^\cF,q_m^\cF,\dots, \phi_1^\cF,q_1^\cF)\]
        the \emph{index} of the factor $\cF$. Fix a $\cP$-factor $\cF$ whose index is lexicographically maximal. 
        For $i \in [m+1]$, let $\fA_i$ be the set of $K_i$-copies in $\cF$, and let $A_i$ be the set of vertices defined by them. Define $\fB_h$ and $B_h$ similarly with respect to $Q_h$, for $h\in [m]$. 
        
        Now, if $\sum_{i=1}^m|A_i| \leq C$, then the conclusion \ref{l:6.2chr(i)} holds. So suppose that this is not the case and let $j\in [m]$ be minimal such that 
        \begin{equation}\label{eq:A_j}
            |A_j|\geq \frac{C}{m}.
        \end{equation}
        Our task is to show that $G$ has an independent set of size roughly $\frac{s}{r}n$. We will show that such a set exists by first proving a number of auxiliary claims concerning the properties of factor $\cF$. 
        

        \begin{claim}
            \label{c:singular-A_m}
            If $t=0$, then $|A_m|< sm$ and hence $j<m$.
        \end{claim}
        
        \begin{proof}
            Assume the opposite, so $\cF$ contains at least $s$ copies of $K_m$. By Remark~\ref{rem:Q_h}, we can replace those $s$ copies of $K_m$ by a single copy of $Q_m$. This increases $q_m^\cF$ without changing $\varphi_{m+1}^{\cF}$ and $\varphi_m^{\cF}$.  Hence, we get a new $\Pc$-factor $\cF'$ with an index larger than that of $\cF$. This contradicts the maximality of $\cF$.  
        \end{proof}

        For vertex sets $X,Y\subset V(G)$, let $e(X,Y)$ denote the number of ordered pairs $(u,v)$ with $u\in X$ and $v\in Y$. In particular, if $X\cap Y=\emptyset$, then $e(X,Y)$ is just the number of edges in $G$ with one endpoint in $X$ and the other in $Y$. On the other hand, $e(X, X) = 2e(X)$, which will be used in the following claim and throughout the proof.
        
        \begin{claim}
            \label{c:p-A-edges-in-[p,m]}
            For any $h\in [j,m]$, if $h<m$ or $t>0$, then 
            $$e(A_j,A_h)\leq (\alpha-\beta)|A_j||A_h|.$$
        \end{claim}
       
        \begin{proof}
            Assume that the edge density between $A_j$ and $A_h$ is greater than $\alpha-\beta$.
            Then there exist $\tilde K_j\in \fA_j$ and $\tilde K_h\in \fA_h$ such that $$e(V(\tilde K_j),V(\tilde K_h))>(\alpha-\beta)jh.$$ 
            Note that
            $\alpha jh - j(h-1) = j\left(h-\frac{sh}{r}-h+1 \right)=j\left(\frac{(m-h)s+t}{ms+t}\right)>0$, since we assumed that $h<m$ or $t>0$. Hence,  $(\alpha-\beta) jh > j(h-1)$ for sufficiently small $\beta$. It follows that $$e(V(\tilde K_j),V(\tilde K_h))\geq j(h-1)+1.$$ By the pigeonhole principle, there is a vertex $x\in \tilde K_j$ such that $x$ is a neighbour of all vertices in $\tilde K_h$. Replace $\tilde K_j$ and $\tilde K_h$ with two cliques $ K_{j-1}'$ and $ K_{h+1}'$, by shifting the vertex $x$ from $\tilde K_j$ to $\tilde K_h$. This increases $\phi_{h+1}^\cF$ while the preceding terms stay intact, so again we get a contradiction to the maximality of $\cF$. 
        \end{proof}

        \begin{claim}
            \label{c:p-A-edges-in-m+1}
            $$e(A_j,A_{m+1})\leq |A_j|\left((\alpha-\beta)|A_{m+1}|+C\right).$$
        \end{claim}

        \begin{proof}
            Assume that this is not the case. Then 
            \begin{equation}\label{eq:A_{m+1}}
                |A_{m+1}|\geq C,
            \end{equation}
            as otherwise $e(A_j,A_{m+1})\leq|A_j|C$, and
            \begin{equation}
            \label{eq:edges-a}
                e(A_j,A_{m+1})> (\alpha-\beta) |A_j||A_{m+1}|.
            \end{equation}
            Construct an auxiliary bipartite graph $H$ with vertex classes $\fA_j$, $\fA_{m+1}$ by joining $\tilde K_j$ to $\tilde K_{m+1}$ if the edge density between those cliques is greater than $\alpha-2\beta$, that is if
            \[e(V(\tilde K_j), V(\tilde K_{m+1}))> (\alpha-2\beta)j(m+1).\]
            Owing to \eqref{eq:edges-a}, by standard double counting, we have 
            \[e(H)\geq \beta |\fA_j||\fA_{m+1}|.\]
            Consider an edge $f=\tilde K_j \tilde K_{m+1}$ of $H$. We claim that 
            \begin{equation}
            \label{eq:edges-K_j_m+1}
            e(V(\tilde K_j), V(\tilde K_{m+1}))\geq mj.
            \end{equation} 
            Indeed, consider the following two cases.
            
            \textbf{Case 1.} If $t>0$, then recalling that $\alpha=1-\frac{s}{r}>1-\frac1m$, we get
            \[\alpha (m+1)j > \frac{(m-1)(m+1)}{m}j=mj-\frac{j}{m}\geq mj-1.\]
            
            \textbf{Case 2.} If $t=0$, then $\alpha=1-\frac1m$, and by Claim~\ref{c:singular-A_m} we have $j<m$. Then, similarly, \[\alpha (m+1)j = \frac{(m-1)(m+1)}{m}j=mj-\frac{j}{m}> mj-1.\]
            In both cases, for sufficiently small $\beta$, we have $(\alpha-2\beta)j(m+1) > mj-1$, which implies~\eqref{eq:edges-K_j_m+1}.
            In other words, there are at most $j$ missing edges between $\tilde K_j$ and $\tilde K_{m+1}$. Hence, there is a partition $V(\tilde K_{m+1})=M_f\cup N_f$ such that $|M_f|=m+1-j$, $|N_f|=j$ and $V(\tilde K_j)\cup M_f$ span a copy of $K_{m+1}$ in $G$. 
            
            Next, notice that by our assumption on $\beta$, we have $\beta<\binom{m+1}{j}^{-1}$. Then since the number of possible partitions is $\binom{m+1}{j}$, for each $\tilde K_{m+1}\in\fA_{m+1}$ we can choose a partition $V(\tilde K_{m+1})=\tilde M \cup \tilde N$ such that there are at least $\beta \deg_H(\tilde K_{m+1})$ edges $f$ in graph $H$ with $M_f=\tilde M$ and $N_f=\tilde N$. Let $H'$ be the subgraph of $H$ consisting of such edges. We have $$e(H')\geq \beta e(H)\geq \beta^2 |\fA_j||\fA_{m+1}|.$$
            Since, by \eqref{eq:A_j}, $|\fA_j|\geq \frac{C}{m^2}$ and, by \eqref{eq:A_{m+1}},  $|\fA_{m+1}|\geq \frac{C}{m+1}$, the theorem of K\H{o}vari, S\'os and Tur\'an~\cite{kst54}\footnote{This application of the K\H{o}vari, S\'os and Tur\'an Theorem to a graph of density $\beta^2$, as well as the subsequent one, are our main restrictions on $C$.} yields a copy of $K_{s-t,(m-j)s+t}$ in $H'$. We can see that the vertices of this $K_{s-t,(m-j)s+t}$-copy form, in turn, a copy of $Q_j$ in $G$ (the $s-t$ copies of $K_j$ form the $L$-part of $Q_j$, while for each copy $\tilde K_{m+1}$ we take the aforementioned partition $\tilde M \cup\tilde N$ which corresponds to the $M$-part and $N$-part). By replacing the $s-t$ copies of $K_j$ and $(m-j)s+t$ copies of $K_{m+1}$ with $Q_j$ in $\cF$, we can increase $q_j^\cF$ without changing $\phi_{m+1}^\cF$ or $\phi_j^\cF$, contradicting the maximality of $\cF$. 
        \end{proof}

        \begin{claim}
            \label{c:p-B-edges}
            If $h\in[j-1]$, then 
            $$e(A_j,B_h)\leq |A_j|((\alpha-\beta)|B_h|+C).$$
        \end{claim}
        
        \begin{proof}
            The proof is similar to the proof of the previous claim. Again we assume for the sake of contradiction that 
            \begin{equation}\label{eq:B_h}
                |B_h|\geq C \quad \text{ and } \quad e(A_j,B_h)> (\alpha-\beta)|A_j||B_h|.
            \end{equation}
            Like in the proof of Claim~\ref{c:p-A-edges-in-m+1}, we construct an auxiliary bipartite graph $H$ between $\fA_j$ and $\fB_h$ by joining $\tilde K_j$ and $\tilde Q_h$ if the edge density between them is greater than $\alpha-2\beta$, that is, if 
            \[e(V(\tilde K_j), V(\tilde Q_h))> (\alpha-2\beta)j|V(\tilde Q_h)|.\]
            Notice that by the same reasoning as before, we have
            \[ e(H) \geq \beta|\fA_j||\fB_h|.\]
            Recall that $|V(Q_h)| = (m+1-h)r$. Fix an edge $f=\tilde K_j \tilde Q_h$ of $H$. 
            Let 
            \[
            X_f=\{x\in V(\tilde Q_h)\mid V(\tilde K_j) \subseteq N(x)\}.
            \]
            Note that 
            \[|X_f|  \geq e(V(\tilde K_j), V(\tilde Q_h))-(j-1)|V(\tilde Q_h)|\geq ((\alpha-2\beta)j-(j-1))(m+1-h)r.\] 
            We claim that, for sufficiently small $\beta$, the right-hand side in the second inequality above is larger than $(m-j)((m-h)s+t).$ Indeed, we have
            $$(\alpha j -j+1)(m+1-h)(ms+t)-(m-j)(ms-hs+t)= s(m-j)+t(j+1-h)>0,$$
            since we assumed that $h \leq j-1$, and since by Claim~\ref{c:singular-A_m} $j<m$ in case $t=0$.
            We conclude that
            \begin{equation}
             \label{eq:size-Xf}
            \begin{split}
            |X_f| %& \geq e(V(\tilde K_j), V(\tilde Q_h))-(j-1)|V(\tilde Q_h)|\\
            %&\geq j((m-1)s+t)(m+1-h)-(j-1)(m+1-h)(ms+t) \\
           % &=(m+1-h)((m-j)s+t)\\
            & > (m-j)((m-h)s+t).
            \end{split}
            \end{equation}

            If some vertex $x\in X_f$ is contained in the $L$- or $N$-set of $\tilde Q_h$, we can replace $\tilde K_j$ and $\tilde Q_h$ with a copy of $K_{j+1}$ (by shifting the vertex $x$ to $\tilde K_j$), $(m-h)s+t$ copies of $K_{m+1}$, one copy of $K_{h-1}$, and $s-t-1$ additional copies of $K_h$. This increases $\phi_{j+1}^\cF$ without decreasing $\phi_{m+1}^\cF$, which contradicts the maximality of $\cF$. For this reason $X_f$ is contained in the union of the $M$-sets of $\tilde Q_h$. 
            
            Now, by \eqref{eq:size-Xf} there is some index $i(f)\in [(m-h)s+t]$ such that $\tilde M_{i(f)}$ intersects $X_f$ in at least $m-j+1$ vertices. This means that there is a partition 
            \[\tilde M_{i(f)}=C_f\cup D_f\]
            with $|C_f|=m-j+1$ and $|D_f|=j-h$, and such that $C_f\subset X_f$. Since for each $\tilde Q_h$ the number of possible pairs $(i(f), C_f)$ as above is at most $((m-h)s+t){m+1-h\choose j-h} < \beta^{-1}$, we can choose one, say $(i(\tilde Q_h), C(\tilde Q_h))$, for which there are at least $\beta \deg_H(\tilde Q_h)$ edges $f$ with $(i(f), C_f) = (i(\tilde Q_h),C(\tilde Q_h))$.
            Thus, like in the proof of the previous claim, there is a~subgraph $H'$ of $H$ satisfying
            \begin{itemize}
                \item $e(H')\geq \beta e(H) \geq  \beta ^2 |\fA_j| |\fB_h|$,
                \item for any $\tilde Q_h$, there is a "special pair" $(\tilde M_i, \tilde N_i)$ in $\tilde Q_h$ for which there is a unique $\tilde C=C(\tilde Q_h)\subseteq \tilde M_i$, such that $(i,\tilde C)=(i(f),C_f)$ for any edge $f=\tilde K_j \tilde Q_h$ in $H'$. 
            \end{itemize}

            
            By \eqref{eq:A_j} we have $|\fA_j|\geq \frac{C}{m^2}$, and by \eqref{eq:B_h}, $|\fB_h|\geq \frac{|B_h|}{|Q_h|}\geq\frac{C}{(m-h)s+t}$. The K\H{o}vari-S\'os-Tur\'an theorem implies now that $H'$ contains a subgraph isomorphic to $K_{s-t,(m-j)s+t}$. 
            We can take this subgraph, which consists of $s-t$ copies of $K_j$ and $(m-j)s+t$ copies of $Q_h$, and replace them with one copy of $Q_j$, $((m-j)s+t)((m-h)s+t-1)$ copies of $K_{m+1}$ and some additional copies of $K_h$ as follows:
            \begin{itemize}
                \item the $s-t$ copies of $K_j$ become the $L$-part of $Q_j$,
                \item we take the $(m-j)s+t$ special pairs $(\tilde M_i, \tilde N_i)$, each from a different copy of $\tilde Q_h$, and split $\tilde M_i$ into $\tilde M_i = \tilde C \cup \tilde D$, where $\tilde C = C(\tilde Q_h)$; the sets $\tilde C$ become the $M$-part of $Q_j$ and the sets $\tilde D \cup \tilde N$ become the $N$-part of $Q_j$, 
                
                \item the remaining pairs $(\tilde M_k, \tilde N_k)$ are just copies of $K_{m+1}$, 
                \item finally, the $L$-part of $\tilde Q_h$ gives the additional copies of $K_h$.
            \end{itemize} 
            
            This increases $q_j^\cF$, while $\phi_{m+1}^\cF$ and $\phi_{j}^\cF$ remain unchanged, which is again a contradiction to the maximality of $\cF$.
        \end{proof}

        \begin{claim}
            \label{c:p-B-non-neighbours-in-L-N}
            If $x\in A_j$, $h\in[j,m]$, and $\tilde Q_h\in \fB_h$, then $x$ has at least one non-neighbour in each $L$- and $N$-set of $\tilde Q_h$. In particular,
            $$|N_G(x)\cap V(\tilde Q_h)|\leq \alpha|V(\tilde Q_h)|.$$
        \end{claim}
        
        \begin{proof}
            Let $\tilde K_j\in \fA_j$ be the copy of $K_j$ that contains $x$. If $x$ is completely joined to any $L$- or $N$-set of $\tilde Q_h$, we can replace $\tilde Q_h$ and $\tilde K_j$ in $\cF$ with one copy of $K_{h+1}$ (containing $x$), $(m-h)s+t$ copies of $K_{m+1}$, one copy of $K_{j-1}$ ($\tilde K_j - x$) and some additional copies of $K_h$. This increases $\phi_{h+1}^\cF$ without changing $\phi_{m+1}^\cF$, contradicting the maximality of $\cF$.
            Additionally, recalling that $|V(\tilde Q_h)|=(m+1-h)r$, we get
            \begin{align*}|N_G(x)\cap V(\tilde Q_h)|&\leq |V(\tilde Q_h)|-((m-h)s+t)-(s-t) = \alpha |V(\tilde Q_h)|.
            \end{align*}
        \end{proof}

        Next, we define a partition $V(G)=X\cup Y$ by taking 
        \begin{align*}X =A_1\cup\dots\cup A_{m+1}\cup B_1\cup \dots \cup B_{j-1} \quad \text{ and } \quad 
        Y = B_j\cup \dots \cup B_m.\end{align*}
        
        \begin{claim}\label{c:X_is_not_too_large} We have
            $$|X|\leq \eta n.$$
        \end{claim}
        \begin{proof}
            Owing to the minimality of $j$, recall that $|A_h| \leq \frac{C}{m}$ for $h<j$. Then, using all the previous claims, we have
            \begin{align*}
                (\alpha-\delta)|A_j|n & \leq  e(A_j,V(G))\\
                &= \sum_{h=1}^{j-1}e(A_j,A_h)+\sum_{h=j}^{m+1}e(A_j,A_h)+\sum_{h=1}^{j-1} e(A_j, B_h)+\sum_{h=j}^m e(A_j, B_h) \\
                %&\leq p|A_p|\frac{C}{m}+(\alpha-\beta)|A_p|\sum_{h=p}^m |A_h| + |A_p|((\alpha-\beta)|A_{m+1}|+C) + |A_p|\sum_{h=1}^{p-1}((\alpha-\beta)|B_h|+C)+\alpha|A_p|\sum_{h=p}^m|B_h|\\
                & \leq |A_j|\left(\frac{(j-1)C}{m} + (\alpha-\beta)|X|+jC +sm+\alpha|Y|\right)\\
                & \leq |A_j|\left((\alpha-\beta) |X|+\alpha (n-|X|)+(m+2)C\right),
            \end{align*}
            where the term $|A_j|sm$ is needed for case $t=0$ and comes from Claim~\ref{c:singular-A_m}.
            In particular,
            \[|X|\leq \frac{\delta n +(m+2)C}{\beta}\leq \frac{2\delta n}{\beta}\leq \eta n.\]
        \end{proof}

        Next, let $W$ be the union of the $L$- and $N$-sets of the copies of $Q_h$ in $\fB_j\cup\dots\cup\fB_m$. We construct an auxiliary bipartite graph $H$ between $A_j$ and $W$ with edges defined as follows. For a vertex $x\in A_j$ and a copy $\tilde Q_h\in \fB_h$, where $h\in [j,m]$:
        \begin{itemize}
            \item if $|N_G(x)\cap V(\tilde Q_h)|<\alpha |V(\tilde Q_h)|$, then $H$ has no edges from $x$ to $V(\tilde Q_h) \cap W$;
            \item otherwise, by Claim~\ref{c:p-B-non-neighbours-in-L-N} we have $|N_G(x)\cap V(\tilde Q_h)|=\alpha |V(\tilde Q_h)|=(m-h+1)s$ and there are exactly this many non-edges between $x$ and $V(\tilde Q_h)\cap W$, one for each $L$- and $N$-set. These $(m-h+1)s$ non-edges will become the edges in $H$ from $x$ to $V(\tilde Q_h)\cap W$. 
        \end{itemize}

        \begin{claim}\label{c:no-two-H-neighbours}
            No vertex $y\in W$ has two $H$-neighbours belonging to the same clique $\tilde K_j \in \fA_j$.
        \end{claim}
        \begin{proof}
            Assume such a vertex $y$ exists, and let $\tilde K_j\in \fA_j$ and $u,v \in V(\tilde K_j)$ be such that $uy,vy\in E(H)$. Let $\tilde Q_h \in \fB_h$ be the copy of $Q_h$ that contains $y$ and $\tilde Y$ be the corresponding $L$- or $N$-set of $\tilde Q_h$ with $y\in \tilde Y$. By the definition of $H$, both $u$ and $v$ are connected to every vertex in $\tilde Y \setminus \{y\}$ in $G$. Let us remove $y$ from $\tilde Y$ and add $u$ and $v$ instead. By doing so, we replace $\tilde K_j$ and $\tilde Q_h$ by a copy of $K_{h+1}$ (spanned by vertices in $\tilde Y\cup\{u,v\}\setminus \{y\}$), $(m-h)s+t$ copies of $K_{m+1}$, along with a copy of $K_{j-2}$ (i.e.~$\tilde K_j - \{u,v\}$) and $K_1$ (i.e.~vertex $y$), and additional copies of $K_h$. This increases $\phi_{h+1}^\cF$, while not decreasing $\phi_{m+1}^\cF$, thus contradicting the maximality of $\cF$.
        \end{proof}
        
        Finally, set 
        \[ Z=\{z\in W\colon d_H(z)\geq j+1\}. \]

        \begin{claim}
            \label{c:Z-is-large}
            We have $|Z|\geq j(1-\xi)(1-\alpha)n$.
        \end{claim}
        \begin{proof}
            Recall $|A_j|\geq \frac{C}{m}$, which implies $|\fA_j|\geq \frac{C}{jm}\geq \frac{C}{m^2}\geq j(m+1-j)+1$, for an appropriate choice of $C$. Pick a subset $\fA_j'\subseteq \fA_j$ of size $|\fA_j'|=j(m+1-j)+1$ and let $A_j'$ be the union of the vertex sets of the cliques in $\fA_j'$. Look at any vertex $x\in A_j'$. Let $\fD_x$ be the set of all $\tilde Q_h\in \fB_h$ such that $h\in [j,m]$ and $|N_G(x)\cap V(\tilde Q_h)|=\alpha |V(\tilde Q_h)|$.
            By Claim~\ref{c:p-B-non-neighbours-in-L-N}, all other $\tilde Q_h\in \fB_j\cup \dots \cup \fB_m$ satisfy $|N_G(x)\cap V(\tilde Q_h)|<\alpha|V(\tilde Q_h)|\leq (\alpha- \beta)|V(\tilde Q_h)|$ for $\beta$ sufficiently small. Let $D_x$ be the union of the vertex sets of the cliques in $\fD_x$. We have 
            \begin{align*}(\alpha-\delta)n\leq |N_G(x)| &\leq |X|+|N_G(x)\cap D_x|+|N_G(x)\cap Y\setminus D_x|\\
            &\leq \eta n + \alpha |D_x|+(\alpha-\beta)(n-|D_x|),
            \end{align*}
            whence 
            \[|D_x|\geq \frac{\beta-\delta-\eta}{\beta}n\geq (1-\xi^2)n.\]
            Now set 
            \[D=\bigcap_{x\in A_j'}D_x \quad \text{ and } \quad \fD=\bigcap_{x\in A_j'}\fD_x .\]
            Since, for $\xi$ sufficiently large, \[|D|\geq (1-|A_j'|\xi^2)n\geq (1-\xi)n,\]
            it suffices to argue that every $\tilde Q_h\in \fD$ has at least \[j(1-\alpha)|V(\tilde Q_h)|=(m+1-h)sj\] vertices belonging to $Z$. Fix $\tilde Q_h\in \fD$ and let $U$ be one of its $L$- or $N$-sets. We will show that $|U\cap Z|\geq j$, which implies the desired bound on $|V(\tilde Q_h)\cap Z|$. Suppose that $|U\cap Z|\leq j-1$. Since $|U|=h\geq j$, there is a set $U'$ of size $j-1$ such that $U \cap Z\subseteq U'\subseteq U$. Look at any $\tilde K_j \in \fA_j'$. Due to $\tilde Q_h\in \fD\subseteq \bigcap_{x\in V(\tilde K_j)}\fD_x$, every vertex of $\tilde K_j$ has an $H$-neighbour in $U$. Moreover, by Claim~\ref{c:no-two-H-neighbours}, $H$ induces a matching from $V(\tilde K_j)$ to $U$. Therefore, there is an $H$-edge from $\tilde K_j$ to $U\setminus U'$. Applying this argument to every $\tilde K_j \in \fA_j'$ we obtain $j(m+1-j)+1$ edges from $A_j'$ to $U\setminus U'$. But $|U\setminus U'|=h-(j-1)\leq m+1-j$, so there is a vertex $u\in U\setminus U'$ which receives at least $j+1$ of those edges. This implies that $u\in U\cap Z \subseteq U'$, which is a contradiction.
        \end{proof}

        \begin{claim}
            \label{c:Z-has-high-mindegree}
            We have
            \[\delta(G[Z])\geq\frac{j-1-2\xi}{j}|Z|.\]
        \end{claim}
        \begin{proof}
            Using the previous claim and the inequality $\alpha =1-\frac{s}{r}\leq 1-\frac{1}{m+1}$, we can see that
            \begin{align*}
                \delta(G[Z]) - \frac{j-1-2\xi}{j}|Z| &\geq (\alpha-\delta)n-(n-|Z|)-\frac{j-1-2\xi}{j}|Z|\\
                &=\frac{1+2\xi}{j}|Z|-(1-\alpha+\delta)n\\
                &\geq \left((1-\alpha)((1+2\xi)(1-\xi)-1)-\delta\right)n\\
                &\geq \left(\frac{1}{m+1}(\xi-2\xi^2)-\delta\right)n\geq 0.
            \end{align*}
        \end{proof}

        \begin{claim}
            \label{c:no-K_(j+1)-in-Z}
            There are no copies of $K_{j+1}$ in $G[Z]$.
        \end{claim}
        \begin{proof}
        Assume $y_1,\dots, y_{j+1}$ is a clique in $G[Z]$. By the definition of $Z$, there exists a~matching $x_1y_1,\dots , x_{j+1}y_{j+1}$ in $H$. As before, we will perturb the factor $\cF$ to increase its index. For each $i\in [j+1]$, remove the vertex $x_i$ from its copy of $K_j$, and  replace $y_i$ by $x_i$ in the copy of $Q_h$ containing $y_i$. This leaves all the copies of $Q_h$ containing one of the vertices $y_i$, $i\in[j+1]$, intact. We can now add the clique $y_1,\dots, y_{j+1}$ to the factor. This increases $\phi_{j+1}^\cF$,  contrary to the maximality of $\cF$.
        \end{proof}

        We can finally finish the proof of Lemma~\ref{l:6.2chr}.
        If $j=1$, then Claim~\ref{c:Z-is-large} and Claim~\ref{c:no-K_(j+1)-in-Z} show that $Z$ is an independent set of size at least $(1-\xi)(1-\alpha)n=\frac{s}{r}(1-\xi)n>(\frac{s}{r}-\gamma)n$. If $j\geq 2$, then the Andr\'asfari-Erd\H{o}s-S\'os~\cite{aes74} theorem together with Claim~\ref{c:Z-has-high-mindegree} and Claim~\ref{c:no-K_(j+1)-in-Z} show that $G[Z]$ is $j$-partite. This in turn means that there is an independent set whose size, due to Claim~\ref{c:Z-is-large}, is at least $\frac{|Z|}{j}\geq (1-\xi)(1-\alpha)n$. In both cases, we see that \ref{l:6.2chr(ii)} holds.
    \end{proof}  
    
    \subsection{Weighted graphs and packing}

    Throughout this subsection, let us assume that $r=ms+t$ with $0\leq t<s$.

    Recall the definition of a weighted graph and an $\cF$-packing from the Introduction.
    Any graph $F$ can be viewed as a~weighted graph with $w_F \equiv 1$. 
    
    The \textit{residue} of a packing is defined as $|V(G)|- \sum_{i \in I}\sum_{v \in V(F_i)}w_{F_i}(v)$. An \emph{$\mathcal{F}$-factor} is an~$\mathcal{F}$-packing with residue 0. If $\cF$ consists of a single weighted graph $F$, we will write $F$-factor or $F$-packing for short.
    
    Let $T$ denote the weighted complete graph on vertices $\sigma_1,\dots,\sigma_m$ and $\tau$, with vertices $\sigma_i$ having each weight $\frac{s}{r}$ and vertex $\tau$ having weight $\frac{t}{r}$. In case $t=0$, one can think of $T$ as the clique $K_m$ with each vertex bearing weight $1/m$ (formally, we could also exclude $\tau$ from the definition of $T$ in this case). 
    
    We can easily see that the graph $K_{m+1}$ has a~$T$-factor consisting of $m+1$ copies of $T$. Indeed, it suffices to take $f_i(\tau)$ to be a~different vertex for each $F_i$, $i\in[m+1]$. Recall also that, in case $t=0$, the graph $Q_m^{(m)}$ consists of $s$ vertex-disjoint copies of $K_m$. This graph naturally has a $T$-factor, since we can map $\tau$ anywhere without enhancing the total weight of a vertex. It turns out that we can also factorize all other graphs $Q_h^{(m)}$, possibly first rescaling the weights of $T$ by an appropriate rational number.
    Recall that for a weighted graph $F$ and $\phi > 0 $, by $\phi\ltimes F$ we denote the weighted graph obtained by multiplying all the vertex weights of $F$ by $\phi$. The weight function is therefore given by $w_{\phi\ltimes F}(x)=\phi \cdot w_F(x)$ for every $x\in V(F)$. 

    \begin{lemma}\label{l:factoring-Qh}
        The clique $K_{m+1}$ has a $T$-factor. Furthermore, for every $h\in [m]$, there is a rational number $q_h\in \QQ_{>0}$ such that $Q_h^{(m)}$ has a $\{q_h\ltimes T\}$-factor.
    \end{lemma}
    
    \begin{proof}
    The first part has been already justified, and we have also shown that in the case $t=0$, the graph $Q_m^{(m)}$ has a $T$-factor.
    So consider the graph $Q_h=Q_h^{(m)}$ and assume that $t>0$ or $h<m$. Set \begin{equation*}
        x=s-t, \quad y=(m-h)s+t,
    \end{equation*} 
    \begin{equation*}
        q=q_h=\frac{r}{s(m-h+1)xy}, \quad q_1=x\cdot q,\quad q_2=y\cdot q.    
    \end{equation*} 
    Since $q_1\ltimes T$ and $q_2\ltimes T$ have a $q\ltimes T$-factor, it suffices to find a $\{q_1\ltimes T, q_2\ltimes T\}$-factor.
        Consider first the weighted graph $F$ with vertices
        \begin{itemize}
            \item $l_1,\dots,l_h$ of weights $\frac 1y$,
            \item $m_1,\dots, m_{m-h+1}$ and $n_1,\dots,n_h$ of weights $\frac 1x$,
        \end{itemize}
        and all edges except $l_in_j$, $i,j \in [h]$.
        We can factor $Q_h$ with $xy$ copies of $F$ by assigning one copy of $F$ to $L_i\cup M_j \cup N_j$ for each $i\in [s-t]$ and $j\in [(m-h)s+t]$. 
        Therefore, it suffices to factor $F$ with copies of $q_1\ltimes T$ and $q_2\ltimes T$. 
        
        For $i\in[m-h+1]$, let $T_i$ be the copy of $T$ together with a map $f_i$ mapping the vertices of $T_i$ to $l_1,\dots, l_h,m_1,\dots ,m_{m-h+1}$, in such a way that vertex corresponding to $\tau$ in $T_i$ is mapped to $m_i$. Similarly, for $i\in[m-h+1]$, let $T_i^*$ be the copy of $T$ with a map $f_i^*$ mapping the vertices of $T_i^*$  to $m_1,\dots ,m_{m-h+1},n_1,\dots ,n_h$ in such a way that vertex corresponding to $\tau$ in $T_i^*$ is again mapped to $m_i$. 
        Then $q_1\ltimes T_i$ and $q_2\ltimes T_i^*$, $i\in[m-h+1]$, factorise $F$. Indeed, each vertex $l_i$ gets the total weight
        \[ q_1(m-h+1)\frac{s}{r} = \frac1y,\]
        similarly, each vertex $n_i$ gets the total weight
        \[ q_2(m-h+1)\frac{s}{r} = \frac1x,\]
        and each vertex $m_i$ gets the total weight
        \begin{align*} 
        (q_1+q_2)\left(\frac{t}{r} + (m-h)\frac{s}{r}\right) & = \left(\frac{1}{x}+\frac{1}{y}\right) \frac{(m-h)s+t}{s(m-h+1)} \\
        & = \frac{1}{x}\left(1+ \frac{s-t}{(m-h)s+t}\right)\frac{(m-h)s+t}{s(m-h+1)} = \frac1x.
        \end{align*}
    \end{proof}
    
    \begin{cor}\label{cor:T_packing}
    For any $m\in\NN$, there exists $b\in \NN$ such that any collection of vertex-disjoint copies of $K_{m+1}, Q_1^{(m)}, \dots, Q_m^{(m)}$ has a $\left\{\frac 1b\ltimes T\right\}$-factor.
    \end{cor}

    \begin{proof}
        By Lemma~\ref{l:factoring-Qh}, any collection of vertex-disjoint copies of $K_{m+1},Q_1^{(m)},\dots,Q_m^{(m)}$ has a $\left\{T,q_1\ltimes T, \dots , q_m\ltimes T\right\}$-factor, for some $q_1,\dots, q_m\in \QQ_{>0}$. Let $b$ be the common denominator of $q_1,\dots , q_m$. For every $h\in [m]$, $q_h\ltimes T$ can be factorized using $b\cdot q_h$ copies of $\frac 1b\ltimes T$, so we can find a $\frac 1b \ltimes T$-factor for our collection.
    \end{proof}

    \subsection{Stability for Shokoufandeh--Zhao Theorem}
        \label{sec:sz-stability}
        In this section, we prove Theorem~\ref{t:sz-stability}. Let us first introduce the \emph{critical chromatic number} $\chi_{cr}$. Let $H$ be an $r$-vertex graph with chromatic number $m+1$, and let $t$ be the smallest class size in any proper $(m+1)$-colouring of $H$. Let $s$ be defined by $r = ms+t$, and note that this is not necessarily an integer. Then the critical chromatic number of $H$ is given by $\chi_{cr}(H) = \frac{r}{s}$.
        
        We will only use the following fact which follows directly from the definition and can be found in~\cite{sz03}. They refer to the graph $B(m, r, t)$ below as a \textit{bottle} graph.

        \begin{fact}
        Let $H$ be as above, and let $B(m,r,t)$ be the complete $(m+1)$-partite graph with $m$ classes each of size $ms$ and one class of size $mt$. Then $B(m,r,t)$ has an $H$-factor, and $\chi_{cr}(B(m,r,t)) = \frac rs$.
        \end{fact}

        We can now proceed with the proof.

        \begin{proof}
            [Proof of Theorem~\ref{t:sz-stability}]
            Let $H$ be the given graph, and let $r, m, t$ and $s$ be as above. As in the previous subsection, let $T$ denote the weighted complete graph on vertices $\sigma_1,\dots,\sigma_m$ and $\tau$, with $m$ vertices of weight $\frac{s}{r}$ and one vertex of weight $\frac{t}{r}$. Since $\chi_{cr}(H) = \frac{r}{s}$, the graph $G$ satisfies the hypothesis of Lemma~\ref{l:6.2chr}. Using Lemma~\ref{l:6.2chr} (conclusion (i)) and Corollary~\ref{cor:T_packing}, there is an integer $b \in \NN$ such that $G$ has a $\left\{ b^{-1} \ltimes T\right\}$-packing with remainder at most $C$. Furthermore, each copy of $T$ has a $\left\{b_1^{-1} \ltimes B(m,s,t)\right\}$-factor, for some integer $b_1$, which in turn yields a $\left\{b_2^{-1} \ltimes H\right\}$-factor for some integer $b_2$, as required.
        \end{proof}
        

          
    
    \subsection{Constructing the absorber}
        
    In this subsection, let $r=ms+t$ with $1\leq t \leq s$. This is a small change of parametrization compared to the previous subsections, as the case $t=s$ will now be parametrised as $r=(m+1)s$ instead of $r=ms$. 

    \begin{definition} 
        \label{def:absorber}
        A clique $K_g$ in $G$ is called \textit{good} for a vertex $v$ if $v$ has at least $g-2$ neighbours in $K_g$. A \textit{$(k, m, \delta)$-absorber} in $G$ is a collection $\absorber$ of vertex-disjoint copies of  $K_{m+k}$ in $G$ covering at most $10(k+m)\delta n$ vertices and such that for each pair of vertices $u,v\in V(G)$, $\absorber$ contains at least $\delta^2 n$ cliques which are good for both $u$ and $v$. 
    \end{definition}
    
    \begin{lemma}\label{l:absorber}
        Let $0<\delta<(4mr)^{-1}$ and suppose that $\delta(G)\geq \left(1-\frac sr -\delta\right)n$.
        \begin{enumerate}[label=(\roman*)]
            \item\label{l:absorber(i)} If $t < s$, then $G$ contains a~$(1, m, \delta_2)$-absorber for some constant $\delta_2=\delta_2(r)>0$. 
            
            \item\label{l:absorber(ii)} Suppose $t = s$, and that every subset $X\subset V(G)$ with $|X|\geq\left(\frac1{m+1}-m\delta\right)n$ satisfies $e_G(X)\geq\delta_1 n^2$ for some $\delta_1>0$. Then $G$ contains a~$(2, m, \delta_2)$-absorber for some constant $\delta_2=\delta_2(r,\delta_1)>0$.
            
        \end{enumerate}
    \end{lemma}
        
    \begin{proof}
        For both cases we will need two auxiliary claims.
        \begin{claim}\label{c:absorber_abstract_claim}
            Let $\ell = n^{O(1)}$, and let $\cK_1, \ldots, \cK_\ell$, be collections of $K_{g}$-cliques in a graph $G$ satisfying $|\cK_i| \geq \xi n^{g}$. Then there is a collection $\absorber$ of vertex-disjoint copies of $K_g$ covering at most $\xi g n$ vertices such that 
            $$|\cA \cap \cK_i| \geq \xi^2 n/50$$
            for $i \in [\ell]$.
        \end{claim}

        \begin{proof}            
            Suppose without loss of generality that $|\cK_i|=\xi n^{g}$ for each $i\in[\ell]$. Let $\cK=\bigcup_i \cK_i$, and note that $|\cK|\leq n^g$ since $\cK$ consists only of copies of $K_g$. Let $\cA'$ be a random subfamily of $\cK$ containing each element of $\cK$ independently with probability $\frac{\xi}{10}n^{1-g}$. We have \[\E[|\cA'|]=\frac{\xi}{10}n^{1-g}|\cK|\leq \frac{\xi}{10}n, \quad \E[|\cA'\cap \cK_i|]=\frac{\xi}{10}n^{1-g}\xi n^g=\frac{\xi^2}{10}n.\] By Chernoff bounds we also have $\pr {|\cA'| >\frac{\xi}{5}n}\leq e^{-\frac{\xi}{20}n}$ and $\pr{|\cA\cap \cK_i|< \frac{\xi^2}{20}n}\leq e^{-\frac{\xi^2}{80}n}$ for each $i\in [\ell]$, so by the union bound, the probability of any of these events occurring is at most $(\ell+1) e^{-\frac{\xi^2}{80}n}=o(1)$. Next, we bound the number of intersecting pairs of copies of $K_g$ in $\cA'$. We have 
            \[\E[|\{(K_1,K_2)\in \cA' \mid K_1\cap K_2\neq \emptyset\}|]\leq n^{2g-1}\left(\frac{\xi}{10}n^{1-g}\right)^2=\frac{\xi^2}{100}n.\]
            By Markov's inequality, with probability at least $1/2$ the number of such pairs is~at most $\left(\xi^2/50\right)n$. Let $\cA$ be the collection of $K_g$ copies achieved by removing one of the cliques in each intersecting pair. Notice that $|\absorber|\leq |\absorber'|< \xi n$, so $\absorber$ covers at most $\xi g n$ vertices. Moreover, with probability at least $1/2-o(1)$, the collection $\cA$ satisfies \[|\cA \cap \cK_i|\geq \frac{\xi^2}{20}n-\frac{\xi^2}{50}n\geq \frac{\xi^2}{50}n\] for each $i\in[\ell]$. Hence, there exists a family $\cA$ which satisfies the statement of the Lemma.
        \end{proof}
        \begin{claim}\label{c:many_common_neighbours_(i)}
            In both case (i) and (ii), if $\delta(G)\geq \left(1-\frac sr -\delta\right)n$, any set $S$ of at most $m$ vertices of $G$ has at least $\left(\frac{t}{r}-m\delta\right)n$ common neighbours.
        \end{claim}
            
        \begin{proof}
            It is enough to upper-bound the number of vertices that are a~non-neighbour of at least one vertex in $S$. This number is at most
        $$mn \left(\frac{s}{r}+\delta \right)= \frac{n(ms+mr\delta)}{r} = n\left(1-\frac{t-mr\delta}{r} \right).$$
        %\geq n\left(1- \frac{t}{2r} \right),$$
               % \nkb{recalling that $\delta m \leq \frac{1}{4r}$}.
        \end{proof}

        Suppose now that we are in case \ref{l:absorber(i)}. We construct the collections $\cK_i$ as follows. 
        Let $\xi_1=\frac{t}{2r}$ and $\xi=\frac{\xi_1^{m+1}}{(m+1)!}$. Note that $\left(\frac{t}{r}-m\delta\right)\geq \xi_1$. Order the pairs of vertices of $G$ arbitrarily and suppose $\{u,v\}$ is the $i$-th pair of vertices. Using Claim~\ref{c:many_common_neighbours_(i)} repeatedly, we can find vertices $w_1,\dots,w_{m-1}$ such that $w_k\in N_G(\{u,v,w_1,\dots,w_{k-1}\})$, for all $k\in[m-1]$. Notice that the set $\{w_1,\dots,w_{m-1}\}$ induces in $G$ a copy of $K_{m-1}$. By a simple counting argument, there are at least $\frac{\xi_1^{m-1}}{(m-1)!}n^{m-1}$ different choices for the set $\{w_1,\dots,w_{m-1}\}$, i.e.~we can find a collection of $\frac{\xi_1^{m-1}}{(m-1)!}n^{m-1}$ different copies of $K_{m-1}$ that are contained in $N_G(u)\cap N_G(v)$. 
        Next, extend those copies of $K_{m-1}$, again using Claim~\ref{c:many_common_neighbours_(i)} twice, to get copies of $K_{m+1}$. Each of these has at least $m-1$ vertices from $N_G(u)\cap N_G(v)$, and there are at least $\frac{\xi_1^{m+1}}{(m+1)!}n^{m+1}$ of them, by the same counting argument as before. This gives a collection $\cK_i$ of $K_{m+1}$-cliques with $|\cK_i|\geq \xi n^{m+1}$, for each $i\in [\binom n2]$. Applying Claim~\ref{c:absorber_abstract_claim} with $\ell={n\choose 2}$ and $g=m+1$, we get a collection $\cA$ which is a $(1,m,\delta_2)$-absorber for $\delta_2\leq \xi/10$.

     
          
          
        
        Next, we prove \ref{l:absorber(ii)}. Let $\xi_1=\min\left\{\delta_1,\frac{t}{2r}\right\}$ and $\xi=\frac{\xi_1^{m}}{(m+2)!}$. Notice that by Claim~\ref{c:many_common_neighbours_(i)} and the assumptions of \ref{l:absorber(ii)}, we have the following property
        \begin{itemize}
            \item[($\ast$)] for any set $S$ of $m$ vertices, $N_G(S)$ spans at least $\delta_1n$ edges in $G$.
        \end{itemize} 
        As before, using Claim~\ref{c:many_common_neighbours_(i)}, we can find $\frac{\xi_1^{m-2}}{(m-2)!}n^{m-2}$ copies of $K_{m-2}$ contained in $N_G(u)\cap N_G(v)$. Using ($\ast$), these $K_{m-2}$-copies can be extended to $K_{m}$-copies in $N_G(u)\cap N_G(v)$, and, once again using ($\ast$), these $K_m$-copies can in turn be extended to $K_{m+2}$-copies which are good for both $u$ and $v$. Hence, we get a collection $\cK_i$ of $K_{m+2}$-copies with $|\cK_i|\geq \xi n^{m+2}$. Applying Claim~\ref{c:absorber_abstract_claim} with $\ell= \binom n2$ and $g=m+2$, we get a collection $\cA$ which is a $(2,m,\delta_2)$-absorber for $\delta_2\leq \xi/10$.
    \end{proof}
    
    Next, we partition our graph with an absorber.

    \begin{lemma}\label{l:covering-with-absorber}
        Suppose $r=ms+t$ with $1<t\leq s$. Let $g = 0$ if $t<s$ and $g=1$ if $t=s$. There exists a constant $C>0$ such that for all $\gamma>0$ the following holds. There exist constants $\delta , \delta_2>0$ and $n_0\in\mathbb{N}$ such that every graph $G$ on $n$ vertices, $n\geq n_0$, with $\delta(G)\geq \left(1-\frac{s}{r}-\delta\right)n$, satisfies one of the following two conditions.
        \begin{enumerate}[label= (\roman*)]
            \item\label{l:covering-with-absorber(i)} There is a collection $\fA$ of vertex-disjoint copies of $K_{m+1+g}, Q_1^{(m+g)}, \dots , Q_{m}^{(m+g)}$ covering all but at most $C$ vertices such that $\fA$ includes a $(1+g,m,\delta_2)$-absorber, or
            \item\label{l:covering-with-absorber(ii)} There is a subset $A\subset V(G)$ with $|A|=\frac{sn}{r}$ such that $e_G(A) < \gamma n^2$. 
        \end{enumerate}
    \end{lemma}
    
    \begin{proof}
        Suppose first that $t<s$, which implies $g=0$.
        Let $C\gg r$ and $n_0^{-1}\ll \delta_{\ref{l:6.2chr}} \ll \gamma, r^{-1}$, where $\delta_{\ref{l:6.2chr}}$ is the constant obtained by applying Lemma~\ref{l:6.2chr} with $\gamma_{\ref{l:6.2chr}}=\gamma/2$. Assume that $2\delta \leq \delta_{\ref{l:6.2chr} }$. Moreover, let $\delta_2$ be sufficiently small so that $30m\delta_2 \leq \delta_{\ref{l:6.2chr} }$ and the conclusion of Lemma~\ref{l:absorber}~\ref{l:absorber(i)} holds. Due to Lemma~\ref{l:absorber}, there is a $(1,m,\delta_2)$-absorber $\fU$ for $G$. Let $G'=G-\fU$ be the graph obtained by removing vertices from the copies of $K_{m+1}$ in $\fU$ and set $n'=|V(G')|$. Since $n'\geq (1-10(m+1) \delta_2)n$ we have $\delta(G')\geq \delta(G) - 10(m+1) \delta_2 n\geq \left(1- \frac sr - \delta_{\ref{l:6.2chr} } \right)n'$. Now apply Lemma~\ref{l:6.2chr} to $G'$. We either get a collection of vertex disjoint copies of $K_{m+1},Q_1^{(m)},\dots ,Q_m^{(m)}$ covering all but at most $C$ vertices of $G'$, which implies \ref{l:covering-with-absorber(i)}, or we get an independent set of size $\left(\frac sr - \frac \gamma 2\right)n'>\left(\frac sr - \gamma\right)n$, which implies \ref{l:covering-with-absorber(ii)}, since we can extend this subset by adding arbitrary $\gamma n$ vertices.
        
        Suppose now that $t=s$, so that $g=1$.
        Similarly to the previous case, take $C\gg r$ and $n_0^{-1} \ll \delta,\delta_2 \ll \gamma, r^{-1}$. Assume that every set of size $\frac{sn}{r}=\frac{n}{m+1}$ contains at least $\gamma n^2$ edges. Applying Lemma~\ref{l:absorber}~\ref{l:absorber(ii)} to $G$ we get a  $(2,m,\delta_2)$-absorber $\fU$ for $G$. Let $G'=G-\fA$ be the graph obtained by removing vertices from the copies of $K_{m+2}$ in $\fU$ and set $n'=|V(G')|$. Since $n'\geq (1-10(m+2)\delta_2)n$ we have $\delta(G')\geq \delta(G) - 10(m+2)\delta_2 n\geq \left(1- \frac {1}{m+1} - \delta-10(m+2)\delta_2\right)n'$. Now apply Corollary~\ref{cor:6.2chr} to $G'$. We either get a collection of vertex disjoint copies of $K_{m+2},Q_1^{(m+1)},\dots ,Q_{m+1}^{(m+1)}$ covering all but at most $C$ vertices of $G'$, which implies \ref{l:covering-with-absorber(i)}, or we get an independent set of size $\left(\frac sr - \frac \gamma 2\right)n'>\left(\frac sr - \gamma\right)n$, which would imply \ref{l:covering-with-absorber(ii)}, leading to a contradiction.
    \end{proof}

%************************* The non-extremal case ****************************
    
    \section{The fully non-extremal case}\label{sec:non-extremal} 
        
    In this section, we prove Theorem~\ref{t:non-extremal-conforming}. To avoid the proof getting too long, we break it down into two arguments. The first part of the argument are standard manipulations in the reduced graph after applying Lemma~\ref{l:covering-with-absorber}. The second part, Section~\ref{sec:balancing}, is more delicate and requires new ideas. %since it is not immediately clear how to deal with the leftover vertices and obtain parts of exactly the appropriate size.
        
    \begin{lemma}\label{l:partition-reduced-graph}
        Let $\gamma, \zeta$ and $r \in \mathbf{N}$ be given with $0 < \zeta \ll \gamma$. Take any constants $\delta=\delta(\gamma)$, $d=d(\gamma)$, $\eps=\eps(d)$ satisfying the following hierarchy \begin{equation}\label{eq:partition-reduced-graph-constants}
            \zeta \ll \eps \ll d \ll \delta,\delta_2 \ll \gamma, r^{-1}. 
        \end{equation}
         Let $n_0=n_0(\eps)$ be sufficiently large. Let $r=ms+t$ with $m\geq 1$, $0< t\leq s$. Let $g = 0$ if $t<s$ and $g=1$ if $t=s$.
            
        Assume that $G$ is an $n$-vertex graph, $n\geq n_0$, with $\delta(G)\geq \left(1-\frac{s}{r}-\frac \delta 2\right)n$, and that every subset of $V(G)$ of size $\frac{sn}{r}$ contains at least $\gamma n^2$ edges. 
        Then $G$ has a vertex partition into sets $S_{i,j}$, $U_{k, \ell}$ and $X$, where $i \in [L_1], j \in [m+1]$, $k\in [L-L_1]$ and $\ell \in [m+1+g]$, $L_1\geq (1-10r\delta_2)L$, with the following properties. Let $R$ be the graph whose vertices are the sets $S_{i,j}$ and $U_{k,\ell}$, where two vertices $A$ and $B$ are connected if $(A, B)$ is an $(\eps, d^+)$-regular pair. Then
	\begin{enumerate}
            
            \item\label{l:part-w-r-f_(i)} For each $i \in [L_1]$ and $j < j' \in [m+1]$, the edges $S_{i,j}S_{i,j'}$ are in $R$. Similarly, for each $k\in [L-L_1]$ and $\ell < \ell' \in [m+1+g]$, the edges $U_{k,\ell}U_{k,\ell'}$ are in $R$. 
			
            \item\label{l:part-w-r-f_(ii)} 
            For each $i \in [L_1]$,
			$$|S_{i,1}|= |S_{i,2} |= \ldots = |S_{i,m}| = \frac {s}{t} \cdot |S_{i,m+1}|,$$
            $$|S_{i,1}\cup S_{i,2}\cup \dots \cup S_{i,m+1}| = \frac{n}{L}(1 \pm \eps),$$
            and $(S_{i,1},\dots, S_{i,m+1})$ is a $(\eps,d^+,d)$-super-regular tuple.

            \noindent
            For each $k \in [L-L_1]$,
			$$|U_{k,1}|= |U_{k,2} |= \ldots = |U_{k,m+1+g}|,$$
            $$|U_{k,1}\cup U_{k,2}\cup\dots\cup U_{k,m+1+g}|= \frac{n}{L}(1 \pm \eps ),$$
            and $(U_{k,1},\dots, U_{k,m+1+g})$ is a $(\eps,d^+,d)$-super-regular tuple.
			
            \item\label{l:part-w-r-f_(iii)} For $A, B \in V(R)$, there are at least $d L$ indices $k \in [L-L_1]$ such that both $A$ and $B$ have at least $m-1+g$ $R$-neighbours in $\{U_{k,1},U_{k,2},\ldots, U_{k,m+1+g} \}$.
			
            \item\label{l:part-w-r-f_(iv)} $|X|< \eps n/2$.
            
	\end{enumerate}
    \end{lemma}
    
    \begin{proof}
        Let $C=C(r)$ be the constant given by Lemma~\ref{l:covering-with-absorber}. 
        Let $d_1 = 16d$, $\eps_1=\eps^2$ and $\ell_0=\lceil 2C/\eps\rceil$. Let $n_0=n_0(\eps)$ %and $L_0=L_0(\eps)$ 
        be given by Lemma~\ref{l:regularity}.
        First, apply the Regularity Lemma (Lemma~\ref{l:regularity}) with $\ell_0$,  $\eps_{\ref{l:regularity}}=\eps_1$, $d_{\ref{l:regularity}}=d_1$ and $\alpha_{\ref{l:regularity}}=1-\frac{s}{r}-\frac{\delta}2$ to $G$, to obtain an $\eps_1$-regular partition $V_0 \cup V_1 \cup \dots \cup V_{\ell'}$ with \begin{equation}\label{eq:ell'}
            \ell' \geq \ell_0 \geq \frac{2C}{\eps_1},
        \end{equation} and let $R$ be the $(\eps_1,d_1)$-reduced graph corresponding to that partition. Note that $\delta(R) \geq \left(1-\frac{s}{r}-\delta\right)\ell'$. We claim that every set of size at least $\frac sr \ell'$ in $R$ spans at least $\frac \gamma 4 \ell'^2$ edges in $R$. Indeed, it is enough to consider sets of size exactly $\frac{s}{r}\ell'$, so suppose that $I$ is such a set, i.e.~$|I| = \frac sr \ell'$, and $I$ spans in $R$ less than $\frac \gamma 4\ell'^2$ edges. Let $V_I = \bigcup_{i\in I} V_i$ and note that for each $i\in I$ we have $(1-\eps_1)\frac n{\ell'}\leq|V_i|\leq\frac{n}{\ell'}$. Then
        \begin{align*} e(G[V_I]) & \leq {|I|\choose 2}|V_1|^2(d_1+\eps_1)+\frac{\gamma}{4}\ell'^2|V_I|^2 + |I|{|V_1|\choose 2} \\ &\leq \frac{1}{2}|I||V_1|^2\left(|I|(d_1+\eps_1)+1 \right)+\frac\gamma4 n^2 \leq \frac34\gamma n^2.
        \end{align*}
        Moreover,
          \[ |V_I| \geq \frac{s}{r} (1-\eps_1)n \geq \left(\frac{s}{r} - \frac{\gamma}8\right)n.\]
          If we now take an arbitrary set of vertices $W \supset V_I$ of size $\frac{sn}{r}$, then 
          \[ e(G[W]) \leq e(G[V_I]) + \frac{\gamma}8n^2 < \gamma n^2,\]
          which contradicts the assumption on $G$.
          Note that this also implies that $R$ has no independent set of size $\left(\frac sr-\frac \gamma4\right)\ell'$.

          Recall that $g = 0$ if $t<s$ and $g=1$ if $t=s$. We can apply Lemma~\ref{l:covering-with-absorber} to $R$ to get a collection $\fA$ consisting of vertex disjoint copies of $K_{m+1+g}, Q_1^{(m+g)}, \ldots, Q_m^{(m+g)}$ and a~$(1+g, m, \delta_2)$-absorber (where $\delta_2 = \delta_2(\gamma) >0$). Let $\fU\subset \fA$ be the subcollection consisting of copies of $K_{m+1+g}$ in the absorber, and let $\fS = \fA \setminus \fU$ be the remaining part. Note that by the definition of the absorber, $|V(\fU)|\leq 10(m+1+g)\delta_2|V(R)|=10(m+1+g)\delta_2\ell'$. Moreover, let $Y$ be the set of vertices of $R$ not covered by $\fA$ and note that $|Y|\leq C$. The next step is to subdivide the sets in $V(\fS)$. To this end, we use the $\left\{\frac 1b\ltimes T\right\}$-factor given by Corollary~\ref{cor:T_packing}.

          The $\left\{\frac 1b\ltimes T\right\}$-factor of $\fS$ is given by a collection of vertex sets
          \[ \left(\{\sigma_1^i,\dots,\sigma_{m+1}^i\}\right)_{i\in[L_1]}\] 
          for a specific $L_1$ to be determined in a moment, and a function
          \[ f: \bigcup_{j\in[m+1],\,i\in[L_1]} \{\sigma_{j}^{i}\} \to V({\fS})\]
          such that $f(\sigma_j^i)f(\sigma_{j'}^i)$ is an edge of $R$ for each $i\in[L_1]$ and $j,j'\in[m+1]$, $j\neq j'$. 
          %(Here, $\sigma_{m+1}^{i}$ corresponds to the vertex $\tau$ in $T$.) 
          Moreover, since the total weight of $\frac 1b\ltimes T$ is $\frac 1b$, and a $\left\{\frac 1b\ltimes T\right\}$-factor has residue $0$, we have $L_1 = b|V(\fS)|$. We now define a family of sets $S'_{i,j}\subseteq V(G)$ corresponding to vertices $\sigma^i_j$, for $i\in[L_1]$ and $j\in[m+1]$. For a set $S\in V(\fS)$, each vertex $x\in S$ randomly chooses one of the sets $S'_{i,j}$ if and only if $f(\sigma^i_j)= S$. Moreover, if this is the case, $x$ chooses $S'_{i,j}$ with probability 
          \[\frac 1b\cdot w(\sigma^i_j) = \begin{cases}
              \frac{s}{br}, & j\in[m] \\
              \frac{t}{br}, & j = m+1.
          \end{cases}\]
        Since $f$ is a $\left\{\frac 1b\ltimes T\right\}$-factor, this selection gives us a probability measure for each vertex $x \in \bigcup_{S\in V(\fS)} S$. Hence, using Chernoff bounds, we get that w.h.p.
        \[ |S'_{i,j}| = \begin{cases}
              \frac{s}{br}|V_1|(1\pm\eps_1), & j\in[m] \\
              \frac{t}{br}|V_1|(1\pm\eps), & j = m+1,
          \end{cases}\]
          and
          \[ |S'_{i,1}\cup S'_{i,2} \cup \dots \cup S'_{i,m+1}| = \frac 1b|V_1|(1\pm\eps_1).\]
         Next, we split the sets in $V(\fU)$ into sets $U'_{k,\ell}$, for $k\in[L_2]$ and $\ell\in[m+1+g]$, using a~similar approach. This time each set $U\in V(\fU)$ is split randomly into $b(m+1+g)$ sets so that we get $L_2 = b(m+1+g)|\fU| = b|V(\fU)|\leq 10b(m+1+g)\delta_2\ell'$ and
         \[ |U'_{k,\ell}| = \frac{1}{b(m+1+g)}|V_1|(1\pm\eps_1),\]
         and
         \[ |U'_{k,1}\cup U'_{k,2} \cup \dots \cup U'_{k,m+1+g}| = \frac 1b|V_1|(1\pm\eps_1).\]
         Note that 
         \[ L := L_1+L_2 = b|V(\fS)| + b|V(\fU)| = b(\ell' - C),\]
         hence
         \[ L_1 \geq \left(1-10r\delta_2\right)L\]
         and
         \[ \frac 1b|V_1|(1\pm\eps_1) = \frac{n}{L}(1\pm 3\eps_1).\]
         
        Next, by Lemma~\ref{l:slicing}, all pairs stemming from the edges of $R$ are $\left(\eps_1/\beta, \left(d_1/4\right)^+\right)$-regular with $\beta = \min\left\{\frac{t}{2br}, \frac{1}{2b(m+1+g)}\right\}\leq\frac12$. This includes all pairs $S'_{i,j}S'_{i,j'}$ and $U'_{k,\ell}U'_{k,\ell'}$, but also some pairs $S'_{i,j}U'_{k,\ell}$. 

        Now to get super-regularity consider first a tuple $(S'_{i,1},\dots,S'_{i,m+1})$. Using Lemma~\ref{l:super-reg-k-tuple}, we can discard at most $(m+1)\frac{\eps_1}{\beta}$ vertices from each part and pass to subsets $S''_{i,j}\subset S'_{i,j}$. In a similar way we can discard at most $(m+1+g)\frac{\eps_1}{\beta}$ vertices from each part $U'_{k,\ell}$ and pass to subsets $U''_{k,\ell}\subseteq U'_{k,\ell}$. After this procedure the pairs in corresponding tuples are $\left(2\eps_1/\beta, \left(d_1/8\right)^+, d_1/8\right)$-super-regular.
          
        Finally, we again discard at most $\eps_1^{3/4}$-fraction of vertices from each part to get the correct sizes, i.e.~we pass to sets $S_{i,j}\subseteq S''_{i,j}$ and $U_{k,\ell}\subseteq U''_{k,\ell}$ satisfying \[|S_{i,1}|= \ldots = |S_{i,m}| = \frac {s}{t} \cdot |S_{i,m+1}|\] 
        for each $i\in[L_1]$ and 
        \[|U_{k,1}|= \ldots = |U_{k,m+1+g}|\] 
        for each $k\in[L_2]$. After this step all super-regular pairs remain super-regular, but with different parameters, namely they are $\left(\sqrt{\eps_1}, \left(d_1/16\right)^+,d_1/16\right)=\left(\eps,d^+,d\right)$-super-regular. Similarly, all other regular pairs remain regular with parameters $\left(\sqrt{\eps_1}, \left(d_1/16\right)^+\right)=\left(\eps,d^+\right)$.
         
         Let $\cS = \bigcup_{i,j}\{S_{i,j}\}$ and $\cU = \bigcup_{k,\ell}\{U_{k,\ell}\}$. We define a new graph $R'$ on vertex set $\cS\cup\cU$ where edges correspond to $\left(\eps, d^+\right)$-regular pairs. Note that by our construction, $R'$ satisfies properties \ref{l:part-w-r-f_(i)} and \ref{l:part-w-r-f_(ii)} from the lemma statement. Intuitively, one can think of $R'$ as a blow-up of $R$.

         We now show property \ref{l:part-w-r-f_(iii)}. Take any $A, B\in \cS\cup\cU$. Let $V_a$ and $V_b$ be the original parts of the initial partition such that $A\subset V_a$ and $B \subset V_b$. Since $\fU$ is the $(1+g,m,\delta_2)$-absorber in $R$, there are at least $\delta_2^2\ell'$ copies $\tilde{K}$ of $K_{m+1+g}$ in $\fU$ which are good for both $V_a$ and $V_b$, that is both $V_a$ and $V_b$ have at least $m-1+g$ $R$-neighbours in $\tilde{K}$. Let $\tilde{U}_\ell$, $\ell\in[m+1+g]$, be the vertices of $\tilde{K}$ and $\tilde{U}_{k,\ell}$, $k\in\left[b(m+1+g) \right]$, $\ell\in[m+1+g]$ be the sets in $\cU$ with $\tilde{U}_{k,\ell} \subset \tilde{U}_\ell$. Then if $V_a\tilde{U}_
        \ell$ is an edge in $R$, that is $(V_a,\tilde{U}_
        \ell)$ is an $(\eps_1, d_1)$-regular pair, then $(A, \tilde{U}_{k,\ell})$ is a $\left(\eps, d^+\right)$-regular pair, which means that $A\tilde{U}_{k,\ell}$ is an edge in $R'$. The same is true for $B$. Hence, taking $d<\frac 12 \delta_2^2,$ there are at least $\delta_2^2\ell'b(m+1+g) \geq dL$ indices $k\in[L_2]$, such that both $A$ and $B$ have at least $m-1+g$ $R'$-neighbours in $\{U_{k,1}, U_{k,2},\dots,U_{k,m+1+g}\}$. 

        Finally, to prove \ref{l:part-w-r-f_(iv)}, let $X$ be the set of all vertices of $G$ that do not belong to any of the sets $S_{i,j}$ nor $U_{k,\ell}$. By our construction, $X$ consists of:
         \begin{itemize}
             \item vertices in $V_0$ (the initial left-over set from the Regularity Lemma),
             \item vertices discarded to get super-regularity,
             \item $\bigcup_{i\in Y}V_i$ (vertices not covered by the collection $\fA$),
             \item vertices discarded to get correct sizes of the parts for condition \ref{l:part-w-r-f_(ii)}.
         \end{itemize}
         Hence, using \eqref{eq:ell'},
         \[ |X|< \eps_1 n + \frac{(m+1+g)\eps_1}{\beta}n + C\frac{n}{\ell'}(1\pm\eps_1) + \eps_1^{3/4}n \leq \sqrt{\eps_1} n/2 = \eps n/2.\]
        \end{proof}



    \subsection{Size adjustment lemmas}


    The following lemmas are all elementary and will be used to adjust the sizes of the partition sets by removing a constant or small fraction of vertices to obtain appropriate size ratios. 
    
    \begin{lemma}\label{l:equalize-(m+1)-tuples-to-s/t}
        Let $S_1,\dots , S_m , T$ be disjoint sets of vertices and let $M\coloneq |S_1\cup\dots\cup S_m\cup T|$. Let $0<\epsilon\ll r^{-1}$. Suppose that $M\gg r$, $M\equiv 0 \pmod{r}$ and $\left(\frac{s}{r}-\epsilon\right) M\leq|S_i|\leq \frac{s}{r}M$ for all $i\in [m]$.
        Consider the following move:
        \begin{itemize}
            \item $P_k$ move: for an index $k\in [m]$ remove $s-1$ vertices from $S_k$, $s$ vertices from each $S_i$, $i\neq k$, and $t+1$ vertices from $T$.
        \end{itemize}   
        Using the above type of move, we can get new sets $S_1'\subseteq S_1,\dots,S_m'\subseteq S_m,T'\subseteq T$ such that $|S_1|=\dots =|S_m|=\frac{s}{t}|T|$ and $|S_1'\cup\dots\cup S_m'\cup T'|\geq (1-mr\epsilon)M.$
    \end{lemma}

    \begin{proof}
        Suppose we can obtain the required sets $S_i'$, $i\in[m]$, and $T'$ by performing $x_k$ moves $P_k$ for $k\in[m]$. This leads to a system of linear equations
        \begin{align}\label{eq:system_x_k}
            |S_1'|&=\frac{s}{t}|T'| \nonumber\\
            &\,\,\,\,\vdots \\
            |S_m'|&=\frac{s}{t}|T'|, \nonumber
        \end{align}
        where $|S_k'|=|S_k|-s\sum_{i=1}^m x_i+x_k$ for $k\in[m]$, and $|T'|=|T|-(t+1)\sum_{i=1}^m x_i$. So it is enough to check if \eqref{eq:system_x_k} has a solution in non-negative integers. Indeed, there is such a~solution given by
        \[x_k=\frac{s}{r}M-|S_k|, \quad \text{ for }k\in[m].\] Since $\frac{s}{r}M-|S_k|\leq \epsilon M$, and in a single move we remove exactly $r$ vertices, the total number of removed vertices is at most $mr\epsilon M$ vertices.
    \end{proof}

     \begin{lemma}\label{l:absorber-divisibility-by-r}
         Let $U_1,\dots  , U_{m+1}$ be disjoint sets of vertices and let $M\coloneq |U_1\cup\dots\cup  U_{m+1}|$. Let $0<\epsilon \ll 1 $. Suppose that $M\gg r$,  $M \equiv 0 \pmod{r}$ and $|U_i|=\frac{M}{m+1}(1\pm\epsilon)$ for all $i\in [m]$.
        Consider the following two types of moves:
        \begin{itemize}
            \item $P_j$ move: for an index $j\in [m]$ remove $t$ vertices from $U_j$ and $s$ vertices from each $U_i$, $i\in[m+1]\setminus\{j\}$.
            \item $P_{j,k}$ move: for indices $j,k \in [m]$ remove $t+1$ vertices from $U_j$, $s-1$ vertices from $U_k$ and $s$ vertices from each $U_i$, $i\in[m+1]\setminus\{j,k\}$.
            \end{itemize} 
        Using the above moves, we can get new sets $U_1'\subseteq U_1,\dots,U_{m+1}'\subseteq U_{m+1}$ such that each $|U_i'|$ is divisible by $r$ and $|U_1'\cup\dots\cup U_{m+1}'|\geq M -(m+1)^2r^2$.
     \end{lemma}

    \begin{proof}
        We will perform the following algorithm using the allowed moves. Suppose that we have reached an $(m+1)$-tuple $W_1,\dots,W_{m+1}$. Associate an $(m+1$)-tuple $(y_1,\dots,y_{m+1})$ of remainders in $\{-(r-1),\dots,r-1\}$ to $W_1,\dots, W_{m+1}$ so that $|W_i|\equiv y_i \pmod r$ and $\sum y_i=0$. Furthermore, assume that this choice is minimal in terms of $\sum |y_i|$. If $\sum|y_i|=0$ then we stop the algorithm. Otherwise, without loss of generality, assume that $y_1\geq \dots \geq y_{m+1}$ and note that $y_1>0>y_{m+1}$. Apply the moves $P_{1,m+1}$, $P_2,\dots ,P_m,P_{m+1}$. With this sequence of moves, we have removed $r+1$ vertices from $W_1$, $r-1$ vertices from $W_{m+1}$, and $r$ vertices from each $W_i$, $i=2,\dots ,m$. Thus, we have reduced $|y_1|$ and $|y_{m+1}|$ each by $1$, and the sum $\sum |y_i|$ by $2$. By iterating this process, we can reach $\sum |y_i|=0$ in a finite number of steps and finish the algorithm.

        Note that at the end of the algorithm we reach an $(m+1)$-tuple with each set of size divisible by $r$.
        Moreover, since in each iteration of the algorithm we remove $(m+1)r$ vertices, and the number of iterations is at most $\frac12(r-1)(m+1)<r(m+1)$, we have removed in total at most $(m+1)^2r^2$ vertices. 
    \end{proof} 

    \begin{lemma}
    \label{l:sing-absorber-divisibility-by-s}
        Suppose that $r=(m+1)s$. Let $U_1,\dots, U_{m+2}$ be disjoint sets of vertices and let $M:=|U_1\cup\dots\cup U_{m+2}|$. Let $0<\epsilon \ll 1$. Suppose that $M\gg r$, $M \equiv 0 \pmod{r}$ and $|U_i|= \frac{M}{m+2}(1\pm \epsilon)$ for all $i\in [m+2]$. Consider the following type of move:
        \begin{itemize}
            \item $P_{j,k}$ move: for indices $j,k\in [m+2]$ remove $1$ vertex from $U_j$, $s-1$ vertices from $U_k$ and $s$ vertices from each $U_i$, $i\in[m+2]\setminus\{j,k\}$.
        \end{itemize}
        Using the above moves, we can get new sets $U_1'\subseteq U_1,\dots , U_{m+2}'\subseteq U_{m+2}$ such that each $|U_i'|$ is divisible by $s$ and $|U_1'\cup \dots \cup U_{m+2}'|\geq M-(m+2)s^2$.
    \end{lemma}

        The proof is essentially the same as the proof of Lemma~\ref{l:absorber-divisibility-by-r} so we omit it.

    \begin{lemma}\label{l:sing-absorber-partitioning}
        Suppose that $r=(m+1)s$. Let $U_1,\dots, U_{m+2}$ be disjoint sets of vertices and $M:=|U_1\cup\dots\cup U_{m+2}|$. Let $0<\epsilon \ll 1$. Suppose that $M\gg r$, $|U_i|\equiv 0 \pmod{s}$, $M\equiv 0 \pmod{r}$ and $|U_i|= \frac{M}{m+2}(1\pm \epsilon)$ for all $i\in [m+2]$. We can partition each set $U_i$ into $m+1$ parts $P_{i,1},\dots, P_{i,i-1}, P_{i,i+1},\dots, P_{i,m+2}$ so that the size of each part is divisible by $s$, and for each $j\in [m+2]$, the sizes $|P_{i,j}|$ are equal for all $i\neq j$.
    \end{lemma}

    \begin{proof}
          Set 
          \[x_i=\frac{1}{m+1}\sum_{j=1}^{m+2}|U_j|-|U_i| = \frac{M}{m+1}-|U_i|\] 
          and note that each $x_i$ is divisible by $s$. Now simply partition each $U_i$ into parts $P_{i,j}$ for $j\neq i$ so that $|P_{i,j}|=x_j$. Note that any partition satisfying the size requirements works.
    \end{proof}

    \subsection{From regular partition to clique factor}
        \label{sec:balancing}
        
    We are now ready to prove Theorem~\ref{t:non-extremal-conforming}.

   
    
    \begin{proof}[Proof of Theorem~\ref{t:non-extremal-conforming}]
    Consider the graph $\tilde G = G\setminus Z$ and note that since $|Z|=\zeta n$ and $\zeta \ll \delta \ll \gamma$, we have $\delta(\tilde G)\geq\left(1-\frac{s}{r} - 2\delta\right)|V(\tilde G)|$ and every subset of size $\frac{s}{r}|V(\tilde G)|$ induces in $\tilde G$ at least $\frac{\gamma}2 |V(\tilde G)|^2$ edges.
    Hence, we can apply to $\tilde G$ Lemma~\ref{l:partition-reduced-graph} with constants $\delta_{\ref{l:partition-reduced-graph}}=2\delta$ and $\gamma_{\ref{l:partition-reduced-graph}}=\frac{\gamma}2$. We thus obtain sets
    \begin{itemize}
        \item $S_{i,j}$, for $i\in[L_1]$, $j\in[m+1]$,
        \item $U_{k,\ell}$, for $k\in[L-L_1]$, $\ell\in[m+1+g]$,
        \item $X$,
    \end{itemize} 
    and a reduced graph $R$ corresponding to $(\eps,d^+)$-regular pairs. Let $\tilde X = X \cup Z$ and note that $|\tilde X|\leq \eps n$. Instead of sampling $\Gnp$ with $p = Cp_s(n)$ at once, we will expose three independent random graphs $G_i \sim G(n, p/4)$ for $i \in [3]$. The union $G_1 \cup G_2 \cup G_3$ can be viewed as a subgraph of $\Gnp$.
    
    The proof consists of four main steps.
        
    \smallskip
 
    \textit{Step 1. Redistributing $\tilde X$.} We redistribute the elements of $\tilde X$ among the sets $S_{i,j}$ and get new parts $S_{i,j}'\supseteq S_{i,j}$ which are of similar size as the original ones. Here we ensure that we maintain super-regularity between all sets in $\{S_{i,j}'\}_{j\in[m+1]}$, for $i\in[L_1]$.

    \smallskip

    \textit{Step 2. Readjusting the sizes of the sets $S_{i,j}'$.} Using appropriate absorbers (in sets $U_{k,\ell}$) and random edges of $G_1$, we take out some copies of $K_r$, each using at most one vertex from $\tilde X$, to return to exact appropriate size ratios. This yields new parts $S_{i,j}^\dagger \subseteq S_{i,j}'$ and $U_{k,\ell}^\dagger \subseteq U_{k,\ell}'$.

    \smallskip

    \textit{Step 3. Ensuring the divisibility of $|U_{k,1}^\dagger\cup\dots\cup U_{k,m+1+g}^\dagger|$.} We create a move $U-A-U'$, where $U$, $U'$ are parts of the regular partition and $A$ is an absorbing $(m+1+g)$-tuple. This move takes out a copy of $K_r$ which uses random edges in $G_2$ and in turn changes $|U|$ and $|U'|$ by $-1 $ and $+1\pmod{r}$, respectively. We can iterate it a constant number of times to ensure that the new parts $U_{k,\ell}^* \subseteq U_{k,\ell}^\dagger$ satisfy $|U_{k,1}^*\cup\dots\cup U_{k,m+1+g}^*| \equiv 0 \pmod{r}$, for all $k\in[L_2]$ $(L_2:=L-L_1)$.

    \smallskip

    \textit{Step 4. Covering the remainder using $G_3$.} In the last step we cover the remaining super-regular tuples of appropriate sizes using random edges in $G_3$.

    \smallskip

    We now explain each step in detail.

    \smallskip

    \textbf{Step 1.~Redistributing $\tilde X$.} 
    
    \smallskip
    
    Call an index $i \in [L_1]$ \textit{good for} $x$ if there is at most one index $j\in[m+1]$ such that $x$ has fewer than $d|S_{ij}|$ neighbours in $S_{ij}$. We claim that for each $x\in \tilde X$, there are at least ${dL}$ good indices. Indeed, assume that the opposite is true. Any index $i\in[L_1]$ that is not good for $x$ must yield at least two indices $j\in[m+1]$ for which $x$ has more than 
    \[(1-d)|S_{i,j}| = \begin{cases}
        (1-d)\frac sr \cdot \frac nL (1\pm \eps), & j\in[m] \\
        (1-d)\frac tr \cdot \frac nL (1\pm \eps), & j=m+1
    \end{cases}\] 
    non-neighbours in $S_{i,j}$. Then, since $L_1\geq (1-10r\delta_2) L$ and $\delta_2 \ll d$, the number of non-neighbours of $x$ is at least
		$$(1-10r\delta_2-d)L\cdot(1-d)\frac{s+t}{r}\cdot\frac{n}{L}(1-\eps)>\frac{ns}{r}.$$ 
	This contradicts the minimum-degree assumption on $x$.

    In order to assign good indices to vertices in $\tilde X$, we will use the following standard argument which will also be used later in the proof.
    \begin{claim}\label{c:even-distribute-X}
        There is an assignment $f: \tilde X \to L_1$ such that for each $x\in \tilde X$, $f(x)$ is good for $x$, and for each $i \in [L_1]$, $|f^{-1}(i)| \leq \frac{\eps n}{dL}$.
    \end{claim}
    
    \begin{proof}
        Consider a random assignment which for each vertex $x\in \tilde X$ chooses a good index $i\in[L_1]$ uniformly at random and sets $i = f(x)$. Now fix an index $i$. Since $\pr{f(x) = i} \leq (dL)^{-1}$, the random variable $|f^{-1}(i)|$ is stochastically dominated by $\Bin(\eps n/2, (dL)^{-1})$, and hence by Chernoff
        $$\pr{|f^{-1}(i)| \geq \frac{\eps}{d}\frac{n}{L}} = e^{-\Omega(n)}.$$ 
        Taking the union bound over all $i\in[L_1]$ implies that $f$ satisfies the conclusion of the claim with positive probability.
    \end{proof}

    Now we will redistribute the vertices of $\tilde X$, creating a new partition that satisfies properties similar to \ref{l:part-w-r-f_(i)}--\ref{l:part-w-r-f_(iv)}. For this purpose, we will use Lemma~\ref{l:slicing-adding}. 

    \begin{claim}
        Let $\eps_1= \sqrt \eps / d$. There is a vertex partition of $G$ into sets $S'_{i,j}$ and $U'_{k,l}$ (with the same indexing as before) with the following properties. Let $R'$ be the graph whose vertices are sets $S'_{i,j}$, $U'_{k,l}$ and two vertices $A$ and $B$ are connected if $(A,B)$ is an $\left(\eps_1,\left(d/2\right)^+\right)$-regular pair.
        \begin{enumerate}[label=(\roman*')]
            \item\label{claim_(i)} For each $i \in [L_1]$ and $j < j' \in [m+1]$, the edges $S'_{i,j}S'_{i,j'}$ are in $R'$. Similarly, for each $k\in [L_2]$ and $\ell < \ell' \in [m+1+g]$, the edges $U'_{k,\ell}U'_{k,\ell'}$ are in $R'$.
            
            \item\label{claim_(ii)} For each $i \in [L_1]$, $j\in[m]$ we have 
            $$|S'_{i,j}| = \frac{s}{r}\cdot\frac{n}{L}(1\pm \eps_1),\quad |S'_{i,m+1}| = \frac{t}{r}\cdot\frac{n}{L}(1\pm \eps_1)$$
            and $(S'_{i,1},\dots, S'_{i,m+1})$ is an $\left(\eps_1,\left(d/2\right)^+,d/2\right)$-super-regular tuple.
            
            \noindent
            For each $k \in [L_2]$, $\ell\in[m+1+g]$ we have 
            $$|U'_{k,\ell}|=\frac {1}{m+1+g}\cdot\frac{n}{L}(1 \pm \eps_1)$$
            and $(U'_{k,1},\dots, U'_{k,m+1+g})$ is an $\left(\eps_1,\left(d/2\right)^+,d/2\right)$-super-regular tuple.
            
            \item\label{claim_(iii)}  For $A, B \in V(R')$, there are at least $d L$ indices $k \in [L_2]$ such that both $A$ and $B$ have at least $m-1+g$ $R'$-neighbours in $\{U'_{k,1}, U'_{k,2}, \ldots, U'_{k,m+1+g}\}$.
            
            \item\label{claim_(iv)} Denote $s_{ij}:=|S_{i,j}'|-|S_{i,j}|$. Then, $$\sum_{i,j}s_{ij}\leq \eps n.$$ 
            
            \item\label{claim_(v)} For all $i\in[L_1]$, $j\in [m+1]$, 
            $$|\tilde X\cap S_{i,j}'|\leq \frac{\eps}{d}|S_{i,j}'|.$$
        \end{enumerate}
    \end{claim}
             
    \begin{proof}
        Recall that 
        \[ |S_{i,j}|=\frac{s}{r}\cdot\frac{n}{L}(1\pm\eps), j\in[m] \quad \text{ and } \quad |S_{i,m+1}|=\frac{t}{r}\cdot\frac{n}{L}(1\pm\eps). \]
        Similarly,
        \[ |U_{k,\ell}|=\frac1{m+1+g}\cdot\frac{n}{L}(1\pm\eps).\]
        Using Claim~\ref{c:even-distribute-X}, we can redistribute the vertices of $\tilde X$ into $S_{i,j}$ to obtain new sets $S_{i,j}'$ as follows. By the definition of good indices, any $x\in \tilde X$ can be added to some part $S_{f(x),j}$, so that it has at least a $d$ proportion of neighbours in each part $S_{f(x),j'}$ for $j' \neq j$. Now taking $\eps_1 = \sqrt{\eps}/d$, Claim~\ref{c:even-distribute-X} asserts that no more than 
        \[\frac{\eps n}{dL} < \eps_1\sqrt{\eps}|S_{i,j}|< \frac{\eps_1}2|S_{i,j}|\] vertices were added to each $S_{i,j}$. Moreover, set $U'_{k,\ell}=U_{k,\ell}$ for all $k\in[L-L_1]$, $\ell\in[m+1+g]$.
        Properties \ref{claim_(ii)} and \ref{claim_(v)} follow immediately from the above observations, where additionally we use the fact that vertices in $\tilde X$ are redistributed in such a way, that the super-regularity is preserved between suitable pairs with parameter $d/2$. 
        
        Next, property \ref{claim_(iii)} follows from property \ref{l:part-w-r-f_(iii)} using Lemma~\ref{l:slicing-adding}, i.e.~adding an $\eps/d$ proportion of vertices to any of the sets $S_{i,j}$ turns $(\sqrt{\eps},d)$-regular pairs into $(\sqrt\eps/d, d')$-regular pairs with $d'\geq d-3\sqrt \eps$. Similarly property \ref{claim_(i)} follows from property \ref{l:part-w-r-f_(i)}. Finally, the sum in \ref{claim_(iv)} is just $|\tilde X|\leq \eps n $.
    \end{proof} 

    \smallskip

    \textbf{Step 2.~Readjusting the sizes of the sets $S_{i,j}'$.}  
    
    \smallskip
    
    In this step, we will remove a small number of vertices from the sets $S_{i,j}'$ and $U'_{k,\ell}$, forming disjoint copies of $K_r$, and obtain a new collection of sets $S_{i,j}^\dagger\subseteq S_{i,j}'$ and $U_{k,\ell}^\dagger\subseteq U_{k,\ell}'$ satisfying
    \begin{equation}\label{eq:correct_proprtion_in_S_dagger}
        |S_{i,1}^\dagger|= |S_{i,2}^\dagger |= \dots = |S_{i,m}^\dagger| = \frac {s}{t} \cdot |S_{i,m+1}^\dagger|
    \end{equation}
    for all $i\in[L_1]$. To this end we define the following procedure. For $i\in [L_1]$, $j\in[m+1]$ and $k\in [L_2]$, in a \textit{basic $(i,j;k)$-move} we remove a copy of $K_r$ in $G \cup G_1$ which contains one vertex from $S_{i,j}'$ and $r-1$ vertices from $U_{k,1}'\cup \dots \cup U_{k,m+1+g}'$. We will slightly abuse the notation by referring to this move as a basic $(i,j;k)$-move even after some vertices have been already removed from the corresponding vertex sets $S_{i,j}'$ or $ U_{k,\ell}'$. 

    Recall that $G_1\sim G(n, Cp_s)$. Before we proceed, we define the following property of the graph $G_1$, which, as we show below, holds w.h.p. Let $d_1$ be a constant with $\eps \ll d_1 \ll d$.
    \begin{description}
	\item[(P1)] For $i=0,1,\ldots,m+3$, let $A_i\in V(R')$ and let $B_i \subset A_i$ with $|B_i| \geq |A_i|/2$. Assume that all pairs $(A_i, A_j)$ are $(\eps^{1/10}, d')$-regular for some $d' \geq d_1$, with (possibly) the exception of $(A_0, A_{1})$ and $(A_0, A_{2})$. Then $G \cup G_1$ contains a clique with exactly one vertex in $B_0$, $s-1$ vertices in each $B_{1}$ and $B_{2}$, and $s$ vertices in $B_i$, $i = 3, \dots, m+3$.
    \end{description}
			
    \begin{claim}\label{c:p1p2}
	$G_1$ satisfies (P1) with probability $1-e^{-\omega(n)}$.
    \end{claim}
            
      
    \begin{proof}
        Let $H$ be a graph with vertex set $U_0\cup U_1 \cup\dots\cup U_{m+3}$, where $|U_0|=1$, $|U_1|=|U_2|=s-1$, and $|U_3|=\dots=|U_{m+3}|=s$, where all pairs apart from $(U_0,U_1)$ and $(U_0, U_2)$ span a complete bipartite graph.
        Let $\mathcal{H}$ be the family of copies of $H$ in $G$ with the part corresponding to $U_i$ contained in $B_i$, $i=0,\dots,m+3$. Using the Counting Lemma we have $|\mathcal{H}|\geq \nu n^{|V(H)|}$, for a specific constant $\nu$. 
        We claim that after revealing the edges of $G_1$, w.h.p.~at least one copy of $H\in \mathcal{H}$ becomes the desired clique. Indeed, note that the complement $\bar{H}$ of $H$ consists of $m+3$ copies of $K_s$, with exactly two copies sharing exactly one vertex and the remaining copies being vertex disjoint. Then by Lemma~\ref{l:many_ks_and_one_ks-ks}, the failure probability that no element in $\mathcal{H}$ spans a~complete graph in $G\cup G_1$ is at most $e^{-\Theta(n^2p)}=e^{-\omega(n)}$. Taking the union bound over all $2^{O(n)}$ choices of sets $B_i$ and $O(1)$ choices of $A_i\in V(R')$, we get that (P1) holds with probability $1-e^{-\omega(n)}$.
    \end{proof}

    Henceforth we assume that $G_1$ indeed satisfies (P1). Notice that (P1) implies the following.
    
    \begin{remark}\label{rem:(P1)}
    If $g=0$ 
    (the regular case $r = ms+t$ with $t < s$), 
    then the graph $G \cup G_1$ contains:
    \begin{itemize}
        \item a~$K_{r}$-copy in $A_0 \cup A_1 \cup \dots \cup A_{m+1}$ with exactly one vertex in $A_0$,
        \item a~$K_{r}$-copy in $A_1 \cup A_2 \cup \dots \cup A_{m+1}$ with exactly $s-1$ vertices in $A_1$ and $t+1\leq s$ vertices in $A_{m+1}$,
        \item a~$K_{r+1}$-copy in $A_3 \cup A_4 \cup \dots \cup A_{m+3}$ with exactly $t+1$ vertices in $A_{m+3}$ and $s$ vertices in the remaining sets.
    \end{itemize} 
        
    If $g=1$ 
    (the singular case $r = ms+t$ with $t = s$), 
    then the graph $G \cup G_1$ contains:
    \begin{itemize}
        \item a $K_{r}$-copy in $A_0 \cup A_2 \cup A_3 \cup \dots \cup  A_{m+2}$ with exactly one vertex in $A_0$,
        \item a~$K_{r+1}$-copy in $A_0\cup A_3 \cup A_4 \cup \dots \cup A_{m+3}$ with exactly $1$ vertex in $A_0$ and $s$ vertices in the remaining sets.
    \end{itemize} 
    \end{remark}
    
    We will now use this property to perform a series of $(i,j;k)$-moves which will lead us to a~partition satisfying \eqref{eq:correct_proprtion_in_S_dagger}. 
        
    \begin{claim}
        It is possible to perform an $(i,j;k)$-move as long as $S_{i,j}'$ forms a $(2\eps_1, d')$-regular pair, for some $d'\geq d/4$, with at least $m-1+g$ sets $U_{k,\ell}'$.
    \end{claim} 
        
    \begin{proof}
        If $g=0$ (regular case), w.l.o.g.~assume that $S_{i,j}'$ is $(2\eps_1,d')$-regular with the sets $U_{k, 3}',\dots,U_{k,m+1}'$. We apply (P1) with $A_0=S_{i,j}', A_1=U_{k,1}', \ldots, A_{m+1}=U_{k,m+1}'$ to remove a copy of $K_r$ with $t\leq s-1$ vertices in $A_1$ and $s-1$ vertices in $A_2$. If $g=1$, we similarly apply (P1) with $A_0=S_{i,j}', A_2=U_{k,1}', A_3=U_{k,2}', \dots, A_{m+2}=U_{k,m+1}'$.
    \end{proof}

    In order to be able to perform a series of $(i,j;k)$-moves, we need to show that while removing a small number of vertices in the copies of $K_r$, regularity is maintained. Note that by removing at most an $\eps_1$ proportion of vertices from any partition set, using Lemma~\ref{l:slicing}, we can ensure that pairs which were $(\eps_1,(d/2)^+)$-regular are $(2\eps_1,(d/4)^+)$-regular in the new partition, and, analogously, pairs that were $(\eps_1,(d/2)^+,d/2)$-super-regular become $(2\eps_1,(d/4)^+,d/4)$-super-regular. Hence, we need to perform the $(i,j;k)$-moves in such a way that each partition set loses at most an $\eps_1$ fraction of its vertices.

    Recall that $s_{ij} = |S_{i,j}'|-|S_{i,j}| \leq \frac{\eps}{d}|S_{i,j}'| \leq \eps_1|S_{i,j}|/2$. For each part $S_{i,j}'$, we will remove $s_{ij}$ vertices by performing a sequence of basic $(i,j;k)$-moves. However, we don't want to use any of the vertices $k\in[L_2]$ too often, so that the total number of vertices removed from any set $U_{k,\ell}'$ is also at most an $\eps_1$ fraction. Recall that by \ref{claim_(iii)}, for each $S_{i,j}'$, there are at least $dL$ indices $k$ for which an $(i,j;k)$-move can be performed. Hence, by the same argument as in Claim~\ref{c:even-distribute-X}, and using \ref{claim_(iv)}, we can assign $s_{ij}$ indices $k\in[L_2]$, possibly with repetitions, to each $S_{i,j}'$ so that none of the indices is assigned more than $\frac{2\eps n}{dL}$ times. Now for each assigned pair $S_{i,j}'$ and $k$, we apply a basic $(i,j;k)$ move. This removes exactly $s_{ij}\leq \eps_1|S_{i,j}|/2$ vertices from each set $S_{i,j}'$ and at most an $4r\eps/d < \eps_1$ proportion of vertices from any given $U_{k,\ell}'$.
    
    Denote the resulting sets by $S_{i,j}^\dagger$ and $U_{i,j}^\dagger$. Then \eqref{eq:correct_proprtion_in_S_dagger} follows by \ref{l:part-w-r-f_(ii)} and the fact that for each $i\in[L_1]$ and $j\in[m+1]$ we have $|S_{i,j}^\dagger|=|S_{i,j}|$. Note also that each of the removed copies of $K_r$ contains at most one vertex from $\tilde X$.
    
    Denote $\Ub_k^\dagger := U_{k,1}^\dagger \cup \dots \cup U_{k,m+1+g}^\dagger$, for $k\in[L_2]$. Let $R^\dagger$ be the graph with vertex set $U_{k,\ell}^\dagger$ and edges corresponding to $(2\eps_1,(d/4)^+)$-regular pairs. Then by our construction and \ref{claim_(i)}, \ref{claim_(ii)} and \ref{claim_(iii)}, we have
    \begin{enumerate}[label=(\roman*$^\dagger$)]
        \item\label{(i)dagger} for each $k\in [L_2]$ and $\ell < \ell' \in [m+1+g]$, the edges $U^\dagger_{k,\ell}U^\dagger_{k,\ell'}$ are in $R^\dagger$;
        \item\label{(ii)dagger} for each $k \in [L_2]$, $\ell\in[m+1+g]$ we have 
            $$|U^\dagger_{k,\ell}|=\frac {1}{m+1+g}\cdot\frac{n}{L}(1 \pm 2\eps_1)$$
            and $(U^\dagger_{k,1},\dots, U^\dagger_{k,m+1+g})$ is a $\left(2\eps_1,\left(d/4\right)^+,d/4\right)$-super-regular tuple;
        \item\label{(iii)dagger} for any $A,B\in V(R^\dagger)$, there are at least $dL$ indices $k\in[L_2]$ such that both $A$ and $B$ have at least $m-1+g$ $R^\dagger$-neighbours in $\{U_{k,1}^\dagger, U_{k,2}^\dagger,\dots, U_{k,m+1+g}^\dagger\}$.
    \end{enumerate}
    

    \smallskip

    \textbf{Step 3. Ensuring the divisibility of $|\Ub_k^\dagger|$.} 

    \smallskip

    In this step, we will reveal the edges of $G_2$. For any $k,q\in[L_2]$, $\ell\in[m+1+g]$, we can define a basic $(k,\ell;q)$-move similarly as in Step 2, that is, in such a move we remove a~copy of $K_r$ in $G \cup G_2$ that contains one vertex from $U_{k,\ell}^\dagger$ and $r-1$ vertices from $\Ub_{q}^\dagger$.
    Now fix a pair $(\Ub_k^\dagger, \Ub_{k'}^\dagger)$. We define a \emph{$(k,k')$-transfer} as follows. Choose an auxiliary index $q$ for which $U_{k,1}$ and $U_{k',1}$ have at least $m-1+g$ $R^\dagger$-neighbours in $\Ub_{q}^\dagger$. Perform one basic $(k,1;q)$-move and $r-1$ basic $(k',1;q)$-moves. Modulo $r$, this changes $|\Ub_k^\dagger|$ by -1, $|\Ub_{k'}^\dagger|$ by $+1$, and preserves $|\Ub_{q}^\dagger|$. Note that by \ref{(iii)dagger}, for any pair $k, k' \in [L_2]$, we can always find an index $q$ such that a $(k,k')$-transfer can be performed using $q$ as an auxiliary index. Using the same method as in the proof of Lemma~\ref{l:absorber-divisibility-by-r}, using no more than $rL_2$ $(k,k')$-transfers, we can pass to parts $U_{k,\ell}^\ddagger\subseteq U_{k,\ell}^\dagger$ with $\Ub_k^\ddagger := U_{k,1}^\ddagger \cup \dots \cup U_{k,m+1+g}^\ddagger$ satisfying $|\Ub_k^\ddagger| \equiv 0 \pmod{r}$. Note that the number of removed vertices in this step is at most $r^3L_2$ and none of them is a vertex from $\tilde X$.

    \smallskip
    
    \textbf{Step 4. Covering the remainder using $G_3$.}

    \smallskip

    We will reveal the graph $G_3$ separately for sets $S_{i,j}^\dagger$ and $U_{k,\ell}^\ddagger$.
    Consider first any $i \in [L_1]$. Recall that $(S_{i,1}^\dagger,S_{i,2}^\dagger,\dots, S_{i,m+1}^\dagger)$ is an $(\eps_1,(d/4)^+,d/4)$-super-regular $(m+1)$-tuple satisfying~\eqref{eq:correct_proprtion_in_S_dagger}. Thus, by Lemma~\ref{l:st-reg-pairs-factors}, w.h.p.~there is an $\tilde X$-conforming collection $\cK_i$ of vertex-disjoint $K_r$-cliques in $G \cup G_3$ covering  $(S_{i,1}^\dagger,S_{i,2}^\dagger,\dots, S_{i,m+1}^\dagger)$. Since the total number of such $(m+1)$-tuples is constant, w.h.p.~we can find such a collection $\cK_i$ for any $i\in[L_1]$. 
    
    Now consider $k\in[L_2]$. Since in Step 3 we removed at most $r^3L$ vertices from any set $U_{k,\ell}^\dagger$ to obtain sets $U_{k,\ell}^\ddagger$, the latter satisfy conditions analogous to \ref{(i)dagger}--\ref{(iii)dagger} with slightly perturbed parameters. In particular we can assume that $|U_{k,l}^\ddagger|=\frac{1}{m+1+g}\cdot\frac{n}{L}(1\pm3\eps_1)$ for each $k\in[L_2]$, $\ell\in[m+1+g]$ and that regularity, as well as super-regularity, is maintained.

    \smallskip

    \textit{The regular case ($g=0$).}

    \smallskip
       
    We first use Lemma~\ref{l:absorber-divisibility-by-r} and property (P1) to pass to an $(m+1)$-tuple $\{U_{k,1}^*, \ldots, U_{k,m+1}^*\}$ in which each part is divisible by $r$. To be more precise, Lemma~\ref{l:absorber-divisibility-by-r} gives a sequence of moves that needs to be performed to obtain the divisibility by $r$, while Remark~\ref{rem:(P1)} tells us that we can find a copy of $K_r$ in $G\cup G_3$ corresponding to a specific move. At this point we have removed at most a constant number of vertices from each part.
    
    Take $x$ to be the maximal integer such that all the parts have at least $xr$ vertices, in particular $x\geq \left\lfloor\frac1r\cdot\frac{1}{m+1}\cdot\frac{n}{L}(1-4\eps_1)\right\rfloor$. Let $|U_{k,\ell}^*| = xr+y_\ell$, noting that all the $y_\ell$ are divisible by $r$. We have 
    \[y_\ell\leq |U_{k,\ell}^*|-\frac{1}{m+1}\cdot\frac{n}{L}(1-4\eps_1)\leq 9\eps_1 xr.\] 
    Next, we partition each set $U_{k,\ell}^*$, $\ell\in[m+1]$, randomly into $m$ parts $U_{k,\ell,i}^*$, $i\in[m+1]\setminus\{\ell\}$, of size $xs$, and one part $U_{k,\ell,\ell}^*$ of size $xt+y_\ell$, so that by Lemma~\ref{l:slicing} all pairs $(U_{k,\ell,i}^*,U_{k,\ell',j}^*)$, $\ell\neq\ell'$, $i,j\in[m+1]$, are $(4r\eps_1, d')$-regular with $d' \geq d/(8r)$. Moreover, using an argument similar to that in the proof of Lemma~\ref{l:st-reg-pairs-factors}, w.h.p.~these pairs are $(4r\eps_1, d', d/(8r))$-super-regular with some $d' \geq d/(8r)$.
    
    Now consider an $(m+1)$-tuple $\{U^*_{k,1,\ell},U^*_{k,2,\ell},\ldots,U^*_{k,m+1,\ell}\}$. Using Lemma \ref{l:equalize-(m+1)-tuples-to-s/t} and Remark~\ref{rem:(P1)}, by removing an additional $4r\eps_1|\Ub_k^*|$ copies of $K_r$ in $G \cup G_3$, we can pass to an $(m+1)$-tuple $\{U_{k,1,\ell},U_{k,2,\ell},\ldots,U_{k,m+1,\ell}\}$ which satisfies the assumptions of Lemma~\ref{l:st-reg-pairs-factors}. Therefore, by Lemma~\ref{l:st-reg-pairs-factors}, w.h.p.~for each $k\in[L_2], \ell\in[m+1]$ there is a~collection of vertex-disjoint $K_r$-copies in $G \cup G_3$ covering the $(m+1)$-tuple $\{U_{k,1,\ell},U_{k,2,\ell},\ldots,U_{k,m+1,\ell}\}$.

    \smallskip

    \textit{The singular case $(g=1)$}

    \smallskip

    The proof is analogous to that in case $g=0$. First, using Lemma~\ref{l:sing-absorber-divisibility-by-s} and Remark~\ref{rem:(P1)} we can pass to an $(m+2)$-tuple $\{U_{k,1}^*, \ldots, U_{k,m+2}^*\}$ in which each part is divisible by~$s$. Next, using Lemma~\ref{l:sing-absorber-partitioning}, we can randomly partition each set $U_{k,\ell}^*$ into $(m+1)$-tuple $\{U_{k,\ell,1},\dots,U_{k,\ell,\ell-1},U_{k,\ell,\ell+1},\dots,U_{k,\ell,m+2}\}$, so that for each $\ell\in[m+2]$, all sets in the $(m+1)$-tuple $\{U_{k,1,\ell},\dots,U_{k,\ell-1,\ell},U_{k,\ell+1,\ell},\dots,U_{k,m+2,\ell}\}$ have the same size and all the involved pairs are $(4(m+1)\eps_1,d',d/(8r))$-super-regular with some $d'\geq d/(8r)$. By Lemma~\ref{l:st-reg-pairs-factors}, w.h.p.~there is a partial $K_r$-factor in $G\cup G_3$ covering these $(m+1)$-tuples. 
    \end{proof}

    We finish this section by providing a proof of Lemma~\ref{l:annoying-non-extremal}, which grants at least $\eta n$ vertex-disjoint copies of $K_{r+1}$ in $G \cup \Gnp$ under suitable assumptions. For simplicity, we deduce the lemma from Lemma~\ref{l:partition-reduced-graph}, but the only important ingredient is that in the singular case ($r = (m+1)s$), $G$ contains a regular $(m+2)$-tuple. There are several other ways to prove the lemma, and the constant $\eta$ is far from optimal.


    \begin{proof}[Proof of Lemma~\ref{l:annoying-non-extremal}]
    Let $G$ be the graph satisfying the assumption of the Lemma, with the given constants $r, s, \gamma$ and $\delta$. Suppose also that 
         $$\eta \ll \eps \ll d \ll \delta =\delta_2 \ll \gamma, r^{-1}.$$ 
        Apply Lemma~\ref{l:partition-reduced-graph} to $G$; specifically, $G$ contains an $(\epsilon, d)$-regular $(m+1+g)$-tuple $\{U_{k,1},\dots,U_{k,m+1+g}\}$. Recall that $g=1$ if $r = (m+1)s$, and $g=0$ otherwise. Apply Remark~\ref{rem:(P1)} (that is, property (P1)) $\eta n$ many times to find w.h.p.~the desired collection of vertex-disjoint $K_{r+1}$-copies in $G\cup \Gnp$.
    \end{proof}
    
%************************* The partitioning lemma ****************************
    
    \section{The partitioning lemma}
	   \label{sec:partition-lemma} 
    
    In this section, we prove the partitioning lemma.
    
    \begin{proof}[Proof of Lemma~\ref{l:4.3chr}] 
    Let $\beta_0 = \frac{1}{10m(m+1)^2}$, and given $\beta\leq \beta_0$ take $\gamma_m \ll \beta$. 
    In particular, we need the following dependencies
    \begin{align*}
        \gamma_m \leq \frac{\beta^2}2, \quad \gamma_m \leq \sqrt{\frac{\beta^2(s/r-3\beta)}{2}}
    \end{align*}
    Let a graph $G$ and a sequence $\gamma_1, \dots, \gamma_m$ as in the statement of the Lemma be given.

    Let $h\in[m]$ be a maximal integer for which there exists a partition
    \begin{equation}\label{eq:partition} 
        V(G) = A_1^* \dcup \dots \dcup A_h^* \dcup A_{h+1}
    \end{equation}
    satisfying the following properties
    \begin{enumerate}[label=(\roman**)]
        \item\label{(i*)} For $i\in[h]$, $|A_i|=(s/r \pm \beta\gamma_i)n$ and $e(A_i^*)\leq \beta^2\gamma_i n^2$.
        \item\label{(ii*)}  Every vertex $x \in V(G) \setminus A_{h+1}$ has at most $4 \beta n$ non-neighbours in $A_{h+1}$.
        \item\label{(iii*)}  Every vertex $x \in A_{h+1}$ has at least $2\beta n$ neighbours in $A_i^*$, for any $i \in [h]$.
    \end{enumerate}
    Note that such $h$ exists, as $h=0$ trivially satisfies the above properties. Properties \ref{(i*)}--\ref{(iii*)} will help us to establish properties \ref{partition(i)}--\ref{partition(iii)} of the lemma. To get the remaining ones, we first deal with property \ref{partition(iv)}.

    \begin{claim}\label{claim:4.4chr}
        If $h<m$ and $X \subset A_{h+1}$ with $|X| =sn/r$, then $e(X) > \gamma_{h+1}^2 n^2$.
    \end{claim}
    \begin{proof}
        Assume for a contradiction that $h<m$ and there exists a subset $X \subseteq A_{h+1}$ with $|X| =sn/r$ and $e(X) \leq \gamma_{h+1}^2 n^2$. Note that 
        \begin{equation}\label{eq:size_of_X}
            |X| \geq n - \delta(G).
        \end{equation}
        We will show that we can change the initial partition (\ref{eq:partition}), keeping the conditions \ref{(i*)}--\ref{(iii*)}, but increasing $h$.

        Set $Y = A_{h+1}\setminus X$ and let
        \[ R = \{x\in X: |Y\setminus N(x)| \geq 3\beta n\}. \]
        Note that by (\ref{eq:size_of_X}) and $e(X) \leq \gamma_{h+1}^2n^2$, the number of missing edges between $X$ and $V(G)\setminus X$ is at most $2\gamma_{h+1}^2 n^2$.
        Moreover, since every vertex in $R$ contributes at least $3\beta n$ missing edges between $X$ and $V(G)\setminus X$, we have
        \[ |R| \leq \frac{2\gamma_{h+1}^2}{3\beta}n \leq \beta\gamma_{h+1}n.\]
        Similarly, we let
        \[ S = \{y\in Y: |N(y)\cap X| \leq 3\beta n\}, \]
        and note that
        \[ |S| \leq \frac{2\gamma_{h+1}^2}{s/r-3\beta}n\leq \frac{\beta^2\gamma_{h+1}}2 n.  \]
        Next, set
        \[ A_{h+1}^* = (X\setminus R) \cup S \quad \text{and} \quad A_{h+2} = (Y\setminus S)\cup R,\]
        and consider the partition
        \[ V(G) = A_1^* \dcup \dots \dcup A_{h+1}^* \dcup A_{h+2}. \]
        We claim that this new partition also satisfies \ref{(i*)}--\ref{(iii*)}, thus contradicting the maximality of $h$. 

        To show \ref{(i*)}, we only need to consider $i=h+1$. Since both sets $S$ and $R$ have size at most $\beta\gamma_{h+1}n$, we have $|A_{h+1}^*| \leq (s/r \pm \beta\gamma_{h+1})n$. Moreover,
        \[ e(A_{h+1}^*) \leq e(X) + |S|n \leq \left(\gamma_{h+1}^2 + \frac{\beta^2\gamma_{h+1}}2\right)n^2 \leq \beta^2\gamma_{h+1}n^2.\]

        To show \ref{(ii*)}, consider a vertex $x \in V(G)\setminus A_{h+2}$. If $x\notin A_{h+1}^*$, then since the initial partition (\ref{eq:partition}) satisfied \ref{(ii*)}, we know that $x$ has at most $4\beta n$ non-neighbours in $A_{h+1}$, thus also at most this many non-neighbours in $A_{h+2} \subset A_{h+1}$. Now suppose $x\in A_{h+1}^* = (X\setminus R)\cup S$. Due to $|A_{h+2}\setminus N(x)| \leq |Y \setminus N(x)| + |R|$ and $|R|\leq \beta n$, it suffices to show that $|Y\setminus N(x)| \leq 3\beta n$. If $x\in X\setminus R$ this is clear by the definition of $R$, so suppose that $x\in S$. Using (\ref{eq:size_of_X}), we get that $|X| \geq |V(G) \setminus N(x)| \geq |X \setminus N(x)| + |Y \setminus N(x)|$, hence, by the definition of $S$, $|Y\setminus N(x)| \leq |X\cap N(x)| \leq 3\beta n$.

        Finally, to show \ref{(iii*)}, let $x\in A_{h+2} = (Y\setminus S)\cup R$ and note that it is enough to consider $i=h+1$. Since $|A_{h+1}^* \cap N(x)| \geq |X \cap N(x)| - |R|$ and $|R| \leq \beta n$, it suffices to show that $|X \cap N(x)| \geq 3\beta n$. If $x\in Y\setminus S$ this is clear by the definition of $S$, so suppose that $x \in R$. Then, using again (\ref{eq:size_of_X}) and the definition of $R$, we have 
        \begin{align*} |X \cap N(x)| & = |X| + |N(x)| - |X \cup N(x)| \geq |V(G)| - |X \cup N(x)| \\
        & \geq |Y \setminus N(x)| \geq 3\beta n.
        \end{align*}
    \end{proof}

    We now adjust the partition to also get the remaining properties \ref{partition(v)} and \ref{partition(vi)}. To this end, choose $t\geq 0$ and a partition
    \[ V(G) \setminus A_{h+1} = A_1 \dcup \dots \dcup A_h\]
    such that
    \begin{enumerate}
        \item[(a)] $\sum_{i=1}^h |A_i^* \bigtriangleup A_i| \leq 2t$ 
        \item[(b)] $\sum_{i=1}^h e(A_i) \leq \sum_{i=1}^h e(A_i^*) - \beta tn$,
    \end{enumerate}
    and, subject to the above conditions, $t$ is maximal. Note that we can always find such a~partition and a parameter $t$, as $t=0$ and $A_i = A_i^*$ satisfies conditions (a) and (b). Moreover, conditions (b) and \ref{(i*)} imply that
    \[ 0 \leq m\beta^2\gamma_h n^2 - \beta t n,\]
    hence
    \begin{equation}\label{eq:t} 
        t \leq m\beta\gamma_h n \leq \frac{\beta}2 n.
    \end{equation}
    We claim that the partition
    \[ V(G) = A_1 \dcup \dots \dcup A_{h+1}\]
    has all desired properties \ref{partition(i)}--\ref{partition(vi)} of the lemma.

    To show \ref{partition(i)}, let $i\in[h]$. Then, using (a), \ref{(i*)} and (\ref{eq:t}), we get
    \begin{align}\label{eq:size_of_A_i}
        |A_i| & = |A_i^*| \pm 2t = \left(\frac{s}{r} \pm \beta\gamma_i\right)n \pm 2m\beta \gamma_hn \nonumber \\
        & = \left(\frac{s}{r} + 3\beta m \gamma_h\right) n = \left(\frac{s}{r} \pm \gamma_h\right)n.
    \end{align}

    Property \ref{partition(ii)} is equivalent to property \ref{(ii*)}.

    To show \ref{partition(iii)}, note that $|A_i^*\setminus A_i| \leq 2t \leq \beta n$, hence due to \ref{(iii*)}, every vertex $x\in A_{h+1}$ has at least $\beta n$ neighbours in any $A_i$.

    Property \ref{partition(iv)} is given by Claim~\ref{claim:4.4chr}.

    To show \ref{partition(v)}, take distinct indices $i,j\in[h]$. If there exists a vertex $x\in A_i$ with $|N(x) \cap A_j| < \frac13|A_j|$, then we can replace sets $A_i, A_j$ by sets $A_i'=A_i\setminus\{x\}, A_j'=A_j\cup\{x\}$, and reach a contradiction to the maximality of $t$. Indeed, for $k\in\{i, j\}$ we have $|A_k^* \bigtriangleup A_k'| \leq |A_k^* \bigtriangleup A_k| + 1$, and if $|N(x)\cap A_j| < \frac13|A_j|$, then $|N(x)\cap A_i| > \frac13|A_i|$, which in turn implies that $e(A_i')+e(A_j') < e(A_i)+e(A_j)$, so we can increase $t$.

    Finally, to show \ref{partition(vi)}, let $i,j\in[h+1]$ be distinct indices with $i < j$. Note that by \ref{(i*)}, (a) and (\ref{eq:t}) we have
    \begin{align*}
        e(A_i) \leq e(A_i^*) + |A_i\setminus A_i^*|n \leq (\beta^2\gamma_h+2\beta m\gamma_h)n^2 \leq 3\beta m \gamma_h n^2.
    \end{align*}
    Let $M$ be the number of missing edges between $A_i$ and $A_j$. Then
    \[ |A_i|\left(1-\frac{s}{r}\right)n \leq \sum_{x\in A_i} d(x) = 2e(A_i) + |A_i|(n-|A_i|) - M,\]
    hence, using (\ref{eq:size_of_A_i}), we can upper bound $M$ by
    \begin{align*}
        M & \leq |A_i| \left( \frac{s}{r}n - |A_i|\right) + 6\beta m\gamma_h n^2 \leq 3\beta m \gamma_h n (|A_i| + 2n) \leq \gamma_h|A_i||A_j|.
    \end{align*}
    
        
    \end{proof}
	
    %\bibliographystyle{amsplain}	
    %\bibliography{main}
  %\end{document}

  
    
	\bibliographystyle{plain}
	\begin{thebibliography}{50}

\bibitem{abcd22}
Peter Allen, Julia B\"ottcher, Jan Corsten, Ewan Davies, Matthew Jenssen, Patrick Morris, Barnaby Roberts, and Jozef Skokan, \emph{A robust {C}orr\'adi-{H}ajnal theorem}, Random Structures Algorithms \textbf{65} (2024), no.~1, 61--130. \MR{4760768}

\bibitem{aes74}
B{\'e}la Andr{\'a}sfai, Paul Erd{\"o}s, and Vera~T S{\'o}s, \emph{On the connection between chromatic number, maximal clique and minimal degree of a graph}, Discrete Mathematics \textbf{8} (1974), no.~3, 205--218.

\bibitem{adrrs21}
S.~Antoniuk, A.~Dudek, Chr. Reiher, A.~Ruci\'{n}ski, and M.~Schacht, \emph{High powers of {H}amiltonian cycles in randomly augmented graphs}, J. Graph Theory \textbf{98} (2021), no.~2, 255--284. \MR{4371440}

\bibitem{akr24}
Sylwia Antoniuk, Nina Kam{\v{c}}ev, and Christian Reiher, \emph{Clique factors in randomly perturbed graphs: the transition points}, arXiv preprint arXiv:2410.11003 (2024).

\bibitem{balog16}
J.~Balogh, T.~Molla, and M.~Sharifzadeh, \emph{Triangle factors of graphs without large independent sets and of weighted graphs}, Random Struct. Algorithms \textbf{49} (2016), no.~4, 669--693. \MR{3570984}

\bibitem{btw18}
J.~Balogh, A.~Treglown, and A.~Z. Wagner, \emph{Tilings in {R}andomly {P}erturbed {D}ense {G}raphs}, Comb. Probab. Comput. \textbf{28} (2018), no.~2, 159--176.

\bibitem{bhkm19}
W.~Bedenknecht, J.~Han, Y.~Kohayakawa, and G.~O. Mota, \emph{Powers of tight {H}amilton cycles in randomly perturbed hypergraphs}, Random Structures Algorithms \textbf{55} (2019), no.~4, 795--807. \MR{4025389}

\bibitem{bfm04}
T.~Bohman, A.~Frieze, M.~Krivelevich, and R.~Martin, \emph{Adding random edges to dense graphs}, Random Structures Algorithms \textbf{24} (2004), no.~2, 105--117. \MR{2035870}

\bibitem{bfm03}
T.~Bohman, A.~Frieze, and R.~Martin, \emph{How many random edges make a dense graph hamiltonian?}, Random Structures Algorithms \textbf{22} (2003), no.~1, 33--42.

\bibitem{bmpp20}
J.~B\"{o}ttcher, R.~Montgomery, O.~Parczyk, and Y.~Person, \emph{Embedding spanning bounded degree graphs in randomly perturbed graphs}, Mathematika \textbf{66} (2020), no.~2, 422--447. \MR{4130332}

\bibitem{bpss23}
J.~B{\"o}ttcher, O.~Parczyk, A.~Sgueglia, and J.~Skokan, \emph{Triangles in randomly perturbed graphs}, Combinatorics, Probability and Computing \textbf{32} (2023), no.~1, 91--121.

\bibitem{fknp21}
K.~Frankston, J.~Kahn, B.~Narayanan, and J.~Park, \emph{Thresholds versus fractional expectation-thresholds}, Ann. of Math. (2) \textbf{194} (2021), no.~2, 475--495. \MR{4298747}

\bibitem{Hajnal1970}
A.~Hajnal and E.~Szemer\'{e}di, \emph{Proof of a conjecture of {P}. {E}rd{\H{o}}s}, Combinatorial theory and its applications, {II} ({P}roc. {C}olloq., {B}alatonf\"{u}red, 1969), North-Holland, Amsterdam, 1970, pp.~601--623.

\bibitem{hmt21}
J.~Han, P.~Morris, and A.~Treglown, \emph{Tilings in randomly perturbed graphs: bridging the gap between {H}ajnal--{S}zemer\'{e}di and {J}ohansson--{K}ahn--{V}u}, Random Structures Algorithms \textbf{58} (2021), no.~3, 480--516. \MR{4234994}

\bibitem{hhp19}
J.~Hladk\'{y}, P.~Hu, and D.~Piguet, \emph{Koml\'{o}s's tiling theorem via graphon covers}, J. Graph Theory \textbf{90} (2019), no.~1, 24--45. \MR{3879960}

\bibitem{jlr00}
S.~Janson, T.~{\L}uczak, and A.~Ruci\'{n}ski, \emph{Random graphs}, Wiley-Interscience Series in Discrete Mathematics and Optimization, Wiley-Interscience, New York, 2000. \MR{1782847}

\bibitem{jkv08}
A.~Johansson, J.~Kahn, and V.~Vu, \emph{Factors in random graphs}, Random Structures Algorithms \textbf{33} (2008), no.~1, 1--28.

\bibitem{jk20}
F.~Joos and J.~Kim, \emph{Spanning trees in randomly perturbed graphs}, Random Structures Algorithms \textbf{56} (2020), no.~1, 169--219. \MR{4052851}

\bibitem{kst54}
T.~K\H{o}vari, V.~S\'{o}s, and P.~Tur\'{a}n, \emph{On a problem of {K}. {Z}arankiewicz}, Colloq. Math. \textbf{3} (1954), 50--57. \MR{65617}

\bibitem{komlos00}
J.~Koml\'{o}s, \emph{Tiling {T}ur\'{a}n theorems}, Combinatorica \textbf{20} (2000), no.~2, 203--218. \MR{1767021}

\bibitem{ks95}
J.~Koml{\'o}s and M.~Simonovits, \emph{Szemer\'edi's {R}egularity lemma and its applications in graph theory}, 1995.

\bibitem{kks17}
M.~Krivelevich, M.~Kwan, and B.~Sudakov, \emph{Bounded-degree spanning trees in randomly perturbed graphs}, SIAM J. Discrete Math. \textbf{31} (2017), no.~1, 155--171. \MR{3595872}

\bibitem{kst06}
M.~Krivelevich, B.~Sudakov, and P.~Tetali, \emph{On smoothed analysis in dense graphs and formulas}, Random Structures Algorithms \textbf{29} (2006), no.~2, 180--193. \MR{2245499}

\bibitem{ko09}
D.~K\"{u}hn and D.~Osthus, \emph{The minimum degree threshold for perfect graph packings}, Combinatorica \textbf{29} (2009), no.~1, 65--107. \MR{2506388}

\bibitem{pp24}
Jinyoung Park and Huy~Tuan Pham, \emph{A proof of the {K}ahn-{K}alai conjecture}, J. Amer. Math. Soc. \textbf{37} (2024), no.~1, 235--243. \MR{4654612}

\bibitem{sz03}
A.~Shokoufandeh and Y.~Zhao, \emph{Proof of a tiling conjecture of {K}oml\'{o}s}, Random Structures Algorithms \textbf{23} (2003), no.~2, 180--205. \MR{1995690}

\bibitem{skala2013hypergeometric}
M.~Skala, \emph{Hypergeometric tail inequalities: ending the insanity}, arXiv:1311.5939, 2013.

\bibitem{szemeredi78}
E.~Szemerédi, \emph{Regular partitions of graphs}, {P}robl\'{e}mes {C}ombinatoires et {T}h\'{e}orie des {G}raphes {C}olloques {I}nternationaux {CNRS} 260 (1978), 399--401.

\end{thebibliography}

\end{document}
	

